% Unified, source-checked transcription of McAulay's three-part treatise.
% This is self-contained; Part I--III bodies are embedded below.
\documentclass[11pt]{article}
\usepackage[T1]{fontenc}
\usepackage[utf8]{inputenc}
\usepackage[margin=1.05in]{geometry}
\usepackage{amsmath,amssymb}
\usepackage{graphicx}
\usepackage{tikz}
\usetikzlibrary{decorations.pathreplacing}
\usepackage{iftex}
\usepackage{mathptmx}
\usepackage{fancyhdr}
\setlength{\parindent}{1.5em}
\setlength{\parskip}{0pt}

% Part III's inverse variance mark is the explicit U+0254 open-o/reversed-c.
\ifPDFTeX
  \newcommand{\invc}{\mathord{\rotatebox[origin=c]{180}{$c$}}}
\else
  \usepackage{fontspec}
  \newfontfamily\invcfont{DejaVu Serif}
  \newcommand{\invc}{\mathord{\text{\invcfont ɔ}}}
\fi
\newcommand{\markpair}[3]{{}^{#1#2}#3}
\newcommand{\invd}{\mathord{\vcenter{\hbox{\rotatebox[origin=c]{180}{\(\mathrm{D}\)}}}}}
\newcommand{\thetaPi}{{}_\theta\invd}
\setlength{\headheight}{14pt}
\fancypagestyle{treatise}{%
  \fancyhf{}
  \fancyhead[L]{\small\itshape Multenions and Differential Invariants}
  \fancyhead[R]{\small\thepage}
  \renewcommand{\headrulewidth}{0.4pt}
  \renewcommand{\footrulewidth}{0pt}%
}
\fancypagestyle{plain}{%
  \fancyhf{}
  \fancyfoot[C]{\small\thepage}
  \renewcommand{\headrulewidth}{0pt}
  \renewcommand{\footrulewidth}{0pt}%
}

\begin{document}
\newcommand{\unifiedtreatise}{}
\pagestyle{treatise}
\raggedbottom
% === embedded transcription bodies ===
\thispagestyle{plain}

\begin{center}\rule{0.34\textwidth}{0.5pt}\end{center}
\begin{center}\textit{Multenions and Differential Invariants.}\end{center}
\begin{center}\textsc{By Alex. McAulay, M.A., Professor of Mathematics, University of Tasmania.}\end{center}
\begin{center}(Communicated by Mr. W. B. Hardy, Sec.R.S.---Received May 13, 1920.)\end{center}

\noindent\textsc{§ 1. Preliminary Explanations.}---Since multenions appear almost to have been designed by Nature to serve as an algebra dealing with such matters as differential invariants and relativity, I have thought it desirable at the present moment to present a summary of their properties. Towards the end I propose to enter into more detail in applying them to differential invariants, and to point out some of the special properties of multenions when $n$, the vector complexity of the system, has the particular value 4.

\indent Clifford first, I believe, applied the subject to elliptic space of three dimensions. His bi-quaternion $q+\omega r$, where $\omega^2=1$, is the general member of a sub-algebra of multenions when $n=4$. In the form of multenions suitable to deal with real questions of Relativity Hamilton's bi-quaternion $q+r\sqrt{-1}$ presents itself in a precisely similar manner, and of this I shall have more to say towards the end of the paper. Joly's \emph{Manual of Quaternions} has its last chapter devoted wholly to multenions. I believe the most elaborate attempt at development is my own paper in \emph{Proc. Roy. Soc. Edin.}, 1907--8, p. 503.

\indent The notation used below is based as closely as possible on quaternion notation. It is practically identical with Joly's so far as the latter goes. It differs in several respects from the notation of my own former paper. The changes are all in the direction of simplicity, and back towards quaternion rules. The same applies to changes in terminology. These changes, which I believe are all for the better, are in the main due to friendly advice received from Prof. Knott in connection with the former paper.

\indent \emph{Law I.} Besides the unit scalar $1$ there are given $n$ generators, which will generally be regarded as representing $n$ unit mutually orthogonal vectors in a

Euclidian space of $n$ dimensions, $i_1$, $i_2$, \ldots{} $i_n$. These are
such that all the ordinary laws of algebra except the commutative law of
multiplication apply to the general integral expression $q$, called a
multenion, where

\begin{equation}
q = \sum (x_0+x_1i_1+x_2i_2+\cdots)(x'_0+x'_1i'_1+x'_2i'_2+\cdots)\cdots,
\tag{1}
\end{equation}

where every $x$ is a scalar. Since any $x$ may be negative we understand this
statement to include the ordinary algebraic uses of the signs $+$ and $-$; but
we do not understand that any assertion has been made about the use of the sign
$\div$ [$q\div r$ can later be understood to be an alternative for $qr^{-1}$].

\noindent\textsc{Law II.} Scalars are commutative not only with each other but
with each of the $n$ primitive unit vectors. [Thus scalars are commutative with
multenions.]

\begin{center}
  \textsc{Law III.} [$i_1^2$ means $i_1i_1$, etc.]
\end{center}
\begin{equation*}
  -1=i_1^2=i_2^2=\cdots=i_n^2.
\end{equation*}

\noindent\textsc{Law IV.} In a product, called a primitive vectorium,
$i_ai_b\ldots$, of primitive unit vectors, if two adjacent \emph{different}
vectors be interchanged, the product merely changes sign.

\indent Such are the assumptions. The rest is development, aided of course by
suitable terminology and symbolisation.

\noindent\textsc{Def. 1.} A primitive unit means any one of the following:
(1) the scalar $1$, (2) any primitive unit vector $i_a$, (3) any other
primitive vectorium $i_ai_b\ldots$. [Strictly (3) includes (2) and in a sense
(1).]

\noindent\textsc{Def. 2.} An expression of the form
\nopagebreak[4]
\begin{equation*}
  p=\sum_{c=1}^{n}x_ci_c
\end{equation*}
where each $x$ is a scalar, is called a vector, or a $V_1$ or a ${}_nV_1$.

\noindent\textsc{Def. 3.} Any multenion is obviously capable of expression in
the form $\sum xi_ai_b\ldots i_k$ where $i_a$, $i_b$, \ldots{} $i_k$ are $m$
different primitive vectors and $m$ ranges from $0$ to $n$, $m$ being given
$V_mq$ means the sum of all parts of the multenion for that value of $m$. Thus

\begin{equation*}
  q=V_0q+V_1q+V_2q+\cdots+V_nq.
\end{equation*}

$V_aq$ is said to be the part of $q$ of homogeneity $a$ and is called a $V_a$
or a ${}_nV_a$ or a homogeneous multenion or a hyper vector. [These alternatives
are intended to invite suggestions for a settlement of the name. I should have
preferred ``an $a$-vector'' but this clashes with relativity meanings of
4-vector and 6-vector. In multenions with $n=4$, there are two kinds of
4-vector ${}_4V_1$ and ${}_4V_3$.]

$V_0q$ will generally in words be called the scalar part of $q$ but in the
present connection we must not use $Sq$ as an alternative for $V_0q$. In
applications we very decidedly need to be able, in one system of multenions
($n=4$) to make

no change in the quaternion notation but yet to use both the quaternion
notation and the full multenion notation as applying to one symbol $q$. Thus
\begin{equation*}
  Sq=V_0q+V_4q, \qquad Vq=V_2q.
\end{equation*}
Endless confusion would result from allowing $V_0$ to be represented by $S$.

\noindent\textsc{Def. 4.} If $\alpha_1$, $\alpha_2$, \ldots{} $\alpha_a$ are
any $a$-vectors then $V_a\alpha_1\alpha_2\ldots\alpha_a$ is called a
vectorium. [Hence we have already called such an expression as $i_1i_2i_3$ a
primitive vectorium.]

\noindent\textsc{Def. 5.} A linear multenion function of a multenion is called
a multenion linity. Similarly for a vector linity or a ${}_nV_a$ linity, etc.
An extended vector linity is a special kind of multenion linity, $\Phi$, defined
from a given vector linity, $\phi$. Its fundamental property is this: If
$\alpha_1$, $\alpha_2$, \ldots{} $\alpha_a$ are any $a$-vectors whatever, $a$
itself ranging from $0$ to $n$, then
\begin{equation}
  \Phi V_a\alpha_1\alpha_2\ldots\alpha_a
    = V_a\phi\alpha_1\phi\alpha_2\ldots\phi\alpha_a.
\tag{2}
\end{equation}

[It is not obvious that this remark covers a possibility in general. The
definition, which will be modified below, is given here as a hint of the
importance of vectoriums and extended linities in our subject.]


\begin{center}
  \textit{Usual Notations.}
\end{center}

\noindent Scalars ($V_0$): $a$, $b$, $c$, $g$, $h$, $k$, $l$, $m$, $n$,
$t$ ($=$ time), $x$, $y$, $z$, $\mathfrak{J}$ ($=\sqrt{-1}$), $\theta$
($=$ arc), $\pi$.

\noindent General multenions: $p$, $q$, $r$, $s$.

\noindent Hyper vectors (homog. mult.) and vectoriums:
\[
\begin{array}{|l|c|c|c|c|c|}\hline
\text{Homogeneities} & \substack{V_a\\a} & \substack{V_b\\b} & \substack{V_c\\c} & \substack{V_n\\n} & \substack{V_2\\2}\\ \hline
\text{Hyper vectors} & u & v & w & x\upsilon & \omega \\[2pt]
\text{Vectoriums} & \breve{\alpha} & \breve{\epsilon} & \breve{\eta} & x\breve{\upsilon} & \breve{\omega} \\ \hline
\end{array}
\]

\noindent (These conventions should be noted and remembered. They immensely
simplify the reading of formulae by their property of minimising number of
suffixes, etc.)

\noindent Vectors ($V_1$): $\alpha$, $\beta$, $\gamma$, $\delta$, $\epsilon$,
$\zeta$, $\xi$, $\eta$, $\lambda$, $\mu$, $\nu$, $\rho$, $\sigma$, $\tau$.

\noindent Linear functions: $\xi$, $\eta$, $\lambda$, $\mu$, $\nu$, $\varpi$,
$v$, $\phi$, $\chi$, $\psi$.

\noindent Vectoriums: $\breve{\alpha}$, $\breve{\epsilon}$, $\breve{\eta}$,
$\breve{l}$, $\breve{\sigma}$, $\breve{\beta}$, $\breve{\upsilon}$, $\breve{\omega}$.

\noindent Differential symbols: $\delta$, $\Delta$, $\mathrm{H}$, $\mathrm{L}$,
$d$, $d_\rho$, $d_{\breve{\rho}b}$, $d_{s\alpha}$, $d\alpha$, $d\breve{b}$.

\noindent Selective linities, etc.: $\mathrm{K}$, $\mathrm{P}$, $\mathrm{Q}$,
$\mathrm{R}$, $\mathrm{S}$, $\mathrm{T}$, $\mathrm{U}$, $\mathrm{V}$.

\noindent\textsc{Def. 6.} Complexes of vectors, hypervectors, multenions,
etc., are sufficiently illustrated by multenions. If $q_1$, $q_2$, \ldots,
$q_a$ are $a$ given multenions, then all multenions of the form
$\sum_{c=1}^{a}x_cq_c$ are said to form a complex. If $q_1$, $q_2$, \ldots,
$q_a$

can be made linearly dependent on $b$ and no fewer, then $b$ is called the complexity of the complex. Thus the complexity of the complex of all vectors is $n$, of all multenions it is $2^n$, of all multenions of homogeneity $a$ it is ${}_nC_a$.

\noindent\textsc{Def. 7.} A unit vector or a unit vectorium is one, $\breve a$, such that $\breve a^2=\pm1$. [The square of every vectorium is a scalar, but $u^2$ is a scalar in general only in the four cases when $a=0,1,n-1$ or $n$.] Thus all the primitive vectoriums are unit vectoriums.

\noindent\textsc{Def. 8.} Two vectors, $\alpha$, $\beta$, are said to be normal to one another when $V_0\alpha\beta=0$. Two vector complexes are said to be completely normal to one another when every vector of the one complex is normal to every vector of the other.

\indent Every one of the symbols $V_a$ is a multenion linity. They are all commutative with one another. They are all idempotent, that is $V_a^2=V_a$, and they all satisfy the equation $V_aV_b=0$ when $a\ne b$.

\indent There is another system of multenion linities of great utility and simplicity. They are $2^n$ in number. They are all commutative with each other and with each of the $n+1$ linities $V_a$. Let $P_a$ be that definite linity of a general multenion, $q$, which simply reverses the sign of every $i_a$ which occurs in the expression (1) for $q$. From the well-known properties of linities a product of the P's such as $P_1P_2P_3$ simply reverses the sign of every $i_1$, $i_2$, and $i_3$. These linities are (1) unipotent, that is, if
\begin{equation}\phi=P_aP_b\cdots P_k,\tag{3}\end{equation}
then $\phi^2=1$; (2) they are proscriptive, by which is meant that
\begin{equation}\phi(qr)=\phi q\,\phi r,\tag{4}\end{equation}
where $q$ and $r$ are any two multenions. It is convenient to use $P$ for the product $P_1P_2\ldots P_n$ of all the linities $P_a$, so that $P$ reverses the sign of every primitive vector.

\indent $K$ is defined, consistently with its quaternion use, in the following manner. $Kq$ is a linity of a general multenion, $q$, which reverses the sequence of the primitive unit vectors in every primitive vectorium $\breve{l}$, of $q$, and \emph{also} reverses the sign of every primitive unit vector. Thus $K(i_1i_2i_3)=(-i_3)(-i_2)(-i_1)=i_1i_2i_3$. $K$ is (1) unipotent, that is $K^2=1$; (2) it is commutative with every $V$ and every $P$; (3) it is retrospective, that is $K(qr)=Kr\,Kq$. [To prove (3) first prove it when $q$ and $r$ are two primitive vectoriums and then generalise. Such a method of proof is very often applicable and may be referred to as proof by reduction to primitive units.]
\begin{equation}K(qr)=Kr\,Kq.\tag{5}\end{equation}

but if $\phi$ be given by (4), $K\phi(=\phi K)$ is also retrospective. More generally, if $\phi$, $\psi$, $\chi$, \ldots{} are any number of multenion linities, every one of which is either proscriptive or retrospective, then also is the product, $\phi\psi\chi\ldots$, proscriptive or retrospective; it is proscriptive when the number of $\phi$, $\psi$, $\chi$, \ldots{} which are retrospective is even, and it is retrospective when that number is odd.

\indent The special retrospective $PK(=KP)$ is of sufficient importance to have a symbol appropriated to it. We shall call it $Q$. Thus the product of any two of the three $P$, $Q$, $K$ is equal to the third.

\indent A caution to users of present methods in relativity and hyperbolic geometry seems desirable. It relates to the discussion of real questions by aid of the imagery
\begin{equation}\mathfrak J=\sqrt{-1}.\tag{6}\end{equation}
\indent Let us illustrate by anticipating our mode of treatment of real questions in relativity. We begin with the multenion system based on the primitive generators $i_1$, $i_2$, $i_3$, $i_4$. Then we put
\begin{equation}i_4=\mathfrak Jl,\tag{7}\end{equation}
and drop out of most of our subsequent work all reference to the original $i_4$ and $\mathfrak J$, using $l$ instead with its property $l^2=1$, as opposed to $i_4^2=-1$. Questions in relativity are then treated by aid of the multiplicative combinations of such expressions as
\begin{equation}xi_1+yi_2+zi_3+tl.\tag{8}\end{equation}
whereby, so long as all the scalars thereby introduced are real, we shall be dealing with real questions in relativity. The general multenion thus constructed contains 16 and not 32 independent real scalars. $l$ is for practical relativity purposes real. We shall call such a multenion system semi-real.

\indent Now return to the meaning above assigned to $K$.
\begin{align*}K(i_1i_2i_3i_4)&=i_1i_2i_3i_4,\\ \text{and also}\quad K(i_1i_2i_3l)&=i_1i_2i_3l,\\ \text{but whereas}\quad K'(i_1i_2i_3i_4)&=(i_1i_2i_3i_4)^{-1}.\tag{9}\end{align*}
it is not true that the corresponding equation holds when we replace $i_4$ by $l$. I have therefore refrained from making the definition of $K$ depend on such equations as (9). It is, nevertheless, a perfectly easy theorem to prove that
\begin{equation}Kl=l^{-1}.\tag{10}\end{equation}
when $l$ is any product of our \emph{original} generators $i_1$, $i_2$, \ldots{} $i_n$. This ceases to be true when $l$ takes the place of $i_4$, simply because $\mathfrak J^{-1}$ is $-\mathfrak J$ and not $+\mathfrak J$. It may be noted, however, that there is always a reproscriptive $K\phi=R$ to be found, which is such that for every semi-real vectorium $\breve{l}$, $R\breve{l}=\breve{l}^{-1}$.

Thus, suppose we replace $i_{a+1}$, \ldots{} $i_n$ by $l_{a+1}$, \ldots{} $l_n$, where
\begin{equation}i_c=\mathfrak J l_c\quad\text{and therefore}\quad l_c^2=1,\qquad c=a+1,a+2,\ldots,n.\tag{11}\end{equation}
Then
\begin{equation}R=KP_{a+1}P_{a+2}\cdots P_n=QP_1P_2\cdots P_a,\tag{12}\end{equation}
is easily shown to have the required characteristic. We now obviously have
\begin{align}P&=P_1P_2\cdots P_n,\tag{13}\\ Pu&=(-1)^a u,\qquad Qu=(-1)^{\frac12a(a-1)}u,\qquad Ku=(-1)^{\frac12a(a+1)}u.\tag{14}\end{align}
Further if we put
\begin{equation}q_c=V_cq,\qquad c=0,1,\ldots,n,\tag{15}\end{equation}
\begin{equation}\left.\begin{aligned}Pq&=q_0-q_1+q_2-q_3+\cdots,\\P_{(0)}q&=\tfrac12(1+P)q=q_0+q_2+q_4+\cdots,\\P_{(1)}q&=\tfrac12(1-P)q=q_1+q_3+q_5+\cdots\end{aligned}\right\},\tag{16}\end{equation}
\begin{equation}Qq=q_0+q_1-q_2-q_3+q_4+q_5-\cdots,\tag{17}\end{equation}
\begin{equation}Kq=q_0-q_1-q_2+q_3+q_4-q_5-\cdots.\tag{18}\end{equation}
Clearly, we might introduce $Q_{(0)}$, $Q_{(1)}$, and $K_{(0)}$, $K_{(1)}$, corresponding to $P_{(0)}$ and $P_{(1)}$. All the linities recently considered ($V_c$, $\phi$, $K\phi$) may be referred to as the capital linities.

\noindent\textsc{\S 2. Theorem of Independence and Allied Algebraic Features.}---If $q_1q_2=-q_2q_1$, then $q_1$ and $q_2$ are said to be anti-commutative. In the following enunciation, a product of $a$ different multenions from among $q_1$, $q_2$, \ldots{} $q_m$ is said to be odd-formed or even-formed, according as $a$ is odd or even.

\indent \emph{If $q_1$, $q_2$, \ldots{} $q_m$ be $m$ multenions such that $q_1^2$, $q_2^2$, \ldots{} $q_m^2$ are all scalars differing from zero, and each pair is anti-commutative; then, if $m$ is even, the $2^m$ multiplicative combinations are linearly independent; and, if $m$ is odd, they are independent unless the product $q_1q_2\ldots q_m$ of all of them is a scalar, and, if it is a scalar, $2^{m-1}$ and only $2^{m-1}$ of the combinations are independent, and these $2^{m-1}$ independent combinations may be taken as the odd-formed combinations, or instead as the even-formed combinations, or instead as the combinations from which one assigned multenion, such as $q_m$, is absent.} The theorem is true of linities as well as of multenions.

\indent Thus, if $n$ is even, a general multenion cannot be expressed by fewer than $2^n$ independent scalars. In the case of $n=4$, it cannot be expressed by fewer than sixteen independent scalars. In the case of $n=3$, it can be expressed by four independent scalars, the same number as in the case of $n=2$. [But to express it thus, we have to impose the condition that $i_1i_2i_3$ is a scalar such as $-1$. This is practically Hamilton's procedure in Quaternions.] Either of these cases might be regarded as the case of

\noindent\textsc{Quaternions.} I prefer to view Quaternions from another side, namely, as forming the subalgebra of even homogeneities for the case $n=3$.

\indent In general multenion theory, it is convenient deliberately to prescribe for the case $n$ odd that the $2^n$ primitive vectoriums shall be independent. This is equivalent to saying that the product $\breve{\upsilon}=i_1i_2\ldots i_n$ of all the primitive vectors is \emph{defined} as independent of the scalar $1$, even when $n$ is odd.

\indent Whatever be the value of $n$, multenions with even homogeneities only form a subalgebra with $2^{n-1}$ units. When $n$ is odd, this subalgebra can always be identified with the next lower multenion system.

\indent When $n$ is even, the number of independent scalars, $2^n$, is a perfect square. In this case the algebra is the equivalent of an algebra of linities with the same number. The characteristic equation, therefore, satisfied by a multenion, may be found. When $n$ is even, the degree of the characteristic equation is $2^{\frac12n}$. When $n$ is odd, the degree is $2^{\frac12(n+1)}$.

\indent In the case of $n=4$ or $3$, the characteristic function may be written in a variety of allied forms, of which the two most important are
\begin{align}chf_q(x)&\equiv(q-x)K(q-x)\cdot(2V_0-1)[(q-x)K(q-x)]\tag{1}\\&\equiv(q-x)Q(q-x)\cdot(2V_0-1)[(q-x)Q(q-x)].\tag{2}\end{align}
The coefficient of every power of $x$ is a scalar which it is not difficult to write down in full. Suppose these scalars are $m$, $m'$, $m''$, $m'''$, thus
\begin{equation}chf_q(x)=x^4-m'''x^3+m''x^2-m'x+m.\tag{3}\end{equation}
This defines the scalars $m$. We now have the identity
\begin{equation}q^4-m'''q^3+m''q^2-m'q+m\equiv0.\tag{4}\end{equation}


\indent In the case of the semi-real system suitable to relativity, $chf_q(0)$ is positive or zero, never negative. The positive value of its fourth root is the natural generalisation of the quaternion tensor $+\sqrt{qKq}$. I have been uncertain whether to denote this positive scalar by $Tq$, thus
\begin{equation}\left.\begin{aligned}Tq&=+\{qKq\cdot(2V_0-1)\cdot(qKq)\}^{\frac14}\\&=+\{qQq\cdot(2V_0-1)\cdot(qQq)\}^{\frac14}\end{aligned}\right\}.\tag{5}\end{equation}
An objection to doing so is that this is not Hamilton's meaning in the case of his biquaternion $q+\mathfrak Jr$. His tensor is then imaginary, namely, one of the values of $[(q+\mathfrak Jr)K(q+\mathfrak Jr)]^{\frac12}$. In our form below, this amounts to saying that the tensor consists of a $V_0$ part $+aV_4$ part.

\indent It suffices for our purpose to consider that $q$ only has a tensor when $qKq$ is a scalar ($V_0$). For our semi-real relativity case this must be a real scalar,

and in that case Hamilton's meaning, and the meaning given by (5), coalesce. [If $qKq$ is a scalar, then
\begin{equation}qKq=(2V_0-1)\cdot qKq,\end{equation}
and therefore $Tq$ above $=\{(qKq)^2\}^{\frac14}$.]

\indent The general value of $q$ for which $qKq$ is a scalar is
\begin{equation}q=xe^p\quad\text{where}\quad(1+K)p=0.\tag{6}\end{equation}
The following theorems will now be easily proved by any reader slightly familiar with quaternions.
\begin{align}q&=\sum i_cV_0i_c^{-1}q=\sum i_c^{-1}V_0i_cq,\tag{7}\\ \rho&=\sum_{c=1}^ni_cV_0i_c^{-1}\rho=\sum_{c=1}^ni_c^{-1}V_0i_c\rho,\tag{8}\\V_0(qV_cr)&=V_0(rV_cq)=V_0(V_cqV_cr),\tag{9}\\V_0qr&=V_0rq.\tag{10}\end{align}
whence the cyclic form
\begin{align}V_0q_1q_2\ldots q_m&=V_0q_mq_1q_2\ldots q_{m-1},\notag\\V_0(qRr)&=V_0(rRq),\tag{11}\\V_0(qKRr)&=V_0(rKRq).\tag{12}\end{align}
Here $R$ is any capital retrospective and $KR$ any capital proscriptive. Thus, for the capital linities, not only are they all commutative, but they are all self-conjugate. Conjugacy is defined and tested as in quaternions. Formally, if $\phi$ is a given multenion linity, its conjugate, $\phi'$, is the unique linity for which
\begin{equation}V_0q\phi r=V_0r\phi' q,\tag{13}\end{equation}
$q$ and $r$ being arbitrary. When $\phi=\phi'$, $\phi$ is self-conjugate. More generally, for any $\phi$, $\tfrac12(\phi+\phi')$ is the self-conjugate part and $\tfrac12(\phi-\phi')$ is the skew part. More general kinds of conjugate present themselves occasionally. For instance, let $R$ be any one of the capital retroscriptives. $Q_R$, the $R$-conjugate of $\phi$, may be defined by $V_0qR\phi r=V_0rR\phi_Rq$. With these meanings of the conjugate, $\phi'$, and the $R$-conjugate, $\phi_R$, of $\phi$, we have
\begin{equation}\left.\begin{aligned}(\phi+\psi)'&=\phi'+\psi',&(\phi\psi)'&=\psi'\phi',&(\phi')'&=\phi,\\(\phi+\psi)_R&=\phi_R+\psi_R,&(\phi\psi)_R&=\psi_R\phi_R,&(\phi_R)_R&=\phi,\\\phi_R&=R\phi'R,&\phi'&=R\phi_RR\end{aligned}\right\}.\tag{14}\end{equation}
The general meaning of a conjugate, $\phi_c$, of $\phi$ may be taken to be that it is a unipotent retrospective linity of the linity $\phi$. Thus $\phi'$ is a linity function of the linity $\phi$: (1) unipotent because $(\phi')'=\phi$; (2) retrospective because $(\phi\psi)'=\psi'\phi'$; (3) linear because $(\phi+\psi)'=\phi'+\psi'$. I have before

me in MS. a full analysis of the properties of such a conjugate, but must not give it here. [$R q$ stands to $q$ as $\phi'$ stands to $\phi$.] The capital linities are self-$R$-conjugate as well as self-conjugate. If $\phi q=pqr$, then
\begin{equation}\phi'q=rqp,\qquad \phi_Rq=Rp\cdot q\cdot Rr,\tag{15}\end{equation}
and similarly if $\phi q=p_1qr_1+p_2qr_2+\cdots$.

\indent Frequently, as in proving the above, we require
\begin{equation}V_0pr=V_0qr,\quad\text{then}\quad p=q,\tag{16}\end{equation}
$r$ being arbitrary, and again
\begin{equation}V_0\alpha\rho=V_0\beta\rho,\quad\text{then}\quad\alpha=\beta,\tag{17}\end{equation}
$\rho$ being arbitrary, and similarly for a ${}_nV_a$ instead of a general multenion or a vector.

\indent When in a semi-real system $R$ is chosen so that $R\breve{i}=\breve{i}^{-1}$ for every $i$, we have that if $q=\sum xi$ and $q'=\sum x'i$, then
\[V_0qRq'=\sum xx',\]
and in particular
\[V_0qRq=\sum x^2.\]

\noindent\textsc{\S 3. The Commutation Theorem.}---If $\breve{i}_a$ is of homogeneity $a$, and $\breve{i}_b$ of homogeneity $b$, and if there are exactly $c$ primitive vectors common to the $a$ vectors of $\breve{i}_a$ and the $b$ vectors of $\breve{i}_b$ it is easy to see that $\breve{i}_a\breve{i}_b$ is of homogeneity $a+b-2c$, and that $\breve{i}_a\breve{i}_b=(-1)^{ab-c}\breve{i}_b\breve{i}_a$. Hence
\begin{equation}V_{a+b-2x}uv=(-1)^{ab-x}V_{a+b-2x}vu.\tag{1}\end{equation}
[Obviously true when $u=\breve{i}_a$, $v=\breve{i}_b$ as we see by considering separately the two cases $x=c$, $x\ne c$. It is therefore true for $u$ and $v$. This is a case of proof by reduction to primitive units.]

\indent (1) may be called a first form of the commutation theorem as it shows us a first effect of commuting $u$ and $v$ in $uv$. (1) at once leads to what may be called the full form
\begin{equation}\left.\begin{aligned}uv&=(V_{a+b}+V_{a+b-2}+\cdots+V_{a-b})uv\\&=(-1)^{ab}(V_{a+b}-V_{a+b-2}+\cdots+[-1]^bV_{a-b})vu\\&=(-1)^{a+b}(V_{a-b}-V_{a-b+2}+\cdots+[-1]^bV_{a+b})vu\end{aligned}\right\}.\tag{2}\end{equation}
The following are easy deductions
\begin{equation}\left.\begin{aligned}\tfrac12[uv+(-1)^{ab}vu]&=(V_{a+b}+V_{a+b-4}+V_{a+b-8}+\cdots)uv,\\\tfrac12[uv-(-1)^{ab}vu]&=(V_{a+b-2}+V_{a+b-6}+\cdots)uv\end{aligned}\right\},\tag{3}\end{equation}
and these may be written in the reverse order with $(-1)^{ab+b}$ in place of $(-1)^{ab}$. From (3)
\begin{equation}\left.\begin{aligned}\tfrac12(uv+vu)&=(V_{a+2}+V_{a-2})uv=(V_{a+2}+V_{a-2})vu,\\\tfrac12(uv-vu)&=V_a\cdot uv=-V_a\cdot vu\end{aligned}\right\},\tag{4}\end{equation}

which has important applications to infinitesimal rotations. Other special cases are
\begin{align}u^2&=(V_0+V_4+V_8+\cdots)u^2,\tag{5}\\\omega^2&=(V_0+V_4)\omega^2,\tag{6}\\\alpha\beta&=(V_0+V_2)\alpha\beta=(V_0-V_2)\beta\alpha,\tag{7}\\V_2\alpha\beta&=-V_2\beta\alpha,\qquad V_0\alpha\beta=V_0\beta\alpha,\tag{8}\\V_2\alpha\beta&=\tfrac12(\alpha\beta-\beta\alpha),\qquad V_0\alpha\beta=\tfrac12(\alpha\beta+\beta\alpha).\tag{9}\end{align}
\noindent\textsc{\S 4. Rotations of Three Kinds.}---Suppose we find $n$ multenions $p_1,p_2\ldots p_n$ such that for all purposes of Laws I to IV they may take the place of $i_1,i_2\ldots i_n$, are we to call the substitution of the multenions $p$ for the vectors $i$ a rotation? In Euclidean space of three dimensions we should not always so affirm. If $i,j,k$ should represent, as sometimes in quaternions they do, three directions in a rigid body we could not move the body so that $i$ became $-i$ and $j$ and $k$ remained unaltered; but they would continue to satisfy the conditions of the laws. We should call the transformation a perversion or reflection rather than a rotation. As the algebra of multenions is ($n$ even) a full linity algebra, or ($n$ odd) a sub-algebra of such a linity algebra, we will assume that given $q$, $q^{-1}$ can be obtained uniquely in general so that $qq^{-1}=1=q^{-1}q$ (though in certain singular cases we have to say instead that $q_0$ can be found so that $qq_0=0=q_0q$). In the general case we will say that $q$ is vertible (\emph{i.e.}, it has an inverse $q^{-1}$); in the singular case that $q$ is non-vertible.

\indent \emph{Theorem concerning multenion rotation.} (1) \emph{When $q$ is a given vertible multenion, then in the four fundamental laws the primitive vectors $i_1,i_2,\ldots i_n$ may be replaced respectively by}
\begin{equation}p_1=qi_1q^{-1},\ldots,\qquad p_n=qi_nq^{-1},\tag{1}\end{equation}
\emph{and the effect of the replacement is to change any multenion $r$ to}
\begin{equation}r'=qrq^{-1}.\tag{2}\end{equation}
\emph{(2) Conversely if $p_1,p_2\ldots p_n$ are $n$ multenions such that}
\begin{equation}-1=p_1^2=p_2^2=\cdots=p_n^2,\tag{3}\end{equation}
\emph{and every pair is anti-commutative, that is}
\begin{equation}p_1p_2+p_2p_1=0,\ \text{etc.},\tag{4}\end{equation}
\emph{and when $n$ is odd}
\begin{equation}p_1p_2\ldots p_n=i_1i_2\ldots i_n,\tag{5}\end{equation}
\emph{then there exists a vertible multenion, $q$, for which (1) and (2) are true; $r'$ now meaning the multenion obtained from the given multenion $r$ by replacing $i_1,i_2,\ldots i_n$ by $p_1,p_2,\ldots p_n$ respectively.}

This kind of rotation is scarcely what we seek. For instance, it does not in general leave $V_c$ unchanged, that is, it is not in general true that
\begin{equation}qV_cr\cdot q^{-1}=V_c(qrq^{-1}).\tag{6}\end{equation}
Multenion rotation $qrq^{-1}$ is operation by a multenion linity $\phi=q()q^{-1}$ such that $\phi'\phi=1=\phi\phi'$, and conversely.

\indent Next, imposing the condition that (6) is to be satisfied, we find that $q$ must be the product of any number of vectors, and that if it is such a product (6) is satisfied. We may next find what is the effect of infinitesimally changing the vectors and attempt to build up the finite continuous group of which that infinitesimal transformation is the element. The infinitesimal rotation is the transformation effected by changing $q$ to $q+\tfrac12(\omega q-q\omega)$, where $\omega$ is infinitesimal and the corresponding finite rotation is $e^\omega()e^{-\omega}$, where $\omega$ is finite. $\omega$ is a perfectly arbitrary ${}_nV_2$, infinitesimal of course for the first case. When $\omega$ is not singular ($\omega^2=0$) its general form may be put
\[\omega=a\epsilon_1\epsilon_2+b\epsilon_3\epsilon_4+\cdots,\]
where $\epsilon_1,\epsilon_2,\epsilon_3,\epsilon_4\ldots$ are orthogonal unit vectors,
\[e^\omega\equiv(\cos a+\epsilon_1\epsilon_2\sin a)(\cos b+\epsilon_3\epsilon_4\sin b)\cdots\]
and the factors here are commutative. Putting the last in the form
\[e^\omega=\epsilon_1(\epsilon_1^{-1}\cos a+\epsilon_2\sin a)\epsilon_3(\epsilon_3^{-1}\cos b+\epsilon_4\sin b)\cdots,\]
we see that $e^\omega$ is the product of an \emph{even} number of vectors. The general form of $\omega$ when it is singular is not simple, but I have a complete solution before me in MS.

\indent To show that the infinitesimal rotation is as stated will not occupy too much space. Suppose
\[r'=\alpha\beta\ldots\lambda\cdot r\cdot\lambda^{-1}\ldots\beta^{-1}\alpha^{-1}=prp^{-1},\]
where the vectors $\alpha,\beta,\gamma,\ldots$ are assumed to be unit vectors, as that does not affect the transformation. Let $\alpha,\beta,\ldots$ change to $\alpha+d\alpha,\beta+d\beta,\ldots$, still remaining unit vectors.
\[dr'=d(prp^{-1})=dpp^{-1}\cdot r'-r'dpp^{-1}.\]
That $dpp^{-1}$ is a ${}_nV_2$ is straightforwardly proved thus
\[dpp^{-1}=d\alpha\alpha^{-1}+\alpha\cdot d\beta\beta^{-1}\cdot\alpha^{-1}+\alpha\beta\cdot d\gamma\gamma^{-1}\cdot\beta^{-1}\alpha^{-1}+\cdots.\]
Every term here is a ${}_nV_2$; for example, consider the third term. Since $\gamma$ and $\gamma+d\gamma$ are unit vectors, $V_0d\gamma\gamma^{-1}=0$, and therefore $d\gamma\gamma^{-1}$ is a ${}_nV_2$. It follows that the third term is a ${}_nV_2$. If we put $u,u'$ for $r,r'$ and $\tfrac12d\omega$ for $dpp^{-1}$,
\[du'=\tfrac12(d\omega u'-u'd\omega)=V_0d\omega u'.\]
This is the very simple, very useful form in which true infinitesimal rotation

presents itself in our present subject. A rotation of vectors $\rho$ is a linity $\phi\rho$ of $\rho$ such that $\phi'\phi=1$.

\indent When without qualification we speak of rotation we shall understand it to mean $e^\omega()e^{-\omega}$ (a continuous group).

\indent The skew part of any vector linity, $\phi$, is of the form of an infinitesimal rotation. Thus, introduce the $\zeta$ notation of quaternions,\footnote{See \emph{Utility of Quaternions}.} by which
\begin{equation}\mathrm F(\zeta,\zeta)=\sum_{c=1}^n\mathrm F(i_c,i_c),\tag{7}\end{equation}
where $\mathrm F$ is bilinear in its two members. We have
\begin{align}\rho&=-\zeta V_0\zeta\rho,\tag{8}\\\phi'\rho&=-\zeta V_0\zeta\phi'\rho=-\zeta V_0\rho\phi\zeta.\tag{9}\end{align}
The skew part is
\begin{equation}\tfrac12(\phi-\phi')\rho=\tfrac12V_1(V_2\zeta\phi\zeta)\rho=V_1\omega\rho.\tag{10}\end{equation}
Here we have assumed that
\[V_1(V_2\alpha\beta)\gamma=\alpha V_0\beta\gamma-\beta V_0\gamma\alpha,\]
much as in quaternions. This is not difficult to prove, and we shall in §7 enunciate much more general theorems of which this is a particular case. The pure part does not similarly reduce. We have
\begin{equation}\tfrac12(\phi+\phi')\rho=-\tfrac12(\zeta V_0\rho\phi\zeta+\phi\zeta V_0\rho\zeta).\tag{11}\end{equation}
\noindent\textsc{§ 5. Vectoriums.} Let
\[\alpha^{(a)}=\alpha_1\alpha_2\ldots\alpha_a,\qquad\beta^{(b)}=\beta_1\beta_2\ldots\beta_b,\qquad\gamma^{(c)}=\gamma_1\gamma_2\ldots\gamma_c.\]
Then in the expression $V_{a+b+c}(\alpha^{(a)}\beta^{(b)}\gamma^{(c)})$ it is obvious by reduction to primitive units that within the brackets we may at will write either $\alpha^{(a)}$ or $V_a\alpha^{(a)}$, and $\beta^{(b)}$ or $V_b\beta^{(b)}$, and $\gamma^{(c)}$ or $V_c\gamma^{(c)}$; the value of the expression is unaltered by such insertion or removal of $V_a$, $V_b$, or $V_c$. The statement, of course, is in general only true of this particular part, $V_{a+b+c}$ of the product of highest homogeneity. Next, it is similarly evident that if we change, say, $\alpha_2$ to $x\alpha_2+y\alpha'_2$ we simply alter the corresponding vectorium similarly, that is
\[V_a\alpha_1(x\alpha_2+y\alpha'_2)\alpha_3\ldots\alpha_a=xV_a\alpha_1\alpha_2\alpha_3\ldots\alpha_a+yV_a\alpha_1\alpha'_2\alpha_3\ldots\alpha_a.\]
Again, if in the vectorium we interchange two consecutive members such as $\alpha_2,\alpha_3$ we merely change the sign of the whole. For we may successively write instead of $\alpha_2\alpha_3$
\[\alpha_2\alpha_3,\qquad V_2\alpha_2\alpha_3,\qquad-V_2\alpha_3\alpha_2,\qquad-\alpha_3\alpha_2.\]
Finally, by a familiar simple process we may interchange any two members by successive interchanges of neighbours, and as a final result, again, the vectorium changes in sign merely. Thus $V_a\alpha^{(a)}$ is ``combinatorial.'' The

common simple properties of a combinatorial multilinear function are probably well known to my readers and may be assumed. [Grassmann apparently first, in a perfectly general manner, appreciated the value of, and fully established these properties. Indeed, $V_a\alpha^{(a)}$ is both algebraically and geometrically identical in meaning with Grassmann's combinatorial ``product.'' I should like for a second time to express my unbounded admiration for Grassmann's great pioneer work, though I cannot help believing that the actual form of the matters he dealt with has been put out of court by later developments.]

\indent We may note the following theorem in passing. It is our form of one of Grassmann's products. If $\breve\rho$ is a vectorium of homogeneity $g$, then
\[V_{a+b+c}(V_{g+a+b}\breve\rho uv)(V_{g+c}w\breve\rho)=V_{a+b+c}(V_{g+a}\breve\rho u)(V_{g+b+c}v\breve\rho w).\]
Of $u,v,w$ the middle one, $v$, may be passed over from the first to the second bracket-pair.


\indent The following is wanted immediately and is constantly in request. It may be thought of as a passage into or out of the operation of a $V_a()$ [whether as pre-factor or post-factor] of either $\breve\upsilon$ or $p_n$. We have
\[\breve\upsilon V_aq=V_{n-a}\breve\upsilon q,\qquad V_aq\cdot\breve\upsilon=V_{n-a}(q\breve\upsilon).\]
Also, if $p_n$ is any multenion of homogeneity $n$, $p_n=x\breve\upsilon$, and, therefore, in the present process, $p_n$ behaves exactly as does $\breve\upsilon$, that is
\[p_nV_aq=V_{n-a}p_nq,\qquad V_aq\cdot p_n=V_{n-a}qp_n.\]

\indent Suppose there are $n$ given linearly independent vectors $\alpha_1,\alpha_2,\ldots\alpha_n$. Let $\alpha^{(a)}$ be the product of any $a$ of them, and $\alpha^{(n-a)}$ the product of the rest in such sequence that $V_n\alpha^{(a)}\alpha^{(n-a)}$ always retains the same value. Let
\begin{equation}\left.\begin{aligned}\breve\alpha&=V_a\alpha^{(a)},\\\breve\alpha&=V_{n-a}\alpha^{(n-a)}\cdot V_n^{-1}\alpha^{(a)}\alpha^{(n-a)}\end{aligned}\right\}.\tag{1}\end{equation}
[By a convenient, strictly speaking indefensible, custom, since $V_n^{-1}$ really $=\infty$, we mean by $V_n^{-1}\alpha^{(a)}\alpha^{(n-a)}$ the reciprocal $(V_n\alpha^{(a)}\alpha^{(n-a)})^{-1}$.]

\indent There are $2^n$ linearly independent [proof of linear independence here omitted] values of $\breve\alpha$, and they will be called the $2^n$ vectoriums formed from $\alpha_1,\alpha_2,\ldots\alpha_n$. [The term may, of course, be used when $\alpha_1,\alpha_2,\ldots\alpha_n$ are not linearly independent.] The $2^n$ corresponding values of $\breve\alpha$, which are also linearly independent, will be called the normal reciprocals of the vectoriums, $\breve\alpha$. They are so called because
\[V_0\breve\alpha\breve\alpha=1,\qquad V_0\breve\alpha_1\breve\alpha=0,\]
$\breve\alpha_1$ being a second (not corresponding to $\breve\alpha$) $\breve\alpha$. [Thus, when $n=4$, there are

sixteen equations of the type $V_0\breve\alpha\hat\alpha=1$, and 240 equations of the type $V_0\breve\alpha_1\hat\alpha=0$.] Both equations are proved quite simply from (1) by passing $p_n=V_n^{-1}\alpha^{(a)}\alpha^{(n-a)}$ out of the operation of $V_0$ in $V_0\breve\alpha\hat\alpha$.

\indent Let now $q=\sum x\breve\alpha=\sum y\hat\alpha$. From the normal reciprocal relations we at once get $x=V_0q\hat\alpha$, $y=V_0q\breve\alpha$. Hence
\begin{equation}\left.\begin{aligned}q&=\sum\hat\alpha V_0q\hat\alpha=\sum V_a\alpha^{(a)}\cdot V_nq\alpha^{(n-a)}\cdot V_n^{-1}\alpha^{(a)}\alpha^{(n-a)},\\&=\sum\hat\alpha V_0q\breve\alpha=\sum V_0(qV_a\alpha^{(a)})\cdot V_{n-a}\alpha^{(n-a)}\cdot V_n^{-1}\alpha^{(a)}\alpha^{(n-a)}\end{aligned}\right\},\tag{2}\end{equation}
the last result in the first line coming from a new application of the $p_n$ theorem. If in (2) we put $q=u$, getting a more particular result, we should at once see that (2) is analogous to the quaternion theorems
\[\rho S\alpha\beta\gamma\equiv\alpha S\beta\gamma\rho+\cdots\equiv V\beta\gamma S\alpha\rho+\cdots.\]
A particular case ($a=1$) of (1) is
\begin{equation}\left.\begin{aligned}\hat\alpha_1&=V_{n-1}\alpha_2\alpha_3\ldots\alpha_n\cdot V_n^{-1}\alpha_1\alpha_2\ldots\alpha_3\\&=-V_{n-1}\alpha_1\alpha_3\ldots\alpha_n\cdot V_n^{-1}\alpha_1\alpha_2\ldots\alpha_3=\text{etc.}\end{aligned}\right\}.\tag{3}\end{equation}
This is clearly the solution of the problem. Given
\[S\alpha_c\hat\alpha_c=1,\quad c=1,2,\ldots,n,\qquad S\alpha_b\hat\alpha_c=0,\quad b\ne c;\ b,c=1,2,\ldots,n,\]
what are the vectors $\hat\alpha_c$ in terms of the vectors $\alpha_c$? And, of course, we have a precisely similar solution for $\alpha_c$ in terms of $\hat\alpha_c$. It appears, then, that the relations between the quantities $\hat\alpha$ and the quantities $\hat\alpha$ are symmetrical. No doubt we have here established this if only it be evident what is the precise significance of the word ``symmetrical'' in this connection. The need of the caution may be illustrated thus. If we put $\hat\alpha=V_a\hat\alpha_1\hat\alpha_2\ldots\hat\alpha_n$, we are tempted to say that $\hat\alpha$ must be $V_a\hat\alpha_1\hat\alpha_2\ldots\hat\alpha_a$, but we shall find that we do not thereby always obtain $S\hat\alpha\hat\alpha=1$. The true value of $\hat\alpha$ is
\begin{equation}\left.\begin{aligned}\hat\alpha&=V_a\hat\alpha_a\hat\alpha_{a-1}\ldots\hat\alpha_1=QV_a\hat\alpha_1\hat\alpha_2\ldots\hat\alpha_a,\\\text{when}\qquad\hat\alpha&=V_a\alpha_1\alpha_2\ldots\alpha_a\end{aligned}\right\},\tag{4}\end{equation}
If the $\alpha_c$ are taken as the $i_c$, it is obvious that the $\hat\alpha$ are the primitive vectoriums, and, from the normal reciprocal relations, it is then obvious that the $\hat\alpha$ are also the primitive vectoriums (with different signs in accordance with) $\hat\alpha=K\hat\alpha$.

[Thus $0'=1$, $1'=-1$, $2'=-1/2!$, $3'=+1/3!$ etc., the signs running $+--++--++---\ldots$.]
\begin{align}\sum\mathrm F(i_a,i_a^{-1})&=a'\mathrm F(V_a\zeta^{(a)},V_a\zeta^{(a)}),\tag{6}\\u&=\sum i_aV_0i_a^{-1}=a'V_a\zeta^{(a)}\cdot V_0i_a\zeta^{(a)}.\tag{7}\end{align}
In (6) put the second $V_a\zeta^{(a)}=\sum\hat\alpha V_0\hat\alpha\zeta^{(a)}$. Thus,
\[a'\mathrm F(V_a\zeta^{(a)},V_a\zeta^{(a)})=a'\sum(V_a\zeta^{(a)}V_0\hat\alpha\zeta^{(a)},\hat\alpha)=\sum\mathrm F(\hat\alpha,\hat\alpha).\]
We might have modified the first $V_a\zeta^{(a)}$ instead of the second. Thus collecting results
\begin{equation}\left.\begin{aligned}a'\mathrm F(V_a\zeta^{(a)},V_a\zeta^{(a)})&=\sum\mathrm F(i_a,i_a^{-1})\\&=\sum\mathrm F(\hat\alpha,\hat\alpha)=\sum\mathrm F(\alpha,\hat\alpha)\end{aligned}\right\}.\tag{8}\end{equation}
(${}_nC_a$ terms in each summation.) Just as in quaternions $\mathrm F(\zeta,\phi\zeta)=\mathrm F(\phi'\zeta,\zeta)$ so now if $\phi$ is a ${}_nV_a$ linity
\begin{equation}a'\mathrm F(V_a\zeta^{(a)},\phi V_a\zeta^{(a)})=a'\mathrm F(\phi'V_a\zeta^{(a)},V_a\zeta^{(a)}).\tag{9}\end{equation}
If $\phi$ is a multenion linity
\begin{equation}\sum_{a=0}^na'\mathrm F(V_a\zeta^{(a)},\phi V_a\zeta^{(a)})=\sum_{a=0}^na'\mathrm F(\phi'V_a\zeta^{(a)},V_a\zeta^{(a)}).\tag{10}\end{equation}
For $a=0$ we must put $a'=1$, and for each $V_a\zeta^{(a)}$ we must also put $1$ [see (7) above].

\noindent\textsc{§ 6. Extended Vector Linities.}---The theorems of §§ 6, 7 find continuous applications in the multenion treatment of differential invariants. Such treatment may be considered to begin in §8, and is explicitly introduced in §9.

\indent Let $\phi$ be a given vector linity. The extended linity (a multenion linity) which may for most purposes, but not for all, also be denoted by $\phi$, is here defined thus. Let $i_a=i_gi_h\ldots i_l$ be any primitive unit of homogeneity $a$, $a$ itself being arbitrary. Then
\[\phi V_ai_a=V_a\phi i_g\phi i_h\ldots\phi i_l.\]
Now change the $n$ primitive vectors $i_1,i_2,\ldots i_n$ to any other $n$ vectors by such steps as replacing $i_1$ by $xi_1+yi_2$ (or, if $i_1,i_2$ have already changed to $\epsilon_1,\epsilon_2$, the change should be from $\epsilon_1$ to $x\epsilon_1+y\epsilon_2$). Such a change simply multiplies each side of the defining equation (when $i_1$ occurs in it) by $x$; the meaning is therefore unaltered. It follows that $\phi V_a\alpha^{(a)}=V_a(\phi\alpha)^{(a)}$ where $a$ is arbitrary and the vectors $\alpha$ are arbitrary. If, then, $(\phi),(\psi)$ are the extended linities of $\phi,\psi$, we see that $(\phi)(\psi)$ is the extended linity of $\phi\psi$, for
\[(\phi)(\psi)V_ai_gi_h\ldots i_l=(\phi)V_a\psi i_g\psi i_h\ldots\psi i_l=V_a\phi\psi i_g\phi\psi i_h\ldots\phi\psi i_l.\]
From this again it follows that the extended linity of $\phi^{-1}$ is $(\phi)^{-1}$.

Let us use either $(\phi)$ or $\phi$ for the extended linity of $\phi$ for the present. What we have just proved may be put
\[(\phi)(\psi)=(\phi\psi),\qquad (\phi)^{-1}=(\phi^{-1}).\]
We are now about to prove a similar proposition concerning $\phi'$ the conjugate, namely
\[(\phi)'=(\phi').\]
By (7) of §5 we have
\begin{equation}\phi u=a'\phi V_a\zeta^{(a)}V_0u\zeta^{(a)},\qquad (\phi)u=a'V_a(\phi\zeta)^{(a)}V_0u\zeta^{(a)}.\tag{1}\end{equation}
Let $u,u_0$ be two arbitrary multenions of homogeneity $a$. Operate on (1) by $V_0u_0()$
\[V_0u_0(\phi)u=a'V_0u_0(\phi\zeta)^{(a)}V_0u\zeta^{(a)}.\]
Modifying the left by the fundamental rule of the meaning of $(\phi)'$, and modifying the right by the rule $\mathrm F(\zeta,\phi\zeta)=\mathrm F(\phi'\zeta,\zeta)$ applicable to a vector linity. We obtain
\[V_0u(\phi)'u_0=a'V_0u_0\zeta^{(a)}V_0u(\phi'\zeta)^{(a)}.\]
Since $u$ is arbitrary, by (16) of §3
\[(\phi)'u_0=a'V_a(\phi'\zeta)^{(a)}V_0u_0\zeta^{(a)},\]
which by comparison with (1) proves the assertion that $(\phi)'=(\phi')$.

\indent The following theorems are also true
\begin{align}\phi V_{a+b}uv&=V_{a+b}\phi u\phi v,\tag{2}\\\phi V_{a-b}uv&=V_{a-b}\phi u\phi'^{-1}v,\tag{3}\end{align}
when $\phi$ is non-vertible we must change $v$ to $\phi'v$. An alternative, less convenient than (3) generally, is thus
\begin{equation}\phi V_{a-b}u\phi'v=V_{a-b}\phi uv,\tag{4}\end{equation}
but this alternative is true without qualification. $\breve\upsilon$ having its usual meaning $(i_1,i_2,\ldots i_n)$ it is easy to see that $\phi\breve\upsilon=m\breve\upsilon$ where $m$ is the determinant of $\phi$. Putting $u=\breve\upsilon$ in (3) and (4) we obtain
\begin{equation}\left.\begin{aligned}\breve\upsilon^{-1}\phi(\breve\upsilon v)&=m\phi'^{-1}v,\\\text{or}\qquad\breve\upsilon^{-1}\phi(\breve\upsilon\phi'v)&=mv\end{aligned}\right\}.\tag{5}\end{equation}
[\emph{Added February, 1921.}---Closely connected with the extension in §6 from vector operation to multenion operation, of a given vector linity, is a second kind of such extension. If $\phi_E$ is the extended (or, as I now prefer to say

\emph{extensive}) form of $\phi$ of §6, its fundamental property, showing its dependence (in a \emph{multiple} manner) on $\phi$, is that
\begin{equation}\phi_E V_a\alpha_1\alpha_2\ldots\alpha_a=V_a\phi\alpha_1\phi\alpha_2\ldots\phi\alpha_a.\tag{1}\end{equation}
Similarly, if $\psi_e$ is the new extensive form of $\phi$, its fundamental property showing its dependence (in an \emph{additive} manner) on $\phi$ is that
\begin{equation}\psi_eV_a\alpha_1\alpha_2\ldots\alpha_a=V_a(\psi\alpha_1)\alpha_2\ldots\alpha_a+V_a\alpha_1(\psi\alpha_2)\alpha_3\ldots\alpha_a+\cdots.\tag{2}\end{equation}
The dependence is called multiple or additive because
\begin{equation}(\phi\Phi)_E=\phi_E\Phi_E,\qquad(\psi+\Psi)_e=\psi_e+\Psi_e.\tag{3}\end{equation}
We may speak briefly of $\phi_E$ as the extensive of $\phi$ and $\psi_e$ as the subextensive of $\psi$ (using extensive and subextensive both as nouns and adjectives; compare offensive and defensive; compare also Whitehead's translation of ``Ausdehnungslehre'' as ``Calculus of extension'').

\indent For both kinds of extensive it is true, when $\phi_E$ and $\psi_e$ operate on a mere vector, that
\begin{equation}\phi_E\rho=\phi\rho,\qquad\psi_e\rho=\psi\rho,\tag{4}\end{equation}
so that quite frequently it is convenient to use $\phi$ for $\phi_E$ and on other occasions to use $\psi$ for $\psi_e$. [We should very rarely with a given $\phi$ want on a single occasion to use both $\phi_E$ and $\phi_e$, but if we do we must have some mark to distinguish them, since they have diverse meanings.]

\indent Although the properties of a subextensive linity may be derived directly from its fundamental property defined above, they may instead be derived from those of an extensive by the methods connecting an infinitesimal with a finite transformation. Thus we may put
\begin{equation}\left.\begin{aligned}\phi&=1+\dot\phi\,dt=1+\psi\,dt,\\(1+\psi\,dt)_E&=\phi_E=(1+\dot\phi\,dt)_E=1+\psi_e\,dt\end{aligned}\right\},\tag{5}\end{equation}
where $dt$ is an arbitrary infinitesimal scalar. Thus from former results we at once have
\begin{align}V_c\psi_e&=\psi_eV_c=V_c\psi_eV_c,\tag{6}\\(\psi_e)'&=(\psi')_e,\tag{7}\\\psi_eV_{a+b}uv&=V_{a+b}(\psi_eu\cdot v+u\psi_ev),\tag{8}\\\psi_eV_{a-b}uv&=V_{a-b}(\psi_eu\cdot v-u\psi'_ev),\tag{9}\end{align}
\begin{align}
\psi_e\breve\upsilon&=\psi'_e\breve\upsilon=-\breve\upsilon\cdot V_0\zeta\psi\zeta,\tag{10}\\
\psi_ev&=-V_b\psi\zeta V_{b-1}\zeta
  =-V_b(V_0\psi\zeta\zeta\cdot v-\zeta V_{b+1}\psi\zeta v),\tag{11}\\
\breve\upsilon^{-1}\psi_e(\breve\upsilon v)&=-(\psi'_e+V_0\zeta\psi\zeta)v
  =-V_b\psi\zeta V_{b+1}\zeta v.\tag{12}
\end{align}
the last being obtained by putting $u=\breve\upsilon$ in (9) above.

If $\phi$ and $\psi$ are given vector linities then $\phi_E\psi_e\phi_E^{-1}$ is a subextensive; that is it is the subextensive of $\phi\psi\phi^{-1}$. For by (11)
\[(\phi\psi\phi^{-1})_ev=-V_b\phi\psi\phi^{-1}\zeta V_{b-1}\zeta
 =-\phi_EV_b\psi\zeta V_{b-1}\zeta\phi_E^{-1}v,\]
by applications of both the equations [see (2) and (3) of §6],
\[\phi_EV_{a+c}uw=V_{a+c}\phi_Eu\phi_Ew,\qquad
\phi'_EV_{c-a}uw=V_{c-a}\phi_E'^{-1}u\phi_Ew.\]
If $\phi$ is of assigned dimensions, $\phi_e$ is of the same, but not $\phi_E$. $\phi_Eu$, in so far as its dimensions depend on $\phi_E$ and not on $u$, are $a$ times the dimensions of $\phi$.]

\noindent\textsc{§ 7. Formulae Analogous to the Quaternion Formulae for $V_\alpha V_\beta\gamma$ and $V_\alpha\beta\gamma$.}---Let $\alpha_1,\alpha_2,\ldots,\alpha_{a+b}$ be $a+b$ given vectors and let $\alpha^{(a)}$ be the product of any $a$ of them and $\alpha^{(b)}$ the product of the rest in such a sequence that $V_{a+b}\alpha^{(a)}\alpha^{(b)}$ always has the same value. Then $u$ being as usual
\begin{equation}V_b\cdot uV_{a+b}\alpha^{(a)}\alpha^{(b)}=V_b(V_{a+b}\alpha^{(b)}\alpha^{(a)}\cdot u)=\sum V_b\alpha^{(b)}V_0\alpha^{(a)}u.\tag{1}\end{equation}
Note that $\alpha^{(a)}$ occupies the middle position throughout. If in this we put $b=n-a$ it easily leads to the last form of (2) of §5.

$u$ and $v$ being as usual and $\rho$ any vector,
\begin{align}V_{1+(a\cdot b)}\cdot uV_{b+1}\rho v&=V_{1+(a\cdot b)}[V_{a-1}u\rho\cdot v+(-1)^a\rho V_{a\cdot b}uv],\tag{2}\\V_{1+(a\cdot b)}(V_{a+1}u\rho\cdot v)&=V_{1+(a\cdot b)}[uV_{b-1}\rho v+(-1)^a\rho V_{a\cdot b}uv],\tag{3}\end{align}
\begin{equation}\left.\begin{aligned}V_{1+(a\cdot b)}u\rho v&=V_{1+(a\cdot b)}[V_{a-1}u\rho\cdot v+uV_{b-1}\rho v+(-1)^a\rho V_{a\cdot b}uv]\\&=V_{1+(a\cdot b)}[V_{a+1}u\rho\cdot v+uV_{b+1}\rho v-(-1)^a\rho V_{a\cdot b}uv]\end{aligned}\right\},\tag{4}\end{equation}
\begin{equation}\left.\begin{aligned}V_{a+b-1}u\rho v&=V_{a+b-1}[V_{a-1}u\rho\cdot v+uV_{b+1}\rho v+(-1)^a\rho V_{a+b}uv]\\&=V_{a+b-1}[V_{a+1}u\rho\cdot v+uV_{b-1}\rho v-(-1)^a\rho V_{a+b}uv]\end{aligned}\right\}.\tag{5}\end{equation}
To make practical applications of multenions it is absolutely essential that facility should be gained in using these formulae. The following comments are to aid in the acquirement of the requisite facility. I do not attempt to work without a list of at least the four equations (4) and (5) before me. Each of those four generates two others. (2) and (3) are those two derived from the first equation (4). If the second term on the right of this equation (call it (4a)) is removed to the left it can be made to merge with the left and equation (2) is produced. Similarly, if the first term on the right of (4a) is removed to the left equation (3) is produced. Now as after a little practice, with (4a) before the eye, the whole of this process can be gone through mentally and the desired modification of an expression under treatment can forthwith be written down, (2) and (3) are strictly speaking not required in the working list. Each of the four equations (4) and (5) is similar in the points described to (4a). The reader who desires to acquire this facility could not do better than begin by writing out all the twelve equations implied.

Attend now only to (4) and (5). If the two equations (4) are added together one reaches a mere identity almost immediately. Therefore, either of the equations (4) readily follows from the other. Exactly similarly either of the equations (5) follows from the other. To obtain one of (4) from one of (5), place $\breve\upsilon$ as a post factor to one of (5) and then pass it into the various operators $V_x()$ till it reaches $v$, noting how those operators must be changed in the process. When $\breve\upsilon$ has come to juxtaposition with $v$, change $v\breve\upsilon$ into $w$ so that $n-b=c$ and get rid of $b$, leaving $a$ and $c$. The completion of the verification is simple. I have verified that it is possible to prove either of (5) by reduction to primitive units. I have found no simple proof. Actually I established one of the twelve equations by taking $\rho$ to be one of the vectors $\alpha_c$, say $\alpha_{a+b}$ in (1). Perhaps this is the simplest method.


\indent Because the space I have assigned myself does not permit otherwise, I must now omit several chapters in the MS. treatise from which I have been summarising, but I find one matter requires insertion before I pass to the part of the subject lending itself to applications to differential invariants.

\indent If $\phi$ is a multenion linity its characteristic function $chf_\phi(x)$ is a certain $n$-ic in the arbitrary scalar symbol $x$. Another name for it is that it is the determinant of $x-\phi$, and from the point of view suggested it may be denoted by $|x-\phi|$. In our present notation this is given by the identity
\begin{equation}\left.\begin{aligned}chf_\phi(x)&\equiv|x-\phi|\\&\equiv n'V_n\zeta^{(n)}\cdot V_n[(x-\phi)\zeta]^{(n)}\\&\equiv x^n-m_1x^{n-1}+\cdots+(-1)^nm_n\end{aligned}\right\},\tag{6}\end{equation}
[$n'$ is the same function of $n$ as $a'$ is of $a$ in (5) of §5, that is $n'=(-1)^{\frac12n(n+1)}/n!$; $m_n$ here and in quaternions is generally written as $m$.] Several forms derived from (6) are useful in the present subject. They are expressions for $m_a$ of which it is to be remembered that $m_n$ is a specially important case. There is (1) a fundamental form in $\zeta$, (2) a form involving $n$ given linearly independent vectors $\alpha_1,\alpha_2,\ldots\alpha_n$, (3) a form dependent on the general expression for $\phi$
\begin{equation}\phi\epsilon=-\beta_1V_0\epsilon\alpha_1-\beta_2V_0\epsilon\alpha_2-\cdots=-\sum\beta V_0\epsilon\alpha,\tag{7}\end{equation}
and (4) the particular case of the last for which
\begin{equation}\phi\epsilon=-V_0\epsilon V.\tag{8}\end{equation}
In the order mentioned, these forms are
\begin{equation}\left.\begin{aligned}m_a&=a'V_0(V_a\zeta^{(a)}\phi V_a\zeta^{(a)})\\&=\sum V_n[V_a(\phi\alpha)^{(a)}\cdot V_{n-a}\alpha^{(n-a)}]\cdot V_n^{-1}\alpha^{(a)}\alpha^{(n-a)}\\&=\sum V_0V_a\alpha^{(a)}KV_a\beta^{(a)}\\&=a'V_0(V_a\nabla^{(a)}V_a\epsilon^{(a)})\end{aligned}\right\}.\tag{9}\end{equation}

The last three are all deducible from the first. Another form could at once be written down from the first, involving the vectors $\alpha$ and the vectors $\hat\alpha$. For our applications below the fourth form is important. For this case
\begin{equation}\mathrm F(\zeta,\phi\zeta)=\mathrm F(\nabla_1,\epsilon_1).\tag{10}\end{equation}
From this
\begin{align}V_0\zeta\phi\zeta&=V_0\nabla\epsilon,\qquad V_2\zeta\phi\zeta=V_2\nabla\epsilon,\tag{11}\\|\phi|&=m_n=n'V_0\cdot V_n\nabla^{(n)}V_n\epsilon^{(n)}.\tag{12}\end{align}
[\emph{Added February, 1921.}---The formulae analogous to the quaternion formula for $V_\alpha V_\beta\gamma$ are here greatly generalised and reduced to a form in which they are easy to remember and to apply in particular cases.]

In accordance with (3) of §3 let
\begin{equation}\left.\begin{aligned}\tfrac12[uv+(-1)^{ab}vu]&=(V_{a+b}+V_{a+b-4}+V_{a+b-8}+\cdots)uv=\Lambda uv,\\\tfrac12[uv-(-1)^{ab}vu]&=(V_{a+b-2}+V_{a+b-6}+\cdots)uv=\Lambda'uv\end{aligned}\right\}.\tag{1}\end{equation}
The parts of $\Lambda uv$, and $\Lambda uv$ itself, may be said to be concordant with $V_{a+b}$, those of $\Lambda'uv$ concordant with $V_{a+b-2}$. Similarly, let the parts of $uvw$, as follows,
\begin{equation}\left.\begin{aligned}\tfrac12(uvw+(-1)^{bc+ca+ab}wvu)&=\Lambda uvw,\\\tfrac12(uvw-(-1)^{bc+ca+ab}wvu)&=\Lambda'uvw\end{aligned}\right\}.\tag{2}\end{equation}
be called concordant with $V_{a+b+c}$ and $V_{a+b+c-2}$ respectively. [That $\Lambda uvw$ as defined by (2) is the same as $(V_{a+b+c}+V_{a+b+c-4}+\cdots)uvw$ may be proved from the properties of $K$.]
\[K(uvw)=KwKvKu=(-1)^{\frac12\Sigma a(a+1)}wvu.\]
Hence
\[2\Lambda(uvw)=\{1+(-1)^{\frac12(a+b+c)(a+b+c+1)}K\}uvw.\]
It is important for present purposes to note the various forms of such an expression as $(\Lambda\text{ or }\Lambda')\cdot u(\Lambda\text{ or }\Lambda')vw$, obtained by mere commutations, and especially what are the commutation factors in the various cases. There are two standard cases from which all others may be readily derived. They are
\[\Lambda vu=(-1)^{ab}\Lambda uv,\qquad\Lambda(\Lambda vw\cdot u)=(-1)^{a(b+c)}\Lambda(u\Lambda vw),\]
and we express them by saying that $(-1)^{ab}$ is the commutation factor of $\Lambda uv$, and that $(-1)^{a(b+c)}$ is the commutation factor for $u$ and $\Lambda vw$ in $\Lambda(u\Lambda vw)$. The addition of an accent to any $\Lambda$ clearly multiplies the commutation factor by $-1$. Thus
\[\Lambda'vu=-(-1)^{ab}\Lambda'uv,\qquad\Lambda(\Lambda'vw\cdot u)=-(-1)^{a(b+c)}\Lambda(u\Lambda'vw),\]
\[\Lambda'(\Lambda'vw\cdot u)=(-1)^{a(b+c)}\Lambda'(u\Lambda'vw),\ \text{etc.}\]
Thus, if $x,y$ are independently equal either to $0$ or $1$, and if we put $\Lambda^{(x)}$ for $\Lambda$ or $\Lambda'$ according as $x$ is $0$ or $1$, and similarly for $\Lambda^{(y)}$, we have by commutations four forms of $\Lambda^{(x)}\cdot u\Lambda^{(y)}vw$, namely,
\begin{equation}\Lambda^{(x)}(u\Lambda^{(y)}vw)=(-1)^{y+bc}\Lambda^{(x)}(u\Lambda^{(y)}wv)=(-1)^{x+y+a(b+c)}\Lambda^{(x)}(\Lambda^{(y)}vw\cdot u)=(-1)^{x+bc+ca+ab}\Lambda^{(x)}(\Lambda^{(y)}wv\cdot u).\tag{3}\end{equation}
It is clear that for all present purposes $u,v,w$ may be generalized in meaning. Hitherto they have meant multenions of the forms $V_ap,V_bq,V_cr$; but they may here be extended to mean multenions concordant with $V_a,V_b,V_c$, that is
\begin{equation}\left.\begin{aligned}u&=(V_a+V_{a+4}+V_{a-4}+V_{a+8}+V_{a-8}+\cdots)p,\\v&=(V_b+V_{b+4}+V_{b-4}+\cdots)q,\quad w=(V_c+V_{c+4}+\cdots)r\end{aligned}\right\}.\tag{4}\end{equation}
Let now $Vq$ stand for any $V_xq$, a single selected part, which is concordant with $\Lambda q$; and $V'q$ for any single part which is concordant with $\Lambda'q$. Then the following 8 statements are identities, and each of the 8 has 64 variants obtained by such commutations as we have just been considering. Thus, given $u,v,w$, there are in all 512 variants, and in each of these $V$, or $V'$, has in general a multitude of values. In each of the 8 identities we may add together the individual parts indicated by the $V$ or $V'$, which means that $V$ or $V'$ may in every case be replaced by $\Lambda$ or $\Lambda'$ respectively. For standard reference forms those written seem preferable. The identities are in meaning (not separately, but taken all together), though not in form, symmetrical in $u,v,w$. This will appear in the course of the proof given later.
\begin{equation}
\left.\begin{aligned}
V(\Lambda uv\cdot w)&=V(u\Lambda vw+(-1)^{ab}v\Lambda'uw)\\
 &=V(u\Lambda'vw+(-1)^{ab}v\Lambda uw),\\[3pt]
V(\Lambda'uv\cdot w)&=V(u\Lambda vw-(-1)^{ab}v\Lambda uw)\\
 &=V(u\Lambda'vw-(-1)^{ab}v\Lambda'uw),\\[3pt]
V'(\Lambda uv\cdot w)&=V'(u\Lambda vw+(-1)^{ab}v\Lambda uw)\\
 &=V'(u\Lambda'vw+(-1)^{ab}v\Lambda'uw),\\[3pt]
V'(\Lambda'uv\cdot w)&=V'(u\Lambda vw-(-1)^{ab}v\Lambda'uw)\\
 &=V'(u\Lambda'vw-(-1)^{ab}v\Lambda uw)
\end{aligned}\right\}.\tag{5}
\end{equation}
The following remarks will probably render these easy to remember and apply: (1) Counting the accent of a $V'$ only once when it occurs in an equation, the number of accents in each of the eight equations of three terms is odd, that is, it is 1 or 3. (2) In these standard forms the sequence $u,v,w$ is adopted as far as practicable, that is to say, wholly in the term on the left and the first term on the right; whereas in the second term on the right

$v$ is forced away from the middle position, the sequence $u,w$ is still maintained. (3) When these rules of sequence are adopted the signs to be given to each term are specially easy of determination; when terms on opposite sides of the equation have the same sequence, or by commutation are brought to the same sequence, they have or are brought to have the same sign; and when on the same side to have opposite signs. Thus, as written in (5), the term on the left and the first term on the right, having the same sequence, have the same (plus) sign; and again, the term on the left and the second term on the right are brought to the same sequence by commuting $u$ and $v$ on the left. This last determines the sign, $+(-1)^{ab}$ in the second term on the right.

\indent We will now throw the eight separate cases of (5) into one form, and henceforth replace $V$ and $V'$ by $\Lambda$ and $\Lambda'$. Let $g,x,y,z$ each independently be equal to $0$ or $1$. Then all the cases of (5) are included in
\begin{equation}\left.\begin{aligned}
\Lambda^{(g)}(\Lambda^{(x)}uv\cdot w)&=\Lambda^{(g)}(u\Lambda^{(x)}vw+(-1)^{x+ab}v\Lambda^{(y)}uw)\\
\text{provided}\qquad g+x+y+z&=1\text{ or }3
\end{aligned}\right\}.\tag{6}\end{equation}
Put $x',y',z'$ for the commutation factors of $\Lambda^{(x)}vw$, $\Lambda^{(y)}wu$, $\Lambda^{(z)}uv$ respectively, that is
\begin{equation}x'=(-1)^{x+bc},\qquad y'=(-1)^{y+ca},\qquad z'=(-1)^{z+ab}.\tag{7}\end{equation}
It is easy from (3) to prove that $-y'z'$, $-z'x'$, $-x'y'$ are the commutation factors for $u$ and $\Lambda^{(x)}vw$ in $\Lambda^{(g)}\cdot u\Lambda^{(x)}vw$; and for the like single symbol and coupled pair in $\Lambda^{(g)}\cdot v\Lambda^{(y)}wu$ and $\Lambda^{(g)}\cdot w\Lambda^{(z)}uv$ respectively.

\indent To see that (6) is symmetrical in $u,v,w$, make the commutations necessary to obtain the cyclic order $u,v,w$ throughout; first with the single symbol in the first place throughout, and secondly in the last place throughout. Also put all the terms on one side of the equation and write the three terms in the order for which in the first term $u$ occupies the middle position, in the second $v$ does so, and in the third $w$. Notice that the coefficients $x',y',z'$ thus come in both arrangements to have their alphabetic order. Thus
\begin{equation}\Lambda^{(g)}(x'v\Lambda^{(x)}uw+y'u\Lambda^{(z)}wv+z'v\Lambda^{(y)}wu)=0,\tag{8}\end{equation}
\begin{equation}\Lambda^{(g)}(x'\Lambda^{(y)}wu\cdot v+y'\Lambda^{(z)}wv\cdot w+z'\Lambda^{(x)}uv\cdot u)=0.\tag{9}\end{equation}
To show that these are equivalent to (6), we have only to verify for any two terms in each equation that when brought to the same sequence they have opposite signs, and this is quite easy to do. [Commute $w$ with $\Lambda^{(x)}uv$ in the first term of (8); it comes to agreement in sequence and disagreement in sign with the second term.]

\indent We will now prove (8). In the first term in (8), $-x'y'$ is one commutation factor and $z'$ is the other. Hence four times this term is
\[\begin{aligned}4x'\Lambda^{(g)}w\Lambda^{(z)}uv&=2x'(w\Lambda^{(z)}uv-x'y'\Lambda^{(z)}uv\cdot w)\\&=2(x'w\Lambda^{(z)}uv-y'\Lambda^{(z)}uw\cdot w)\\&=x'\cdot wuv+z'x'\cdot wvu-y'\cdot uvw-y'z'\cdot vuw.\end{aligned}\]
We see by cyclic change from this that four times the second and third terms are
\[y'\cdot uvw+x'y'\cdot uwv-z'\cdot vwu-z'x'\cdot wvu\]
and
\[z'\cdot vwu+y'z'\cdot wuv-x'\cdot uvw-x'y'\cdot uwv.\]
The sum of these three expressions is identically zero, as is seen by inspection.

\indent We can reproduce the equations (5) in other notations. Thus, let
\[vu=Cuv,\qquad wvu=Cuvw,\]
\[\tfrac12(1+C)()=\bar C(),\qquad \tfrac12(1-C)()=\bar C'(),\]
\[\tfrac12(1+K)q=\bar Kq,\qquad \tfrac12(1-K)q=\bar K'q,\]
\[\tfrac12(1+Q)q=\bar Qq,\qquad \tfrac12(1-Q)q=\bar Q'q.\]
Then, similar to (6), the following is true
\begin{equation}\left.\begin{aligned}\bar C^{(g)}(\bar C^{(z)}uv\cdot w)&=\bar C^{(g)}(u\bar C^{(z)}vw+(-1)^{(z)+ab}v\bar C^{(y)}uw)\\\text{provided}\qquad g+x+y+z&=1\text{ or }3\end{aligned}\right\}.\tag{10}\end{equation}
The application of the rule of signs is rather improved, but in most cases the reverse is decidedly the case in the application of the proviso. In (10) \emph{within the brackets} each $C$ may be replaced by $K$, or, instead, by $Q$, but outside the brackets $\bar C^{(g)}$ must be retained; also with $K$ or $Q$ in place of $C$ the index $z$ of $(-1)^z$ must be changed to
\[z+\tfrac12[a(a\pm1)+b(b\pm1)],\]
the upper sign being taken with $K$ and the lower with $Q$.

(10) is more closely analogous than (6) to the quaternion formula for $V_\alpha V_\beta\gamma$, but we must remember that in quaternions
\[\alpha\beta+\beta\alpha=2S\alpha\beta,\qquad \alpha\beta\gamma+\gamma\beta\alpha=2V\alpha\beta\gamma,\]
\[\alpha\beta-\beta\alpha=2V\alpha\beta,\qquad \alpha\beta\gamma-\gamma\beta\alpha=2S\alpha\beta\gamma,\]
so that inside the brackets $\bar C,\bar C'$ are analogous to $S,V$, whereas outside they are analogous to $V,S$.

\indent Of these different forms (5) is without doubt the best for general use, but the others are occasionally convenient.
\[\Lambda-\Lambda'=L\qquad\text{and}\qquad Lu=u=Cu,\]

we have that $C,K,Q$, are all strictly retroscriptive that is
\[Cuvw=CwCvCu,\qquad Kpqr=KrKqKp,\quad\text{etc.},\]
but $L$ is not, though it has the analogous property
\[Luv=(-1)^{ab}LvLu,\qquad Luvw=-(-1)^{bc+ca+ab}LwLvLu.\]
These last properties are easily extended to any number of symbols, such as $u,v,w$. To pass to higher numbers than three with (5), or say (6), is much more troublesome. Here is an example of passing to four symbols. Let $u',v'$ be multenions concordant with $V_{a'},V_{b'}$. In $\Lambda^{(g)}(\Lambda^{(h)}uv\Lambda^{(h')}v'u')$, [where outside the brackets $\Lambda^{(g)}$ is concordant with $V_{a+b+a'+b'-2g}$] first treat $\Lambda^{(h')}v'u'$ as a single symbol, and secondly treat $\Lambda^{(h)}uv$ in that manner. We obtain
\begin{equation}\left.\begin{aligned}
\Lambda^{(g)}(\Lambda^{(h)}uv\Lambda^{(h')}v'u')&=\Lambda^{(g)}(u\Lambda^{(x)}\!\cdot v\Lambda^{(h')}v'u'+(-1)^{h+ab}v\Lambda^{(y)}\!\cdot u\Lambda^{(h')}v'u')\\
&=\Lambda^{(g)}(\Lambda^{(x')}\!\cdot[\Lambda^{(h)}uv\cdot v']u'+(-1)^{h'+a'b'}\Lambda^{(y)}\!\cdot[\Lambda^{(h)}uv\cdot v']v')\\
\text{provided both}\qquad g+(h+h')+(x+y)&\text{ are odd,}\\g+(h+h')+(x'+y')&\text{ are odd}
\end{aligned}\right\}.\tag{11}\end{equation}

In the equation of four terms provided by (11) there would be thirty-two cases corresponding to the eight of (5), because double values are assignable independently to $g,h,h',x,x'$. In each of the equations of three terms provided by (11) there would be sixteen cases.]

\noindent\textsc{§ 8. Integration.}---Let $x_1i_1+x_2i_2+\cdots+x_ni_n=p$ stand for an independent variable position vector of a point in Euclidian space of $n$ dimensions. Let a point trace a curve of which an element is $dp(=dp_1)$. Let this curve move and trace a two-dimensional tract. In the process let an element of the path of the element $dp(=dp_1)$ be $dp_2$. The element of the two-dimensional region traced hereby is given by $V_2dp_1dp_2$. Let this be continued till a $(b+1)$-dimensional tract has been traced.

\indent Let us consider a space integral over the complete boundary (a $b$-tract) of the $(b+1)$-tract and express it as a space integral over the $(b+1)$-tract. To do this, starting from any point of the $(b+1)$-tract, let an elementary closed $b$-tract expand till it has traversed the whole $(b+1)$-tract by reaching its boundary. The process could be described in much more detail. We suppose the $(b+1)$-tract to be eventually filled with parallelepipedal elements. First, for each such element the connection between integral over boundary (of element) is expressed as a multiple of $V_{b+1}dp_1dp_2\ldots dp_{b+1}$, denoting that element of $(b+1)$-tract. These are summed and in this sum contributions from the common boundary of two elements cancel. The direction indicated by $dp_{b+1}$ is to be into the region bounded by the element
\[d\breve p_b=V_bdp_1dp_2\ldots dp_b.\]

\begingroup\normalsize

With these conventions the following is true

\begin{equation}\int_{(b)}\!\int \phi\,d\breve p_b=\int_{(b+1)}\!\int \phi_gV_b\,d\breve p_{b+1}\nabla g,\tag{1}\end{equation}
where
\begin{equation}\nabla=\sum_{c=1}^ni_cD_c.\tag{2}\end{equation}
$d\breve p_b$ is the element of boundary-integral due to the element of boundary $d\breve p_b$, and $d\breve p_{b+1}$ is an element of the $(b+1)$-tract. The form of $\phi$ is a function of position. The suffix $g$ means that $g$ acts on this function $\phi$.

\indent We will now put the integral (1) in a second form. Let $\breve v=U d\breve p_n$ ($U d\breve p_nT d\breve p_n=d\beta_n$, $(T d\breve p_n)^2=\pm(d\breve p_n)^2$, $T d\breve p_n$ a positive real scalar). Also let $d\breve p_b=\breve vds_{n-b}$. $ds_0$ is thus a positive scalar, which measures an element of hyperbulk. $ds_1$ is a vector which measures an element of hyperarea. Write $db$ for $ds_0$ and $da$ for $ds_1$; then
\begin{equation}\left.\begin{aligned}db&=V_0da\,dp_n\\d\breve p_n&=V_nd\breve p_{n+1}dp_n\end{aligned}\right\}.\tag{3}\end{equation}
follows from
\[dp_n=V_ndp_{n+1}dp_n.\]
Since $d\beta$ is positive and $dp_n$ points inwards, $da$ points outwards in harmony with usual convention. Equation (1) in the $d_s$ notation becomes
\begin{equation}\int^{(n-a)}\!\int\psi\,d_s{}_{n-a}=\int^{(n-a+1)}\!\int\psi_gV_a\,d_s{}_{a-1}\nabla g.\tag{4}\end{equation}
The extreme cases of (1) and (4) are when $b=1$ and $a=1$, that is
\begin{equation}\int\phi\,dp=\iint\phi_gV_1d\breve p_2\nabla g,\tag{5}\end{equation}
\begin{equation}\int^{(n-1)}\!\int\psi\,da=\int^{(n)}\!\int\psi_g\nabla g.\tag{6}\end{equation}
In (1) put $\phi d\breve p_b=$ a scalar $=V_0w\,d\breve p_b$ and make similar changes in the other three. Thus (1) and (4) become
\begin{equation}\int^{(b)}\!\int V_0w\,d\breve p_b=\int^{(b+1)}\!\int V_0d\breve p_{b+1}V_{b+1}\nabla w,\tag{7}\end{equation}
\begin{equation}\int^{(n-a)}\!\int V_0wk\,d_s{}_a=\int^{(n-a+1)}\!\int V_0d_s{}_{a-1}V_{a-1}\nabla(kw).\tag{8}\end{equation}
$wk$ has been written in place of a mere $w$ for subsequent convenience. $k$ is supposed to be a given scalar function of position of the nature of density. Thus, if by a general strain $p$ becomes $p'$, $dp$ becomes $dp'$, and $db$ becomes $db'$, then $k$ becomes $k'$ where
\begin{equation}kdb=k'db'.\tag{9}\end{equation}
$k$ is used in the transformation of (6). (5) and (6) become
\begin{equation}\int V_0\sigma\,dp=\iint V_0d\breve p_2V_2\nabla\sigma,\tag{10}\end{equation}
\begin{equation}\int^{(n-1)}\!\int V_0\tau k\,da=\int^{(n)}\!\int V_0\nabla(k\tau)\,db.\tag{11}\end{equation}
We are about to apply these theorems to differential invariants. It may surprise some readers that there is so much concerning invariants which is
\endgroup

independent of any conception of differential quadratic forms. I was acquainted with all the theorems here given to the end of §10 years before I had heard of quadratic differential forms. My first acquaintance with such forms was after the MS. treatise from which I extract them had been completed. That first (and up to this date my only) acquaintance with such forms was through Prof. Wright's delightful Cambridge tract on the subject. How then did they originate? From a study of Maxwell's suggestions concerning intensities and fluxes. Till the last three months I never looked at them from any other point of view. I gave elaborate applications of their three-dimensional form, twenty-seven years ago, to electrical problems.\footnote{\emph{Phil. Trans.}, A, 1892, pp. 685--780.}

\noindent\textsc{§ 9. Covariants and Contravariants.}---Let $p'$ be a vector function of the independent vector $p$. We may think of this as representing a displacement of every point of a medium filling our $n$-dimensional Euclidean space from the position $p$ to the position $p'$, or we may think of it in its pure mathematical aspect as a mere mathematical transformation of the coordinates $x_c=-V_0i_cp$ into the coordinates $x'_c=-V_0i_cp'$.

\indent Corresponding to an arbitrary increment $dp$ of $p$, there is an increment $dp'$ of $p'$ given by
\begin{equation}dp'=-V_0dp\nabla\cdot p'=\chi dp.\tag{1}\end{equation}
We will use $\chi$ to denote the linity extended from the vector linity $\chi$ which is given by (1). Invariance has to do with relations which remain unchanged in form, in some defined sense or senses, when $p'$ is taken as the independent variable instead of $p$. Suppose $h$ is a scalar function of $p$. Then it is also a function of $p'$. Then
\begin{equation}dh=-V_0dp\nabla h=-V_0dp'\nabla'h,\tag{2}\end{equation}
where $\nabla'=\sum_{c=1}^ni_cD_{x'_c}$. This is a relation of the kind just mentioned and we want to systematise the treatment of such relations. Anybody who reads my paper of 1892 will see how it came about that it should seem a very natural course to use the theorems of integration of §8 above for the purpose.

\indent We scarcely need to analyse the details of the strain after what has already been said. The ordinary notions of curl, $V_2\nabla\sigma$, of an intensity, that is a covariant vector, $\sigma$, are given by (10) of §8. The ordinary notions of convergence of a vector flux, $k\tau$, are given by (11) of §8. [$k\tau$ is a vector flux; $\tau$ is a contravariant vector.] The ordinary notions of the half curl, $\tfrac12V_2\nabla\sigma$, representing the rate of rotation of a fluid whose velocity is $\sigma$, are just the same in $n$ dimensions as in three [see (10) of §4 and what has been said above about infinitesimal rotations]. But more:---it cannot surprise

anybody who has seen fluxes and intensities used as a mathematical method that there is no change in the same uses of the same expressions (curl, convergence, etc.) when we pass to an $n$-manifold of the most general kind. $V_2\nabla\sigma$ and $V_0\nabla(k\tau)$ will have just the same significance in the manifold as in the original Euclidean space; but this is anticipating.

\begin{center}\emph{Synopsis of Notation.}\end{center}
\begin{center}\begin{tabular}{|l|c|c|}\hline
& Covariant & Contravariant\\\hline
(1) Vector & $\sigma$ & $\tau$\\\hline
(2) Hypervector & $v$ & $u$\\\hline
(3) Multenion & $s$ & $r$\\\hline
(4) Multenion linity & $\xi$ & $\xi^{-1}$\\\hline
\end{tabular}\end{center}

\indent We have given an instance of an invariantive statement in (2). Next, consider the integrals (10) and (7) of §8, which we here reproduce,
\begin{equation}\left.\begin{aligned}\int V_0\sigma\,d\rho&=\iint V_0d\breve\rho_2V_2\nabla\sigma,\\\int^{(b)}\!\int V_0v\,d\breve\rho_b&=\int^{(b+1)}\!\int V_0d\breve\rho_{b+1}V_{b+1}\nabla v\end{aligned}\right\}.\tag{3}\end{equation}
When $p'$ is the independent variable, precisely these same integrals are to be expressed in the same form, thus
\begin{equation}\left.\begin{aligned}\int V_0\sigma'\,d\rho'&=\iint V_0d\breve\rho'_2V_2\nabla'\sigma',\\\int^{(b)}\!\int V_0v'\,d\breve\rho'_b&=\int^{(b+1)}\!\int V_0d\breve\rho'_{b+1}V_{b+1}\nabla'v'\end{aligned}\right\}.\tag{4}\end{equation}
Now we are compelled to identify our present $dp$ and $d\breve p_b$ with the $dp$ and $d\breve p_b$ of (1), so that $dp'=\chi dp$. Again, $d\breve p_b$ meant $V_bdp_1dp_2\ldots dp_b$, where each $dp_c$ must be identified with the $dp$ of (1), and therefore the corresponding $dp'_c$ with the $dp'$ of (1). Hence
\begin{equation}d\rho'=\chi d\rho,\qquad d\breve{\rho}'_b=\chi d\breve{\rho}_b.\tag{5}\end{equation}
When we say that the integrals (4) are the same as the integrals (3), we assert that each of the four elements of integration in (3) is equal to the corresponding element in (4). In other words, we assert the existence of four invariants,
\begin{equation}\left.\begin{aligned}
V_0\sigma\,d\rho&=V_0\sigma'\,d\rho',&V_0d\breve\rho_2V_2\nabla\sigma&=V_0d\breve\rho'_2V_2\nabla'\sigma',\\
V_0v\,d\breve\rho_b&=V_0v'\,d\breve\rho'_b,&V_0d\breve\rho_{b+1}V_{b+1}\nabla v&=V_0d\breve\rho'_{b+1}V_{b+1}\nabla'v'.
\end{aligned}\right\}.\tag{6}\end{equation}
From (5) and the fundamental property of the conjugate of $\chi$ these at once give
\begin{equation}\left.\begin{aligned}\sigma'&=\chi'^{-1}\sigma,&V_2\nabla'\sigma'&=\chi'^{-1}V_2\nabla\sigma,\\v'&=\chi'^{-1}v,&V_{b+1}\nabla'v'&=\chi'^{-1}V_{b+1}\nabla v\end{aligned}\right\}.\tag{7}\end{equation}

(5) and (7) suggest multenions of types $r,s$, such that
\begin{equation}r'=\chi r,\qquad s'=\chi'^{-1}s.\tag{8}\end{equation}
$r$ and $s$ are multenion functions of $p$, and $r'$ and $s'$ are associated functions of $p'$ which produce the same integrals when $p'$ is taken as the independent variable in place of $p$. When (8) gives the rule of the association of the functions, $r$ and $s$ are said to be contravariant and covariant respectively.

\indent It is well to place here the obvious consequence of (8)
\begin{equation}V_0rs=V_0r's',\tag{9}\end{equation}
that is, $V_0pq$ is always an invariant if one of the two, $p,q$, is covariant and the other contravariant.

\indent By exactly similar reasoning (a little complicated by the presence of $k$), we may treat of equations (11) and (8) of §8. For the sake of saving space, I merely state the results, $k$ being given by (9) of §8, and $\tau,u,d_s$ (including as particular cases $da$ and $db$) being as in (11) and (8) of §8: the following are covariant:---
\[k\breve v,\qquad kdb,\qquad ka,\qquad kds_a;\]
and the following are contravariant
\[kdb,\qquad \tau,\qquad u,\qquad k^{-1}V_0\nabla(k\tau),\qquad k^{-1}V_{a-1}\nabla(ku).\]
$kdb$ is both covariant and contravariant. This merely means that it is invariant ($kdb=k'db'$). Remembering that $\chi h=h=\chi'^{-1}h$ when $h$ is a scalar, (8) shows at once that, for a scalar, covariance and contravariance both mean precisely the same as invariance.

\indent Remembering that $\nabla'h=\chi'^{-1}\nabla h$ when $h$ is a scalar, (8) shows at once that, for a scalar, covariance and contravariance both mean precisely the same as invariance. The persistent manner in which $k$ thrusts itself into these results tempts one to wonder whether covariant and contravariant terminology is superior to intensity and flux terminology. After some experience of using both, I am inclined to believe that the accepted covariant and contravariant standard scheme is superior to the other.

\indent It is not difficult to show that $k\breve v$ satisfies the test of (8) for covariance, that is
\[k'\breve v'=\chi'^{-1}(k\breve v).\]
If in (2) and (3) of §6 we read $\chi$ in place of $\phi$ we get theorems closely connected with multiplication and composition in the absolute differential calculus. Let $v_a,v_b$ be covariant multenions of homogeneities $a,b$. Then
\begin{equation}v_{a+b}=V_{a+b}v_av_b\tag{9}\end{equation}
is covariant (multiplication). If $u_a,u_b$ are similar contravariants, then
\begin{equation}u_{a+b}=V_{a+b}u_au_b\tag{10}\end{equation}
is contravariant (multiplication). If $u$ is contravariant (homogeneity $a$) and $v$ is covariant (homogeneity $b$), then
\begin{equation}u_{a-b}=V_{a-b}uv\tag{11}\end{equation}
is contravariant (composition) when $a>b$. Also
\begin{equation}v_{b-a}=V_{b-a}uv\tag{12}\end{equation}
is covariant (composition) when $b>a$. Lastly, when $a=b$, each is invariant, which is a special case of (9). A special case arises when we put the larger of the two, $a,b$, equal to $n$.

\indent $\xi$, any multenion linity, is said to be covariant when it produces a covariant multenion from a contravariant, thus
\begin{equation}s=\xi r,\qquad s'=\xi'r'.\tag{13}\end{equation}
Conversely since then
\begin{equation}r=\xi^{-1}s,\qquad r'=\xi'^{-1}s',\tag{14}\end{equation}
$\xi^{-1}$ is said to be contravariant. From (8) we at once get the associated transformation laws. From this, when $\xi$ is covariant, its conjugate $\xi'$ is also covariant; therefore both the self-conjugate and skew parts of $\xi$ are covariant. Similarly for $\xi^{-1}$ and contravariance.
\begin{equation}\xi'=\chi'^{-1}\xi\chi^{-1},\qquad \xi'^{-1}=\chi\xi^{-1}\chi'.\tag{15}\end{equation}
From this, when $\xi$ is covariant, its conjugate $\xi'$ is also covariant, and therefore both the self-conjugate part and the skew part of $\xi$ are covariant. Similarly for $\xi^{-1}$ and contravariance. Also, when $\xi$ is self-conjugate, $\chi$ is self-conjugate. If $\xi$ is a covariant vector linity, the extended linity is also covariant. Also, in this case, the rotation-$V_2$ linity $\tfrac12V_2\xi\xi'\xi$ is covariant. To prove this last, we have to show that
\[\chi'V_2\xi\xi'\xi=V_2\xi\xi'\xi.\]
Now
\[\begin{aligned}\chi'V_2\xi\xi'\xi&=V_2\chi'\xi\chi'\xi^{-1}\xi &&\text{[(2) §6]}\\&=V_2\xi\chi'\xi\chi &&\text{[(9) §5]}\\&=V_2\xi\xi'\xi &&\text{[(15) above],}\end{aligned}\]
as required.

\indent Let $\sigma_1,\sigma_2,\ldots,\sigma_n$ be $n$ independent covariant vectors. [They might be taken as $\nabla h_1,\nabla h_2,\ldots,\nabla h_n$ where the Jacobian of the quantities $h$ does not vanish, that is, where $V_n(\nabla h)^{(n)}\ne0$.] From the relations
\begin{equation}\left.\begin{aligned}1&=V_0\sigma_1\hat{\sigma}_1=V_0\sigma_2\hat{\sigma}_2=\cdots,\\0&=V_0\sigma_1\hat{\sigma}_2=V_0\sigma_1\hat{\sigma}_3=\cdots\end{aligned}\right\},\tag{16}\end{equation}
it at once follows that $\hat{\sigma}_1,\hat{\sigma}_2,\ldots,\hat{\sigma}_n$ are $n$ independent contravariant vectors, and $\sigma_c,\hat{\sigma}_c$ might just as properly have been denoted by $\hat{\tau}_c,\tau_c$. Any covariant or contravariant vector can be expressed at once in terms of these and invariants. Also, from the $2^n$ vectoriums formed from each set,
any covariant or contravariant multenion may be similarly expressed. Thus the expressions for $\sigma,\tau,v,u$ are [see (2) of §5]
\begin{equation}\left.\begin{aligned}\sigma&=\sum_{c=1}^n\sigma_cV_0\sigma\hat{\sigma}_c,&\tau&=\sum_{c=1}^n\hat{\sigma}_cV_0\tau\sigma_c,\\v&=\sum V_b\sigma^{(b)}\cdot V_0v\,\hat{\sigma}^{(b)},&u&=\sum V_a\hat{\sigma}^{(a)}\cdot V_0uV_a\sigma^{(a)}\end{aligned}\right\}.\tag{17}\end{equation}
\noindent\textsc{§ 10. Differentiation.}---Our integration theorems have suggested certain systematisations, now to be given in the use of $\nabla$. We have
\[\nabla w=V_{c+1}\nabla w+V_{c-1}\nabla w.\]
If $w$ is covariant, one part on the right, viz., the first, is covariant and the other not. A similar statement applies to contravariance. We cannot, as seems desirable, resolve the symbolic vector $\nabla$ into two parts, to which each of these two discordant parts on the right belong. But we can so resolve the symbolic linity $\nabla()$. The first part of the linity, connected with $V_{c+1}\nabla w$ produces a part of $\nabla w$ of homogeneity one unit higher than that of $w$, and may be thought of as the ascending derivative linity, and will be denoted by $H$. The other part, $L$, produces a lower homogeneity, and may be thought of as the descending derivative. Thus, formally, we define the linities $H$ and $L$ thus
\begin{equation}\left.\begin{aligned}\nabla()&=H+L\\H&=\sum_{c=0}^{n-1}V_{c+1}\cdot\nabla V_c(),&L&=\sum_{c=1}^{n}V_{c-1}\cdot\nabla V_c()\end{aligned}\right\}.\tag{1}\end{equation}
Otherwise
\begin{equation}\left.\begin{aligned}\nabla w&=Hw+Lw\\Hw&=V_{c+1}\nabla w,&Lw&=V_{c-1}\nabla w\end{aligned}\right\}.\tag{2}\end{equation}
There are two, wholly separate, aspects of $H$ and $L$ as operators. Each is both a linear operator and a differential operator. In (2) the subject of these two operations is the one symbol $w$, but in (4), below (where $H'$ is the conjugate of $H$), in $\phi_gH'_gd\breve\rho_{b+1}$, $H'$ as linear operator affects the immediately succeeding symbol $d\breve\rho_{b+1}$, while as differential operator it affects $\phi$, and this is indicated by suffixes in the usual manner. This example illustrates some of the facilities of operation provided by $H$.


\indent In considering the fundamental properties of $H$ and $L$, think of replacing
\[
\begin{array}{r@{\qquad}c@{\qquad}c@{\qquad}c}
  \nabla & \nabla w & V_{c+1}\nabla w & V_{c-1}\nabla w \\
  \text{by}\quad\alpha & \alpha w & V_{c+1}\alpha w & V_{c-1}\alpha w,
\end{array}
\]
where $\alpha$ is an ordinary vector instead of a symbolic vector. Let $H',L'$

be the conjugates of $H,L$ as usual. From the fundamental relation,
\[
V_a p\phi q=V_a q\phi'p,
\]
we at once have
\begin{equation}H'w=V_{c-1}w\nabla,\qquad L'w=V_{c+1}w\nabla.\tag{3}\end{equation}
Thus the integration theorems (1) and (4) of §8 read
\begin{align}
\int^{(b)}\!\int\phi\,d\breve\rho_b
  &=\int^{(b+1)}\!\int\phi_gH'_g\,d\breve\rho_{b+1},\tag{4}\\
\int^{(n-a)}\!\int\psi\,ds_a
  &=\int^{(n-a+1)}\!\int\psi_gL'_g\,ds_{a-1},\tag{5}
\end{align}
and (7) and (8) of §8 read
\begin{align}
\int^{(b)}\!\int V_0v\,d\breve\rho_b
  &=\int^{(b+1)}\!\int V_0d\breve\rho_{b+1}Hv,\tag{6}\\
\int^{(n-a)}\!\int V_0uk\,ds_a
  &=\int^{(n-a+1)}\!\int V_0ds_{a-1}L(ku).\tag{7}
\end{align}
Further we have
\begin{equation}H^2=0=L^2,\tag{8}\end{equation}
and by taking conjugates we have
\[
H'^2=0=L'^2.
\]
From (8) consider the various powers $\nabla^2(),\nabla^3()$ of $\nabla()$. Clearly we get
\begin{equation}\left.\begin{aligned}
\nabla^2()&=HL+LH,\\
\nabla^3()&=HLH+LHL,\ \text{etc.}
\end{aligned}\right\},\tag{9}\end{equation}
and in (9) we may take conjugates [noting that, since the conjugate of $\nabla()$ is $()\nabla$, the conjugate of $\nabla^2()$ is $()\nabla^2$, etc.].

\indent Notice that $H'$ \emph{lowers} the homogeneity, while $L'$ \emph{heightens} it, so that we have the following Table:---
\begin{center}
\begin{tabular}{|l|c|c|}\hline
& \begin{tabular}{c}Heightens\\homogeneity.\end{tabular}
& \begin{tabular}{c}Lowers\\homogeneity.\end{tabular}\\\hline
Prefixes $\nabla\ldots\ldots\ldots\ldots\ldots\ldots$ & $H$ & $L$\\
Postfixes $\nabla\ \ldots\ldots\ldots\ldots\ldots\ldots$ & $L'$ & $H'$\\\hline
\end{tabular}
\end{center}

\indent Either of the two, $H$ and $L$, can be expressed in terms of the other thus
\begin{equation}Lq=H(q\breve{\upsilon})\cdot\breve{\upsilon}^{-1},\qquad Hq=L(q\breve{\upsilon})\cdot\breve{\upsilon}^{-1}.\tag{10}\end{equation}
Certain identities satisfied by $H$ and $L$ are given by the theorems of integration. As an example, $\int dp'=0$ for a closed curve. This can be expressed in terms of $p$ thus:
\[
0=\int\chi\,dp=\iint\chi_gV_1d\breve p_2\nabla_g.
\]
Since the whole tract may be taken as the element $d\breve p_2$ we get
\begin{equation}\chi_gV_1d\breve p_2\nabla_g=0.\tag{11}\end{equation}

If we operate by $V_0\gamma()$ we get
\[
V_0d\breve p_2V_2\nabla\chi'\gamma,
\]
or since $d\breve p_2$ is arbitrary,
\begin{equation}V_2\nabla\chi'\gamma=0.\tag{12}\end{equation}
Both (11) and (12) are simple enough to be established directly from the form of $\chi$, namely, $\chi\gamma=-V_0\gamma\nabla\cdot p'$, but the method of deriving them, just given, suggests how to prove other results not nearly so easy to establish directly. Instead of $\int dp'=0$ we may take either of
\[
\int^{(b)}\!\int d\breve p'_b=0,\qquad
\int^{(n-a)}\!\int ds'_a=0,
\]
and modify them in the same way. Thus from the first,
\[
0=\int^{(b+1)}\!\int\chi_gH'_g\,d\breve p_{b+1}.
\]
From this and by taking conjugates we get
\begin{equation}\chi_gH'_gq=0,\qquad H_g\chi'_gq=0,\tag{13}\end{equation}
where $q$ is an arbitrary multenion. If we put $q=w$ (13) becomes
\begin{equation}\chi_gV_{c-1}w\nabla_g=0,\qquad V_{c+1}\nabla_g\chi'_gw=0.\tag{14}\end{equation}
Similarly we have
\[
0=\int^{(n-a)}\!\int ds'_a
 =\int^{(n-a)}\!\int m\chi'^{-1}\,ds_a.
\]
$m$ is here $k/k'$, \emph{i.e.}, it is the Jacobian of $p'$, or by (11) of §7
\[
m=n'V_0V_n\nabla^{(n)}V_n\epsilon^{(n)}.
\]
Thus
\[
0=\int^{(n-a+1)}\!\int(m\chi'^{-1})_gL'_g\,ds_{a-1}.
\]
Hence, and by taking conjugates
\begin{equation}(m\chi'^{-1})_gL'_gq=0,\qquad L_g(m\chi'^{-1})_gq=0.\tag{15}\end{equation}
Putting $q=w$ we get
\begin{equation}(m\chi'^{-1})_gV_{c+1}w\nabla_g=0,\qquad V_{c-1}\nabla_g(m\chi'^{-1})_gw=0.\tag{16}\end{equation}

\indent Let us now find expressions explicitly giving $\nabla',H',L'$ in terms involving $\nabla,H,L$. By (2), §(9), $\nabla$ is a symbolic covariant vector, so that
\begin{equation}\nabla'=\chi'^{-1}\nabla,\tag{17}\end{equation}
where on the right the differentiations are not to affect $\chi'^{-1}$.

\indent For $H$ we have that $Hs$ is covariant when $s$ is. Hence
\[
H's'=\chi'^{-1}Hs=\chi'^{-1}H(\chi's').
\]
Now, by the second of (13) in this expression we may show the differentiations of $H$ explicitly as affecting or as not affecting the immediately following $\chi'$, because the terms implied by those differentiations are identically zero. Hence we may write
\begin{equation}H'=\chi'^{-1}\cdot H\cdot\chi',\tag{18}\end{equation}
bearing this permissibility of interpretation in mind. The relations between
$H'$ and $\chi$, on the one hand, and $H'$ and $\chi^{-1}$, on the other, are precisely the same, so that we may write
\begin{equation}H=\chi'\cdot H'\cdot\chi'^{-1},\tag{19}\end{equation}
with permissibility of interpretation as in (18).

\indent Similarly from (18) and (19) by means of (5) of §7
\begin{align}
L'&=(m^{-1}\chi)\cdot L\cdot(m\chi^{-1}),\tag{20}\\
L&=(m\chi^{-1})\cdot L'\cdot(m^{-1}\chi),\tag{21}
\end{align}
with the same permissible interpretation as to the incidence of the application of the differentiations of $L$ and $L'$ on the right; they may be supposed to affect the immediately following $(m\chi^{-1})$ and $(m^{-1}\chi)$ or not, at our pleasure.

\noindent\textsc{§ 11. \emph{The fundamental covariant vector linity, $\eta$.}}---The meaning of a manifold as based on a quadratic differential form is here taken for granted, and this, whether, in the ordinary terminology, all or only some of its dimensions are real. Let
\begin{equation}\left.\begin{aligned}
\iota_a&=i_a, &&(a=1,2,\ldots,c),\\
\iota_b&=i_b\sqrt{-1}, &&(b=c+1,c+2,\ldots,n),\\
\rho&=\sum_{h=1}^n x_h\iota_h.
\end{aligned}\right\}\tag{1}\end{equation}
$c$ having any given value from $1$ to $n$. In the present paper we shall deal with multenions and multenion linities based on a system of vectors $\rho$ for which all the co-ordinates $x_h$ and the similar scalars are real. In the case of relativity $n=4$, $c=3$.

\indent The fundamental quadratic form is taken to be
\begin{equation}ds^2=V_0d\rho\eta d\rho=V_0d\rho'\eta'd\rho',\tag{2}\end{equation}
the sign which is given to $ds^2$, a scalar, being convenient in applications to relativity, but inconvenient in applications to pure geometry. $\eta$, a self-conjugate vector linity, is a given function of position, $\rho$; $\rho'$, a vector, an arbitrary function of $\rho$; and since the transformation from $\eta$ to $\eta'$ is given by (2) the fundamental linity $\eta$ is said to be covariant. [$\eta$ and $ds^2$ are real, but sometimes $ds$ and sometimes $ds\sqrt{-1}$ is real; when $ds$ is real it is ``a time-like interval''; when $ds\sqrt{-1}$ is real it is a length; when the question of reality is left undecided $ds$ is an interval.] At any given position $\rho'$ may be so chosen that $\eta'd\rho'=d\rho'$, that is $\eta'=1$. This imposes that $\eta$ is without singularities; it has $n$ mutually orthogonal principal directions and all its roots are real and positive, excluding zero values. Thus $\eta^{\frac12}$ is assumed to have the same orthogonal system and its roots are taken to be the positive values of the square roots of $\eta$ so that $\eta^{\frac12}$ is real and without ambiguity.

\indent We may, and from here on shall, take our previously introduced scalar $k$ (§8) to be $|\eta^{\frac12}|$, the determinant of $\eta^{\frac12}$. [We have
\[
\eta=\chi'\eta'\chi
\]
so that
\[
|\eta|=|\eta'|\times|\chi|^2=|\eta'|\times(k/k')^2,
\]
because $k/k'$ was defined by
\[
kdb=k'db'\qquad\text{or}\qquad k/k'=|\chi|.\quad\text{]}
\]

It will be seen that when, in the usual notation of relativity, it is said that with Galilean co-ordinates
\[
(g_{11},g_{22},g_{33},g_{44})=(-1,-1,-1,+1),
\]
we have not that $\eta(\iota_1,\iota_2,\iota_3,\iota_4)=(-\iota_1,-\iota_2,-\iota_3,+\iota_4)$ but instead the simpler statement that $\eta=\eta^{\frac12}=1$ or
\[
\eta(\iota_1,\iota_2,\iota_3,\iota_4)=(\iota_1,\iota_2,\iota_3,\iota_4),
\]
for it will be observed that this gives
\[
ds^2=-dx_1^2-dx_2^2-dx_3^2+dx_4^2.
\]
\indent A little care is required in expressing a linity, its conjugate, reciprocal, etc., by scalars or co-ordinates. It seems sufficient to illustrate from vector linities. When we say that $\phi$ is given by the square array of $16$ scalars
\[
\left|\begin{array}{cccc}a&b&c&d\\e&f&g&h\\l&m&n&p\\q&r&s&t\end{array}\right|
\]
we imply that the co-ordinates of $\phi\iota_1$ are found in the first \emph{column} thus
\[
\phi\iota_1=a\iota_1+e\iota_2+l\iota_3+q\iota_4
\]
and so for the other columns. In the square array of $n^2$ scalars specifying $\phi$, let $a_{xy}$ be the scalar in the $x$th row and $y$th column. Thus $\phi\iota_y=\sum_{x=1}^n a_{xy}\iota_x$, or more conveniently
\begin{equation}a_{xy}=V_0\iota_x^{-1}\phi\iota_y.\tag{3}\end{equation}
Except when $c=n$ the square array of a self-conjugate linity is not symmetrical about the principal diagonal. Let $a_{xy},a'_{xy}$ be the square arrays of $\phi$ and its conjugate $\phi'$. By the meaning of conjugate we have
\[
V_0\iota_x\phi\iota_y=V_0\iota_y\phi'\iota_x.
\]
Thus instead of $a'_{yx}=a_{xy}$ (as we have when $c=n$) we have by (3)
\begin{equation}a'_{yx}=\iota_x^2\iota_y^2a_{xy},\tag{4}\end{equation}
and in particular, when $\phi$ is self-conjugate $a_{yx}=\iota_x^2\iota_y^2a_{xy}$. If we would retain the $16$ scalars $g_{ab}$ with their usual meanings as in treatises on relativity we may use $g_{(ab)}$ for the corresponding scalars specifying $\eta$. Thus from the fact that $ds^2=\sum g_{ab}dx_adx_b$ and from (3)
\[
g_{ab}=V_0\iota_a\eta\iota_b,\qquad g_{(ab)}=V_0\iota_a^{-1}\eta\iota_b,
\]
whence
\begin{equation}g_{ab}=\iota_a^2g_{(ab)}.\tag{5}\end{equation}
More details with regard to these notations and the theory of tensors will be found below in §17.

\indent \emph{Note on Notations in the Present Paper.}---There will be many exceptions below to the conventions here to be described. Especially do exceptions occur in very frequently recurring symbols. Thus none of the following receive the ``cross'' notation although they should do according to the conventions,
\[
\lambda,\ E_a,\ \bar E_a,\ C_\omega,\ C.
\]
$\lambda$ is a covariant vector, and of the other four $C_\omega,C$ are contramixt linities while $E_a,\bar E_a$ are portions of contramixt linities. The explanations of the technical terms will be made as they come to be used below.

\indent (1) Contravariant multenions will have no distinguishing mark. Covariant multenions will usually be distinguished by a ``cross'' (which, because of its position, is scarcely likely to be confused with a sign of multiplication), thus
\[
\alpha^\times,\ u^\times,\ p^\times,\ q^\times,
\]
read ``alpha cross,'' etc.

\indent (2) Covariant and comixt linities will have no distinguishing mark. Contravariant and contramixt linities will usually be distinguished by a cross.

\indent (3) Contravariant and covariant linities will be denoted by Greek letters; contramixt and comixt linities by Roman capitals.

\indent (4) $\alpha,\beta,\gamma,\alpha^\times,\beta^\times,\gamma^\times$, will generally denote ``dummy'' vectors. (Eddington uses ``dummy'' in an analogous sense for a positive integer.) Dummies are always (below, though not of necessity) treated as constants in that they are not affected by the differentiations of $\nabla$.

\indent The duty of a dummy contravariant or covariant multenion may be generally described as ``to occupy, and thereby indicate, a situation into which an actual, previously defined, contravariant or covariant multenion, may legitimately be inserted.''

\noindent\textsc{§ 12. \emph{The Invariant Forms of any Multenion Expression.}}---$\eta$ has been defined as a fundamental vector linity. It is also used for the corresponding extensive linity, a multenion linity; and similarly for $\eta^{\frac12}$. Since in (2), of §11, $ds^2=(\eta^{\frac12}d\rho)^2$, $\eta^{\frac12}d\rho$ in our manifold takes the place of the $d\rho$ of multenions in a Galilean manifold. [In a Euclidean manifold all the dimensions are real. In a Galilean manifold some are imaginary.] More generally, at a given position of the manifold, $\eta^{\frac12}q$ where $q$ is a contravariant multenion, or $\eta^{-\frac12}q^\times$ where $q^\times$ is a covariant multenion, takes the place of $q$ in a Galilean manifold. $\eta^{\frac12}q$ and $\eta^{-\frac12}q^\times$ will be spoken of as the neutral forms of $q$ and $q^\times$. Indeed, the three $q$ (contravariant), $\eta q$ (covariant) and $\eta^{\frac12}q$ (neutral), are to be regarded as only different mathematical representations of the same intrinsic property, of the manifold; just as in an ordinary dynamical system, the whole system of velocities or instead the whole system of momenta equally will serve to specify the instantaneous motion.

Corresponding to any multenion expression such as $V_a\cdot pqV_b(rs)$ which has a Galilean interpretation there are always the three forms ($p,q,r,s$ now being contravariant) in our manifold (contra., co., neut.).
\[
\eta^{-\frac12}V_a\cdot[\eta^{\frac12}p\eta^{\frac12}qV_b(\eta^{\frac12}r\eta^{\frac12}s)],\qquad \eta^{\frac12}V_a\cdot[\ ],\qquad V_a\cdot[\ ].
\]
The first and second of these three can \emph{always} be expressed in such a shape that only $\eta$ and not $\eta^{\frac12}$ occurs explicitly; though the expression may be very complicated as in the example chosen. It is these final forms from which $\eta^{\frac12}$ has been explicitly banished, that are regarded as invariantive.

\indent To show this, first note as an example, that the product $pqr$ can be first modified by expressing $qr$ as a sum of terms such as $V_{a+b-2c}\cdot V_aqV_br$ and then $p\cdot qr$ can be similarly modified. We have then merely to find the invariant forms (two, namely, covariant and contravariant) of $V_{a+b-2c}uv$. [The reader is reminded that in our notation $u,v,w$, have homogeneities $a,b,c$, respectively, or have the forms $u=V_ap,\ v=V_bq,\ w=V_cr$.]

\indent Let $a$ be greater or equal to $b$ and let $u,v$ be contravariant. Then most frequently all we require to note are the two cases $c=0,c=b$:
\[
\eta^{-\frac12}V_{a+b}\eta^{\frac12}u\eta^{\frac12}v=V_{a+b}uv,\qquad \eta^{\frac12}V_{a+b}\eta^{\frac12}u\eta^{\frac12}v=V_{a+b}\eta u\eta v,
\]
\[
\eta^{-\frac12}V_{a-b}\eta^{\frac12}u\eta^{\frac12}v=V_{a-b}u\eta v,\qquad \eta^{\frac12}V_{a-b}\eta^{\frac12}u\eta^{\frac12}v=V_{a-b}\eta u v.
\]
and here of course we may have $u=\eta^{-1}u^\times$ or $v=\eta^{-1}v^\times$ in each of the forms.

\indent To deal with the general value of $c$ in $V_{a+b-2c}uv$ we use $\zeta$. Now we know that if $\tau_1,\tau_2,\ldots\tau_n$ are $n$ independent contravariant vectors at a point, then their normal reciprocals $\hat{\tau}_1,\hat{\tau}_2,\ldots\hat{\tau}_n$ are $n$ independent covariant vectors. Also we know that if $(\alpha,\beta^\times)$ is any expression---bilinear in two vectors
\[
(\zeta,\zeta)=-\sum_{a=1}^n(\tau_a,\hat{\tau}_a).
\]
Hence if we insert $\zeta,\zeta$ into two places in an invariant form, of which one is occupied by a contravariant vector and the other by a covariant vector, the resulting expression will be invariantive (meaning, ``of invariant form''). Now it is easy to see by reduction to primitive units that
\begin{equation}V_{a+b-2c}uv=[(-1)^{\frac12c(c+1)}/c!]V_{a+b-2c}\cdot V_{a-c}u\zeta^{(c)}V_{b-c}\zeta^{(c)}v,\tag{1}\end{equation}
where as usual $\zeta^{(c)}$ stands for $\zeta_1\zeta_2\ldots\zeta_c$ and $\zeta_c^{(c)}$ for $V_c\zeta^{(c)}$. Hence the contravariant form of $V_{a+b-2c}uv$ is
\begin{equation}[(-1)^{\frac12c(c+1)}/c!]V_{a+b-2c}\cdot V_{a-c}u\zeta_c^{(c)}V_{b-c}\zeta_c^{(c)}v,\tag{2}\end{equation}
from which the covariant form may be written down, and either $u^\times=\eta u$ or $v^\times=\eta v$ or both may be used in place of $u$ and $v$. In my MS. work I


am in the habit of using $\zeta_K^{(c)},\zeta_Q^{(c)}$ for $K\zeta_c^{(c)},Q\zeta_c^{(c)}$. Thus (2) would be $(c!)^{-1}V_{a+b-2c}\cdot V_{a-c}u\zeta_c^{(c)}V_{b-c}\eta\zeta_K^{(c)}v$. A special case of (2) often required is that the contravariant form of $V_2\omega\omega'$ where $\omega$ and $\omega'$ are contravariant and ${}_nV_2$ is
\begin{equation}V_2(V_1\zeta\omega)(V_1\eta\zeta\omega').\tag{3}\end{equation}
or more generally, for infinitesimal rotations generally, the contravariant form of $V_a\omega\iota$ is
\begin{equation}V_a(V_1\zeta\omega)(V_{a-1}\eta\zeta\iota).\tag{4}\end{equation}
\noindent\textsc{§ 13. \emph{Normal and Incident Components in Invariant Form.}}---These components occur very frequently in relativity. Thus the energy-momentum linity is density multiplied by an incident component, at any rate in its kinetic or primitive form. Normal and incident components serve to separate: (1) vector potential from electrostatic potential; (2) momentum from energy; (3) force from power; (4) current density from charge density; (5) magnetic induction from electrostatic force; (6) magnetic field strength from displacement. One should always be prepared to recognise these resolutions.

\indent In quaternions if $\alpha$ is a given vector and $\sigma$ an arbitrary one, $\sigma$ has two components with reference to $\alpha$, namely the normal component
\[
N_0\sigma=\alpha^{-1}V\alpha\sigma=\alpha V\alpha\sigma/\alpha^2
\]
and the incident component
\[
I_0\sigma=\alpha^{-1}S\alpha\sigma=\alpha S\alpha\sigma/\alpha^2.
\]
As will be familiar to users of the quaternion method, it is these components each multiplied by $\alpha^2$, rather than the components themselves, which we shall expect to meet with frequently in analytical expressions. The two
\[
N_\alpha\sigma=N\sigma=\alpha V\alpha\sigma,\qquad I_\alpha\sigma=I\sigma=\alpha S\alpha\sigma,
\]
may be referred to as the normal and incident \emph{expressions} of $\sigma$ with reference to $\alpha$. With regard to these four linities $N_0,I_0,N,I$ we have as follows
\begin{equation}\left.\begin{aligned}N_0^2&=N_0,&I_0^2&=I_0,&N_0I_0&=0=I_0N_0,\\N_0'&=N_0,&I_0'&=I_0,&N_0+I_0&=1=N_0'+I_0'\end{aligned}\right\},\tag{1}\end{equation}
\begin{equation}\left.\begin{aligned}N^2&=\alpha^2N,&I^2&=\alpha^2I,&NI&=0=IN,\\N'&=N,&I'&=I,&N+I&=\alpha^2=N'+I'\end{aligned}\right\}.\tag{2}\end{equation}
Similarly in a Galilean manifold any multenion $q$ has two components $N_0q,I_0q$ normal and incident to a given vector $\alpha$. Let
\[
q=\sum_{c=1}^n w,\qquad\text{where}\quad w=V_cq;
\]

then when the linities $Nw,N_0w$, etc., have been defined, $Nq,N_0q$ have also been defined. We have
\begin{equation}\left.\begin{aligned}Nw&=\alpha V_{c+1}\alpha w,&Iw&=\alpha V_{c-1}\alpha w,\\N_0w&=\alpha^{-1}V_{c+1}\alpha w,&I_0w&=\alpha^{-1}V_{c-1}\alpha w\end{aligned}\right\}.\tag{3}\end{equation}
[If $\epsilon$ is a vectorium of homogeneity $b$, a vector $\sigma$ has a normal component $\epsilon^{-1}V_{b+1}\epsilon\sigma$ and an incident component $\epsilon^{-1}V_{b-1}\epsilon\sigma$ with reference to $\epsilon$. With reference to $\epsilon$, $w$ has many, not merely two, components given by
\[
w=\epsilon^{-1}V_{b+c}\epsilon w+\epsilon^{-1}V_{b+c-2}\epsilon w+\cdots.
\]
The terms on the right of this equation are all, separately, but not obviously, of homogeneity $c$.]

\indent Next let $\alpha$ be a given contravariant vector at a given position $\rho$ of our manifold and $q$ an arbitrary contravariant multenion at the same position and let $w=V_cq$. Also let $\alpha_v,q_v,w_v$ be the neutral forms and $\alpha^\times,q^\times,w^\times$ the covariant forms of $\alpha,q,w$ respectively, that is to say,
\begin{equation}\left.\begin{aligned}\eta^{\frac12}\alpha&=\alpha_v=\eta^{-\frac12}\alpha^\times,\\\eta^{\frac12}w&=w_v=\eta^{-\frac12}w^\times,\\\eta^{\frac12}q&=q_v=\eta^{-\frac12}q^\times\end{aligned}\right\}.\tag{4}\end{equation}
The invariant form of $\alpha^2$ is $\alpha_v^2$ or
\begin{equation}V_0\alpha\eta\alpha=\alpha_v^2=V_0\alpha^\times\eta^{-1}\alpha^\times.\tag{5}\end{equation}
With our present meanings (3) is not invariant in form. There are four modes of choosing the invariant forms of $N$ and $I$, $N_0$ and $I_0$ and generally of any multenion linities, (1) covariant, (2) contravariant, (3) comixt, (4) contramixt. We will take number (3) of these for the basis of the notation in connection with $N$ and $I$. A comixt linity means one which acting on a covariant multenion produces a covariant multenion. For this type of linity it is obvious that the invariant forms of (3) are given by
\begin{equation}\left.\begin{aligned}
Nw^\times&=V_c\alpha V_{c+1}\eta\alpha w^\times=V_c\eta^{-1}\alpha^\times V_{c+1}\alpha^\times w^\times,\\
Iw^\times&=V_c\eta\alpha V_{c-1}\alpha w^\times=V_c\alpha^\times V_{c-1}\eta^{-1}\alpha^\times w^\times,\\
N_0w^\times&=Nw^\times/V_0\alpha\eta\alpha=Nw^\times/V_0\alpha^\times\eta^{-1}\alpha^\times,\\
I_0w^\times&=Iw^\times/V_0\alpha\eta\alpha=Iw^\times/V_0\alpha^\times\eta^{-1}\alpha^\times
\end{aligned}\right\}.\tag{6}\end{equation}
(2) comes with this notation to be modified in two ways only; $\alpha^2$ is replaced by $V_0\alpha\eta\alpha$ or $V_0\alpha^\times\eta^{-1}\alpha^\times$ and in place of $N'$ and $I'$ we have a new kind of conjugate $\eta N'\eta^{-1},\eta I'\eta^{-1}$ involving the fundamental linity $\eta$. Thus we now have
\begin{equation}\left.\begin{aligned}
N^2&=V_0\alpha\eta\alpha\cdot N,&I^2&=V_0\alpha\eta\alpha\cdot I,&NI&=0=IN,\\
\eta N'\eta^{-1}&=N,&\eta I'\eta^{-1}&=I,&N+I&=V_0\alpha\eta\alpha=\eta(N'+I')\eta^{-1}
\end{aligned}\right\}.\tag{7}\end{equation}
[Note that if we put $\eta\phi'\eta^{-1}=\phi_c$ we obtain for this new conjugate the familiar properties
\[
(\phi+\psi)_c=\phi_c+\psi_c,\qquad(\phi\psi)_c=\psi_c\phi_c,\qquad(\phi_c)_c=\phi.]
\]
The types of linity we have to recognise are given in the following short list.
\begin{center}\textsc{Linity Types.}\end{center}
\begin{center}\begin{tabular}{|l|c|c|c|c|c|}\hline
Name & Co-variant. & Contra-variant. & Co-mixt. & Contra-mixt. & Neutral.\\\hline
Symbol & $\xi$ & $\xi^\times$ & $X$ & $X^\times$ & $X_v=\phi$\\
Meaning & $p^\times=\xi q$ & $p=\xi^\times q^\times$ & $p^\times=Xq^\times$ & $p=X^\times q$ & $p_v=X_vq_v$\\\hline
\end{tabular}\end{center}
Thus $\xi$ converts a $q$ (contravariant) into a $p^\times$ (covariant) and similarly for the other types.

\indent The two mixed linities convert a multenion of a given type to one of the same type; the other two convert a given type to the \emph{other} type. [In our present method the mixed linities could scarcely be more unhappily named.]
\indent Corresponding to any one of the types we have, of course, by means of $\eta$, symbols of the other types with the same intrinsic significance. Thus, corresponding to the neutral form $\phi$ we have
\begin{equation}\left.\begin{aligned}\xi&=\eta^{\frac12}\phi\eta^{\frac12},&\xi^\times&=\eta^{-\frac12}\phi\eta^{-\frac12},\\X&=\eta^{\frac12}\phi\eta^{-\frac12},&X^\times&=\eta^{-\frac12}\phi\eta^{\frac12}\end{aligned}\right\}.\tag{8}\end{equation}
From the list and from (8) the following statements, in which $\xi_1,\xi_2$ mean two linities of the type $\xi$, etc., are easily established:---
\par\indent (1) $\xi_1\xi_2$ has no invariant significance, nor has $\xi_1^\times\xi_2^\times$.
\par\indent (2) $X_1X_2$ is of type $X$ and $X_1^\times X_2^\times$ is of type $X^\times$.
\par\indent (3) Although by (1), for integral values of $c$, $\xi^c$ has in general no invariant significance; in particular when $c=-1$ it has significance. $\xi^{-1}$ is of type $\xi^\times$, and $\xi^{\times-1}$ is of type $\xi$. $X^{-1}$ is of type $X$, and $X^{\times-1}$ is of type $X^\times$.
\par\indent (4) A rational function $f(X)$ is of type $X$; and $f(X^\times)$ is of type $X^\times$; the corresponding neutral form in each case being $f(\phi)$.
\par\indent (5) $\xi'$ is of type $\xi$, $\xi^{\times\prime}$ of type $\xi^\times$; so that a similar remark applies to the self-conjugate and skew parts of $\xi$ and $\xi^\times$. $\xi'$ and $\xi^{\times\prime}$ correspond to $\phi'$.
\par\indent (6) $X'$ is of type $X^\times$, $X^{\times\prime}$ of type $X$, both corresponding to $\phi'$. Again $\eta X'\eta^{-1}$ is of type $X$ and $\eta^{-1}X^{\times\prime}\eta$ of type $X^\times$, both corresponding to $\phi'$.
\par\indent (7) The identity satisfied by $\xi$ or by $\xi^\times$ has no significance. The
identity satisfied by $\phi$ is the same as the identity satisfied by each of the following
\[
X=\xi\eta^{-1}=\eta\xi^\times,\qquad X^\times=\eta^{-1}\xi=\xi^\times\eta,
\]
and all their conjugates and $\eta$-conjugates.

\indent (8) In the neutral form a finite rotation is effected by a vector linity (and its extensive) $\phi$ for which $\phi'\phi=1$. This is equivalent to
\[
\xi'\eta^{-1}\xi\eta^{-1}=\eta X'\eta^{-1}X=1.
\]
From the above it will be seen that if in a Galilean manifold we would deal with $\phi',\phi'\pm\phi$, etc., then probably in the general manifold it is best to deal with $\xi$ and $\xi^\times$; if in the Galilean we would use $\phi^2,f(\phi)$ etc., then in the general manifold the use of $X,X^\times$ is indicated.

\indent We need not write down the twenty-four alternative forms of the eight equations (6), but we may note how in terms of $N,I$ the other types appear:
\begin{equation}\left.\begin{aligned}
N,I&\text{ are of type }X,\\
N'=\eta^{-1}N\eta,\quad I'=\eta^{-1}I\eta&\text{ are of type }X^\times,\\
N\eta,\ I\eta&\text{ are of type }\xi,\\
\eta^{-1}N,\ \eta^{-1}I&\text{ are of type }\xi^\times
\end{aligned}\right\}.\tag{9}\end{equation}
After the fundamental linity $\eta$, of type $\xi$, has presented itself the other three may be regarded as arising from
\[
\xi^{-1},\qquad\xi_1\xi_2^{-1},\qquad\xi_1^{-1}\xi_2.
\]
\noindent\textsc{§ 14. \emph{Riemann's and Weyl's Manifolds.}}---Suppose $\rho,\rho+\alpha,\rho+\beta$ are for our manifold the position vectors of three infinitesimally near points, P, A, B. We assume that for such small regions the metrics are intrinsically the same as in a Galilean system, and in particular that there is a definite shift of PA along PB, and of PB along PA, each called a translation, by which if A in the first shift comes to C, B also in the second shift will come to C. With arbitrary co-ordinates we cannot of course affirm that C will have the position vector $\rho+\alpha+\beta$, but instead that it will have
\[
\rho+\alpha+\beta+d_\alpha^T\beta
\]
where $d_\alpha^T$ is to be regarded, not as a symbol of operation, but merely as the abbreviation of the words
\begin{equation}d_\alpha^T=\text{``increment due to translation along }\alpha\text{.''}\tag{1}\end{equation}
From these preliminaries it is clear that
\begin{equation}d_\alpha^T\beta=d_\beta^T\alpha=-E_a\beta,\tag{2}\end{equation}
where $E_a$ is a vector linity which is a function of the position vector $\rho$, and is also a function of $\alpha$, linear in the latter; also $E_a\beta$ is symmetrical in $\alpha$ and $\beta$.
The word translation implies something about constancy or otherwise in interval or in ``the square measure'' $ds^2$ of equation (2) of §11. In Weyl's view we are to assume that in general a \emph{translation} is such that it is impossible for the square measure to remain unaltered when by successive steps a closed path has been traced. In such a closed path beginning and ending at a given point, there will appear a change in the square measure depending on the point and the path.

\indent The square measure then is assumed to change according to the equation
\begin{equation}d_\alpha^T V_0\beta\eta\beta=-V_0\beta\eta\beta\cdot V_0\lambda\alpha,\tag{3}\end{equation}
where $\alpha$ and $\beta$ are arbitrary and with the same meanings as before, and $\lambda$ is a vector, a given function of $\rho$. Thus $\alpha,\beta$ must be contravariant, and $\lambda$ covariant. $E_a\beta$ is a vector neither covariant nor contravariant, which need not surprise us, because more than one position of the manifold has been used in its description.

\indent It is clear that $\eta$ is arbitrary as to a scalar factor and $\lambda$ correspondingly arbitrary. More definitely, nothing is altered in the meanings given if (1) $\eta$ is multiplied by a positive scalar $h$, any function of position, and (2) $\lambda$ is altered to $\lambda+\nabla\log h$, as appears from equation (3).


\indent \emph{Important qualification of the last sentence.}---$\eta$ here means the original \emph{vector} linity denoted by $\eta$. Whenever we speak of altering gauge by multiplying $\eta$ by $h$ this restriction is to be understood. It is clear that $\eta w$ thereby changes, not to $h\eta w$, but to $h^c\eta w$.

\indent ``Into all expressions or formulae, which exhibit analytically true metrical properties of the manifold, $\eta$ and $\lambda$ must enter in such a manner that invariance subsists (1) when there is arbitrary transformation of co-ordinates ($\rho$ replaced by $\rho'$); and (2) when $\eta,\lambda$ are replaced by $h\eta,\lambda+\nabla\log h$, whatever positive scalar function of $\rho,h$ may be.''\footnote{Weyl, \emph{Raum, Zeit, Materie}, 3rd ed., p. 110. English translation, 1922, of 4th ed. (Brose), last sentence of p. 123.}

Since (3) is true, and also what (3) becomes both when $\beta$ is changed to $\gamma$ and when changed to $\beta+\gamma$, we get the corresponding bilinear form
\begin{equation}d_\alpha^T V_0\beta\eta\gamma=-V_0\beta\eta\gamma\cdot V_0\lambda\alpha,\tag{4}\end{equation}
where $\alpha,\beta,\gamma$ are all arbitrary. In the translation denoted by $d_\alpha^T$, $\eta$ suffers the increment $-V_0\alpha\nabla\cdot\eta$ and $\beta$ and $\gamma$ the increments $-E_a\beta,-E_a\gamma$. Hence
\[
-V_0\alpha\nabla\eta\cdot V_0\beta\eta\gamma-V_0E_a\beta\eta\gamma-V_0\beta\eta E_a\gamma=-V_0\beta\eta\gamma\cdot V_0\lambda\alpha.
\]
Put
\begin{equation}\eta E_a=\Gamma_a,\tag{5}\end{equation}
so that $\Gamma_a\beta$, like $E_a\beta$, is symmetrical in $\alpha,\beta$. Thus
\[
-V_0(\gamma\Gamma_a\beta+\beta\Gamma_a\gamma)=V_0\alpha(\nabla\eta-\lambda)\cdot V_0\beta\eta\gamma.
\]
Write down the two equations obtained by cyclically changing $\alpha,\beta,\gamma$, change the signs of the second and third, and then add the three equations. We obtain
\[
2V_0\alpha\Gamma_\beta\gamma=V_0\alpha\bigl[(\nabla\eta-\lambda)V_0\beta\eta\gamma-V_0\beta(\nabla\eta-\lambda)\cdot\eta\gamma-V_0\gamma(\nabla\eta-\lambda)\cdot\eta\beta\bigr],
\]
or
\begin{equation}\Gamma_\beta\gamma=\tfrac12\bigl[-V_0\beta(\nabla\eta-\lambda)\cdot\eta\gamma-V_0\gamma(\nabla\eta-\lambda)\cdot\eta\beta+(\nabla\eta-\lambda)V_0\beta\eta\gamma\bigr].\tag{6}\end{equation}
For Riemann's manifold we have merely to put $\lambda=0$. Since frequent reference must be made to this case of Riemann's manifold, we will put for the values of $E_a$ and $\Gamma_a$ when $\lambda=0$, $\bar E_a$ and $\bar\Gamma_a$ respectively.

\indent From (6)
\begin{equation}(\Gamma_\beta+\Gamma_\beta')\gamma=-V_0\beta(\nabla\eta-\lambda)\cdot\eta\gamma,\tag{7}\end{equation}
\begin{equation}(\Gamma_\beta-\Gamma_\beta')\gamma=V_1\bigl[V_2(\nabla\eta-\lambda)\eta\beta\cdot\gamma\bigr].\tag{8}\end{equation}
Thus $\Gamma_\beta'\gamma$ is obtained from (6) by reversing the signs of the second and third terms on the right. It further follows that $\Gamma_\beta'\gamma$ regarded as a linity of $\beta$, instead of $\gamma$, is self-conjugate.

\indent This, and its converse, can be shown to be true of any such function $Y_{\alpha\beta}$ which is symmetrical in $\alpha$ and $\beta$. Thus we may put for the $\iota_c$ component of $Y_{\alpha\beta}$, $\iota_cV_0\alpha\psi_c\beta$ where $\psi_c$ is a self-conjugate vector linity, or
\begin{equation}\left.\begin{aligned}
Y_{\alpha\beta}&=\sum_{c=1}^n\iota_cV_0\beta\psi_c\alpha\\
Y_\alpha'\beta&=\sum_{c=1}^n\psi_c\alpha V_0\iota_c\beta
\end{aligned}\right\},\tag{9}\end{equation}
showing that $Y_\alpha'\beta$ is self-conjugate with respect to $\alpha$. Each $\psi_c$ involves $\tfrac12n(n+1)$ scalars, so that $Y_{\alpha\beta}$ involves $\tfrac12n^2(n+1)$ scalars; when $n=4$ this number is 40.

\indent When $Y_{\alpha\beta}$ is not symmetrical (9) still holds, $\psi_c$ now being not self-conjugate, and the number of scalars being $n^3$.

\indent We shall from this point onwards take for granted, as is not difficult to prove, that \emph{at a given point it is always possible so to choose our co-ordinates that all the $\tfrac12n^2(n+1)$ differential coefficients $(\nabla\eta,\eta)$ are zero; and independently to choose the gauge $(h)$ so that $\lambda$ is zero; therefore also to make $E_a$ zero for all values of $\alpha$.} [Brose calls $h$ the calibration ratio. Selection of a definite $h$ he calls calibration.]

\indent\textsc{§ 15. \emph{Absolute Differentiation.}}---Let $\tau$ be a contravariant vector given over an assigned path or over a region of $n$ or fewer dimensions. If the region is of fewer than $n$ dimensions, we may still assign arbitrary values over a region of $n$ dimensions containing the given region, so that we can continue to apply $\nabla$ to it.

\indent In passing from the point $\rho$ to the point $\rho+\alpha$, where $\alpha$ is infinitesimal, the
\emph{ordinary} increment $-V_0\alpha\nabla\cdot\tau$ will be denoted by $d_\alpha^0\tau$, and the excess of this above $d_\alpha^T\tau$ will be called the absolute increment and denoted by $d_\alpha^A\tau$. Thus we have
\begin{equation}\left.\begin{aligned}
d_\alpha^A\tau&=d_\alpha^0\tau-d_\alpha^T\tau=d_\alpha^0\tau+E_a\tau\\
&=-V_0\alpha\nabla\cdot\tau+E_a\tau
\end{aligned}\right\}.\tag{1}\end{equation}
We are about to show, as might be expected, that this absolute increment, $d_\alpha^A\tau$, is contravariant. If
\begin{equation}\chi=-V_0(\ )\nabla\cdot\rho',\tag{2}\end{equation}
and the change from $\rho$ to $\rho'$ causes the multenion function $q$ of position to change to $q'$, and the multenion linity $\phi$ to change to $\phi'$; then the necessary and sufficient conditions that $q$ should be (1) contravariant, (2) covariant, are
\[
(1)\quad q'=\chi q,\qquad (2)\quad q'=\chi'^{-1}q.
\]
The necessary and sufficient conditions that $\phi$ should be (1) covariant, (2) contravariant, (3) comixt, (4) contramixt, are that $\phi'$ should be equal to
\[
(1)\quad\chi'^{-1}\phi\chi^{-1},\qquad (2)\quad\chi\phi\chi',\qquad (3)\quad\chi'^{-1}\phi\chi',\qquad (4)\quad\chi\phi\chi^{-1}.
\]
All these may be put in infinitesimal forms which are generally easier to apply. Let
\[
\rho'=\rho+\delta\rho=\rho+\epsilon,\qquad \chi'=1+\delta\chi=1+\chi_0,
\]
\[
q'=q+\delta q,\qquad \phi'=\phi+\delta\phi.
\]
Then the above six conditions (necessary and sufficient) for the several types of invariance become
\[
(1)\quad\delta q=\chi_0q,\qquad (2)\quad\delta q=-\chi_0'q,
\]
and in the case of $\phi$, $\delta\phi$ is equal to
\[
(1)\ -\chi_0'\phi-\phi\chi_0,\qquad (2)\ \chi_0\phi+\phi\chi_0',\qquad (3)\ -\chi_0'\phi+\phi\chi_0',\qquad (4)\ \chi_0\phi-\phi\chi_0;
\]
and in order to make the tests we have
\begin{equation}\chi_0=-V_0(\ )\nabla\cdot\epsilon,\qquad \chi_0'=-\nabla_9V_0(\ )\epsilon_9.\tag{3}\end{equation}
[\,$\chi$ is applied to general multenions by the rules applicable to extensives; $\chi_0$ by those applicable to subextensives.\,]

\indent First apply the test to $E_a\tau$; is $\delta E_a\tau=\chi_0E_a\tau$? Instead of working this directly it will be found easier first to find $\delta V_0\alpha\Gamma_\beta\gamma$ from the equation preceding (6) of §14, in which observe that, since $V_0\alpha\nabla=V_0\alpha'\nabla'$, $\delta$ has no effect on $V_0\alpha\nabla,V_0\beta\nabla,V_0\gamma\nabla$. Further for $\alpha,\beta,\gamma$,
\[
\delta(\alpha,\beta,\gamma)=\chi_0(\alpha,\beta,\gamma),
\]
and for $\eta$
\[
\delta\eta=-\chi_0'\eta-\eta\chi_0.
\]
It is thus easy to prove that
\[
\delta V_0\alpha\Gamma_\beta\gamma=-V_0\beta\nabla_9\cdot V_0\gamma\nabla_9\cdot V_0\eta\alpha\epsilon_9,
\]
and thence that
\begin{equation}\delta E_\beta\gamma=\chi_0E_\beta\gamma-V_0\beta\nabla\cdot V_0\gamma\nabla\cdot\epsilon,\tag{4}\end{equation}
so that $E_a\tau$ is not contravariant. Next we have
\[
\begin{aligned}
\delta\cdot d_{\alpha}^{0}\tau-\chi_0d_{\alpha}^{0}\tau&=-\delta(V_0\alpha\nabla\cdot\tau)+\chi_0(V_0\alpha\nabla\cdot\tau)\\
&=-V_0\alpha\nabla\cdot\chi_0\tau+\chi_0(V_0\alpha\nabla\cdot\tau)
\end{aligned}
\]
because $\delta$ does not affect $V_0\alpha\nabla$, and $\delta\tau=\chi_0\tau$. Thus
\[
\delta\cdot d_{\alpha}^{0}\tau=\chi_0d_{\alpha}^{0}\tau+V_0\alpha\nabla_9\cdot V_0\tau\nabla_9\cdot\epsilon_9
\]
whence finally
\[
\delta d_{\alpha}^{A}\tau=\chi_0d_{\alpha}^{A}\tau,
\]
or $d_{\alpha}^{A}\tau$ is contravariant.

\indent In one respect our taking of equation (1) to define $d_{\alpha}^{A}\tau$ has been misleading. If we follow on from it and equation (4) §14 in a natural manner to define $d_{\alpha}^{A}\sigma^\times$ where $\sigma^\times$ is a covariant vector, the meaning arrived at will be on the whole inconvenient. For the moment we may take it as a sufficient reason for this statement that if equation (4) §14 is to hold for any scalar invariant function of position, $x$, then $d_{\alpha}^{T}x$ being different from zero, if $d_{\alpha}^{A}x$ is to be equal to $d_{\alpha}^{0}x-d_{\alpha}^{T}x$, absolute and ordinary increments of invariants must differ. We reject the suggestion of equation (1), then, except for the increment of a \emph{contravariant} vector. It is convenient to introduce a new increment, the invariantive increment, $d_{\alpha}^{I}q$. For $\tau$ we will take
\begin{equation}d_{\alpha}^{I}\tau=d_{\alpha}^{T}\tau=-E_a\tau,\tag{5}\end{equation}
but to fix the meaning for other cases we will take that
\begin{equation}d_{\alpha}^{I}x=0\tag{6}\end{equation}
when $x$ is an invariant scalar function of position. This gives
\[
0=d_{\alpha}^{I}V_0\sigma^\times\tau=V_0(d_{\alpha}^{I}\sigma^\times)\tau-V_0\sigma^\times E_a\tau.
\]
Since $\tau$ is arbitrary we get

\begin{equation}d_{\alpha}^{I}\sigma^\times=E_a'\sigma^\times.\tag{7}\end{equation}
Finally in all cases we put
\begin{equation}d_{\alpha}^{A}=d_{\alpha}^{0}-d_{\alpha}^{I}=-V_0\alpha\nabla\cdot-d_{\alpha}^{I}.\tag{8}\end{equation}
The strongest reason for this meaning of $d_{\alpha}^{A}$ remains to be given. It is the only increment of the kind which is \emph{truly} invariantive for \emph{truly} contravariant and covariant vectors; vectors which besides their usual relations to change of co-ordinates are also invariant for change of gauge. If $\tau$ is thus truly contravariant $\eta\tau$ cannot be truly covariant, and therefore we ought not to base our conception of the standard increment of $\sigma^\times$ on supposing it possible to put $\sigma^\times=\eta\tau$.

\indent Nevertheless the following up of the meaning of $d_{\alpha}^{A}\sigma^\times$ from equation (4)
of §14 has its uses, and we shall make one important application of it when later we consider curvature. It is not difficult to show that
\begin{equation}d_{\alpha}^{I}\sigma^\times=E_a'\sigma^\times-\sigma^\times V_0\alpha\lambda,\qquad d_{\alpha}^{I}\eta=d_{\alpha}^{0}\eta=-V_0\alpha\nabla\cdot\eta\tag{9}\end{equation}
(so that if we had continued to take $d_{\alpha}^{A}=d_{\alpha}^{0}-d_{\alpha}^{I}$ we should have found, as in the Riemann manifold, that $d_{\alpha}^{A}\eta=0).$ With the meaning of (8) above $d_{\alpha}^{A}\eta$ is not zero, as in a Riemann manifold, but instead
\[
d_{\alpha}^{A}\eta=-V_0\alpha\lambda\cdot\eta
\]
when $\eta$ is used in its primitive sense as a vector linity. When used as an extensive linity in $\eta w$ we get
\begin{equation}d_{\alpha}^{A}\eta\cdot w=-cV_0\alpha\lambda\cdot\eta w.\tag{10}\end{equation}
We will now find the absolute increments of the two kinds of multenions and the four of multenion linities. Mixed linities interpreted as subextensive, have invariant meanings, but not the other two types. Thus if $\phi$ is contramixt
\[
\phi V_3\alpha\beta\gamma=V_3\phi\alpha\beta\gamma+V_3\alpha\phi\beta\gamma+V_3\alpha\beta\phi\gamma
\]
is of invariant form, because then $\phi\alpha,\beta,\gamma$ are of one type, and is not of such form when $\phi$ is covariant, because then $\phi\alpha,\beta,\gamma$ are not all of one type.

\indent Let $E_a$ and $E_a'$ have the subextensive meanings as well as their primitive vector linity meanings. Thus
\begin{equation}E_aw=-V_cE_a\xi V_{c-1}\xi w,\qquad E_a'w=-V_cE_a'\xi V_{c-1}\xi w.\tag{11}\end{equation}
We at once have
\begin{equation}d_{\alpha}^{A}p=d_{\alpha}^{0}p-d_{\alpha}^{I}p=-V_0\alpha\nabla\cdot p+E_ap,\tag{12}\end{equation}
\begin{equation}d_{\alpha}^{A}q^\times=d_{\alpha}^{0}q^\times-d_{\alpha}^{I}q^\times=-V_0\alpha\nabla\cdot q^\times-E_a'q^\times.\tag{13}\end{equation}
For the four linities $\xi,\xi^\times,X,X^\times$, treat of the first as an example of all four. The meaning of $d_{\alpha}^{A}\xi$ is given by
\[
d_{\alpha}^{A}\xi\cdot p=d_{\alpha}^{A}(\xi p)-\xi d_{\alpha}^{A}p.
\]
Here $\xi p$ is a $q^\times$ so that
\[
d_{\alpha}^{A}(\xi p)=-V_0\alpha\nabla\cdot(\xi p)-E_a'\xi p.
\]
Hence
\[
d_{\alpha}^{A}\xi\cdot p=-V_0\alpha\nabla_\eta\cdot\xi\eta p-E_a'\xi p-\xi E_ap.
\]
Thus
\begin{equation}\left.\begin{aligned}
d_{\alpha}^{A}\xi&=-V_0\alpha\nabla\cdot\xi-E_a'\xi-\xi E_a\\
d_{\alpha}^{A}\xi^\times&=-V_0\alpha\nabla\cdot\xi^\times+E_a\xi^\times+\xi^\times E_a'\\
d_{\alpha}^{A}X&=-V_0\alpha\nabla\cdot X-E_a'X+XE_a'\\
d_{\alpha}^{A}X^\times&=-V_0\alpha\nabla\cdot X^\times+E_aX^\times-X^\times E_a
\end{aligned}\right\}.\tag{14}\end{equation}
It is instructive to re-write in the present notation the four kinds of $\delta\phi$ above
\begin{equation}\left.\begin{aligned}
\delta\xi&=-\chi_0'\xi-\xi\chi_0\\
\delta\xi^\times&=\chi_0\xi^\times+\xi^\times\chi_0'\\
\delta X&=-\chi_0'X+X\chi_0'\\
\delta X^\times&=\chi_0X^\times-X^\times\chi_0
\end{aligned}\right\}.\tag{15}\end{equation}
Doubtless the complete parallelism between (14) and (15) has some simple explanation. Of the four results in each set, the second follows from the first by putting $\xi^\times=\xi^{-1}$ and the fourth from the third by putting $X^\times=X'$.

\indent We obtain six very general formulae showing how to convert such expressions as $(\nabla_9,q_9^\times)$ or $(\nabla_9,p_9)$ to invariant form [provided $(\alpha^\times,q^\times)$ or $(\alpha^\times,p)$ are already bilinear and invariantive] by one common process. Introduce a second notation for $d_\alpha^A$, thus
\[
d_\alpha^Aq=q_\alpha,\qquad d_\alpha^A\phi=\phi_\alpha,\quad\text{etc.}
\]
Put $\zeta$ in place of $\alpha$ in (12), (13), and (14). Thus
\begin{equation}(\zeta,p_\zeta)=(\nabla_9,p_9)+(\zeta,E_\zeta p),\tag{12a}\end{equation}
\begin{equation}(\zeta,q_\zeta^\times)=(\nabla_9,q_9^\times)-(\zeta,E_\zeta'q^\times),\tag{13a}\end{equation}
\begin{equation}
\left.\begin{aligned}
(\zeta,\xi_\zeta)&=(\nabla_9,\xi_9)+(\zeta,-E_\zeta'\xi-\xi E_\zeta)\\
(\zeta,\xi_\zeta^\times)&=(\nabla_9,\xi_9^\times)+(\zeta,E_\zeta\xi^\times+\xi^\times E_\zeta')\\
(\zeta,X_\zeta)&=(\nabla_9,X_9)+(\zeta,-E_\zeta'X+XE_\zeta')\\
(\zeta,X_\zeta^\times)&=(\nabla_9,X_9^\times)+(\zeta,E_\zeta X^\times-X^\times E_\zeta)
\end{aligned}\right\}.\tag{14a}
\end{equation}
(12), (13), (14) are particular cases of (12a), (13a), (14a).

\indent As an example of the use of these equations let us suppose we have with Galilean co-ordinates established the equation $T_9\nabla_9=0$ where $T$ is the (complete) energy momentum linity; a linity which in the general manifold is best taken as comixt; a self-conjugate \emph{vector} linity (self-conjugacy means for a \emph{comixt} linity $T=\eta T'\eta^{-1}$, that is $\eta^{-1}T$ or $T\eta$ is self-conjugate in the ordinary sense). The (truly) invariant form of $T_9\nabla_9$ is given as $T_\zeta\zeta$ by the third of (14a). For generality we will at first not assume $T$ to be self-conjugate. We have
\[
T_\zeta\zeta=T_9\nabla_9-E_\zeta'T\zeta+TE_\zeta'\zeta.
\]
To evaluate $E_\zeta'T\zeta$, write down $E_\alpha'\beta^\times=\Gamma_\alpha'\eta^{-1}\beta^\times$ from (6) of §14, change $\alpha,\beta^\times$ to $\zeta,T\zeta$, and simplify. Thus
\begin{equation}E_\zeta'T\zeta=\tfrac12\eta_9(\eta^{-1}T-T'\eta^{-1})(\nabla_9-\lambda)-\tfrac12(\nabla_9-\lambda)V_0\zeta\eta_9\eta^{-1}T\zeta.\tag{16}\end{equation}
When $\eta^{-1}T$ is self-conjugate we get
\begin{equation}E_\zeta'T\zeta=-\tfrac12(\nabla_9-\lambda)V_0\zeta\eta_9\eta^{-1}T\zeta.\tag{17}\end{equation}
Further, since $|\phi|^{-1}\cdot d|\phi|=-V_0\zeta\phi^{-1}d\phi\zeta$, putting $T=1$ we get
\begin{equation}E_\zeta'\zeta=k^{-1}\nabla k-\tfrac12n\lambda.\tag{18}\end{equation}
In equation (16), the skew part of $\eta^{-1}T$ alone occurs in the first term, and the self-conjugate part alone in the second. $\eta^{-1}T$ is contravariant. Let us put $\eta^{-1}T=\xi^\times$ and use the second of (14a) for the further reduction
\[
\xi_\zeta^\times\zeta=\xi_9^\times\nabla_9+E_\zeta\xi^\times\zeta+\xi^\times E_\zeta'\zeta.
\]
We may use (18) to modify $E_\zeta'\zeta$. $E_\zeta\xi^\times\zeta$ happens to be particularly simple as it involves only the self-conjugate part of $\xi^\times$; let $\xi^\times=\xi_0^\times+\xi_1^\times$ where $\xi_0^\times$ is the self-conjugate part and $\xi_1^\times$ the skew part. We find that
\begin{equation}E_\zeta\xi^\times\zeta=\eta^{-1}\bigl(\eta_9\xi_0^\times+\tfrac12V_0\zeta\eta_9\xi_0^\times\zeta\bigr)(\nabla_9-\lambda),\tag{17a}\end{equation}
and
\begin{equation}\xi_\zeta^\times\zeta=k^{-1}(k_\zeta\xi^\times)_9(\nabla_9-\tfrac12n\lambda)+\eta^{-1}\bigl(\eta_9\xi_0^\times+\tfrac12V_0\zeta\eta_9\xi_0^\times\zeta\bigr)(\nabla_9-\lambda).\tag{17b}\end{equation}
Linities of the form $\phi+\tfrac12V_0\zeta\phi\zeta$ frequently present themselves in relativity, $\phi$ being mixed and generally self-conjugate in the mixed sense.

\indent Note that $\xi_\zeta\zeta$ and $X_\zeta^\times\zeta$ have no invariant meaning because $\xi_\zeta\beta^\times$ and $X_\zeta^\times\beta^\times$ have none. We thus do not seek generalisations of $\xi_9\nabla_9$ and $X_9^\times\nabla_9$. $\nabla$ is a symbolic \emph{covariant} vector because $V_0d\rho\nabla=V_0d\rho'\nabla'$.

\indent (17a) is found again below and numbered (22), in the course of illustrating a generally useful process.

\indent Returning to the generalisation required of $T_9\nabla_9=0$, when $T$ is comixt and $\eta^{-1}T$ is self-conjugate, we get
\begin{equation}(kT)_9(\nabla_9-\tfrac12n\lambda)+\tfrac12(\nabla_9-\lambda)V_0\zeta\eta_9\eta^{-1}(kT)\zeta=0.\tag{19}\end{equation}
It is clear that $kT$, a vector linity \emph{density} as it is called, rather than $T$ itself, should be regarded as the fundamental function. (19) is our standard form on the hypothesis that $T$ is \emph{truly} invariantive. If it were of power $a$ in the sense about to be explained, the first $\nabla_9$ in (19) should be replaced by $\nabla_9-a\lambda$.

\indent If an invariantive symbol has such a meaning that on change of gauge $(h\eta$ for $\eta,\ \lambda+h^{-1}\nabla h$ for $\lambda)$ the symbol is multiplied by $h^a$ then it is said to be of power $a$. $\eta$ is of power $1$ when used in its primitive sense as a vector linity; in $\eta w$, $\eta$ is of power $c$; $k$ is of power $\tfrac12n$. It is because of these powers that we find when $\nabla$ applies to $\eta$ (power $1$) as in $E_a$ and $\Gamma_a$, $\nabla$ appears only in the combination $\nabla-\lambda$; when $\nabla$ applies to $k$ (power $\tfrac12n$) we always find it in the combination $\nabla-\tfrac12n\lambda$. Thus in (18), although in a Riemann manifold we should find $E_\zeta'\zeta=k^{-1}\nabla k$, we must expect in the Weyl manifold to find instead
\[
E_\zeta'\zeta=k^{-1}(\nabla-\tfrac12n\lambda)k.
\]
The reason for this is easily seen. Thus if $q$ is of power $a$, $(\nabla_9,q_9)$ is of no
definite power at all, but as is quite easy to verify $(\nabla_9-a\lambda,q_9)$ is of power $a$.
To be \emph{truly} invariantive a symbol must be of power zero.

\indent An important case of this kind occurs in connection with the earliest invariant expressions we met, namely,
\[
k^{-1}V_{a-1}\nabla(ku),\qquad V_{b+1}\nabla v^\times.
\]
Our standard forms for these are based on the supposition that $u$ and $v^\times$ are each of power zero so that $V_{b+1}\nabla v^\times$ has not to be modified; but in the other case $ku$ is of power $\tfrac12n$ and our new invariant form should be $k^{-1}V_{a-1}(\nabla-\tfrac12n\lambda)(ku)$. Now (12a) and (13a) show us that our standard forms in these two cases should be
\[
V_{a-1}\nabla u+V_{a-1}\zeta E_\zeta u,\qquad V_{b+1}\nabla v^\times-V_{b+1}\zeta E_\zeta'v^\times.
\]
These must be the same as those we have just been considering so that we must have
\begin{equation}V_{a-1}\zeta E_\zeta u=k^{-1}V_{a-1}\bigl[(\nabla-\tfrac12n\lambda)k\mathbin{\cdot}u\bigr],\tag{20}\end{equation}

\begin{equation}V_{b+1}\zeta E_\zeta'v^\times=0.\tag{21}\end{equation}
These identities can be established directly, but not very easily.

The specially simple form of $\Gamma_a+\Gamma_a'$ given in (7) of \S14 furnishes a second useful result from (16) for it gives
\[
(\Gamma_\zeta'+\Gamma_\zeta)\eta^{-1}T\zeta=\eta_9\eta^{-1}T\zeta(\nabla_9-\lambda)
\]
where $\Gamma_\zeta'\eta^{-1}=E_\zeta'$. Thus from (16)
\begin{equation}\Gamma_\zeta\eta^{-1}T\zeta=-\tfrac12\eta_9(\eta^{-1}T+T'\eta^{-1})(\nabla_9-\lambda)+\tfrac12(\nabla_9-\lambda)V_0\zeta\eta_9\eta^{-1}T\zeta.\tag{22}\end{equation}
First we have the specially simple case $T=1$, namely
\begin{equation}\Gamma_\zeta\eta^{-1}\zeta=\eta_9\eta^{-1}\nabla_9-k^{-1}\nabla k+\tfrac12(n-2)\lambda.\tag{23}\end{equation}
Another very simple particular case is when $\eta^{-1}T$ is skew. The right of (22) is then zero, and $\eta^{-1}T$ is of the form $V_{1\omega}(\ )$ where $\omega$ is a contravariant ${}_nV_2$. Thus
\[
\Gamma_\zeta V_{1\omega}\zeta=0.
\]
Since $\Gamma_a=\eta E_a$ we also have $E_\zeta V_{1\omega}\zeta=0$. This last may be regarded as a special case of (21) and suggests two other forms of (20), (21). Operate on (20) by $V_{0\omega^\times}(\ )$ taking $c=a-1$; and on (21) by $V_{0\omega}(\ )$ taking $c=b+1$. We thus get
\begin{equation}E_\zeta'V_{c+1}\omega^\times\zeta=k^{-1}V_{c+1}\omega^\times(\nabla-\tfrac12n\lambda)k,\tag{24}\end{equation}
\begin{equation}E_\zeta V_{c-1}\omega\zeta=0,\tag{25}\end{equation}
(18) is a particular case of (24).

\indent \S16. \emph{Curvature.}---If $(\alpha,\beta)$ is any function of two vectors, linear in each then
\begin{equation}(\alpha,\beta)-(\beta,\alpha)=(\zeta,V_1[V_2\alpha\beta\mathbin{\cdot}\zeta]).\tag{1}\end{equation}
as is seen by putting $(\alpha,\beta)=-(\zeta,\beta V_0\alpha\zeta)$ and $(\beta,\alpha)=-(\zeta,\alpha V_0\beta\zeta)$. Thus when
\begin{equation}(\alpha,\beta)=-(\beta,\alpha)\quad\text{then}\quad(\alpha,\beta)=\tfrac12(\zeta,V_1[V_2\alpha\beta\mathbin{\cdot}\zeta]).\tag{2}\end{equation}
The application of these two equations should be constantly borne in mind in what follows, even when attention is not called to it.

\indent Similarly for the various invariant forms of infinitesimal rotations of a vector,
\[
V_1\omega\alpha^\times,\qquad V_1\omega^\times\alpha,\qquad V_1\omega\eta\alpha,\qquad V_1\omega^\times\eta^{-1}\alpha^\times.
\]
Let us calculate the increase in $\gamma$, a contravariant vector, due to translation, starting from the point $p$, along the four vector sides of an infinitesimal parallelogram in the order $\alpha,\beta,-\alpha,-\beta$, $\alpha$ and $\beta$ being contravariant. Afterwards let us calculate the similar increments of $\eta\gamma$. We shall put $V_2\alpha\beta=\omega$, and the increments due to translation round the parallelogram will be denoted by $d_\omega^T\gamma$, $d_\omega^T\eta\gamma$. Since the difference of two contravariant vectors at a point $p$ is itself contravariant it follows that $d_\omega^T\gamma$ is contravariant, and similarly $d_\omega^T\eta\gamma$ is covariant; and it is obvious that they must both be linear in both $\omega$ and $\gamma$. We shall, therefore, put
\begin{equation}d_\omega^T\eta\gamma=\eta d_\omega^T\gamma=\eta C_\omega\gamma=\Omega_\omega\gamma,\tag{3}\end{equation}
$C_\omega$ being a contramixt vector linity and $\Omega_\omega$ the corresponding covariant linity, which may be called the curvature linities. Along the side $+\alpha$ the increment in $\gamma$ is $-E_\alpha\gamma$. Along the side $-\alpha$, $\gamma$ has changed to $\gamma-E_\beta\gamma$ and $d_{-\alpha}^T=-d_\alpha^T$ has received the increment $-V_0\beta\nabla\mathbin{\cdot}(\ )$. Thus, these two sides contribute increments
\[
-E_\alpha\gamma+(E_\alpha-V_0\beta\nabla\mathbin{\cdot}E_\alpha)(\gamma-E_\beta\gamma),
\]
\[
=-V_0\beta\nabla\mathbin{\cdot}E_\alpha\gamma-E_\alpha E_\beta\gamma.
\]
The contributions of the other two sides are obtained by interchange of $\alpha$ and $\beta$ and reversal of the sign. Hence, noting that $E_\alpha$ is linear in $\alpha$,
\begin{equation}C_\omega\gamma=d_\omega^T\gamma=-E_\zeta\mathbin{\cdot}V_{1\omega}\nabla\gamma+(-E_\alpha E_\beta+E_\beta E_\alpha)\gamma.\tag{4}\end{equation}
\indent [Since $d_\alpha^A=d_\alpha^0-d_\alpha^T$, and since $d_\alpha^0$ taken round the parallelogram produces zero we get that $C_\omega\tau=\tau_{\alpha\beta}-\tau_{\beta\alpha}$ which shows that absolute differentiations in different directions are not commutative.]

\indent In treating of $\eta\gamma$ in place of $\gamma$, instead of $-E_\alpha$ we have to write $E_\alpha'-V_0\lambda\alpha$ for $d_\alpha^T$. Noting that if $x,y$ are two scalars $(-V_0\lambda\alpha,-V_0\lambda\beta)$
\[
-(E_\alpha'+x)(E_\beta'+y)+(E_\beta'+y)(E_\alpha'+x)=-E_\alpha'E_\beta'+E_\beta'E_\alpha'
\]
we see at once how (4) has to be modified. Thus
\begin{equation}\Omega_\omega\gamma=d_\omega^T\eta\gamma=E_\zeta\mathbin{\cdot}V_{1\omega}\nabla\eta\gamma+(-E_\alpha'E_\beta'+E_\beta'E_\alpha')\eta\gamma-\eta\gamma\mathbin{\cdot}V_0\omega\omega^\times,\tag{5}\end{equation}
where
\begin{equation}\omega^\times=V_2\nabla\lambda.\tag{6}\end{equation}
[$\omega$ and $\omega^\times$ are of very different natures. $\omega$ is a contravariant \emph{dummy}, $\omega^\times$ is a given truly covariant function of position; in relativity $\omega^\times$ is magnetic induction cum electrostatic intensity, $\lambda$ being vector potential cum electrostatic potential.]

\indent From (4) on putting $\omega=V_2\alpha\beta$, and noting that $E_\alpha\gamma$ is symmetrical in $\alpha$ and $\gamma$, the following ``cyclic rule'' is seen to be true for each of the two parts of $C_\omega\gamma$ separately
\begin{equation}C_{V_2\alpha\beta}\gamma+C_{V_2\beta\gamma}\alpha+C_{V_2\gamma\alpha}\beta=0.\tag{7}\end{equation}
Remembering that $\Omega_\omega=\eta C_\omega$, inspection of (4) and (5) shows that
\begin{equation}
\left.\begin{aligned}
\Omega_\omega+\Omega_\omega'&=-V_0\omega\omega^\times\mathbin{\cdot}\eta\\
\text{or}\qquad C_\omega+\eta^{-1}C_\omega'\eta&=-V_0\omega\omega^\times
\end{aligned}\right\}.\tag{8}
\end{equation}
The equation in $\Omega_\omega+\Omega_\omega'$ is even simpler than that in $\Gamma_a+\Gamma_a'$; and for a Riemann manifold $\Omega_\omega+\Omega_\omega'=0$.

\indent The first of (8) shows that $\Omega_\omega=-\tfrac12V_0\omega\omega^\times\mathbin{\cdot}\eta+V_1\varpi'(\ )$ where $\varpi'$ is a covariant linear function of $\omega$. Hence
\begin{equation}\Omega_\omega\gamma=-\tfrac12V_0\omega\omega^\times\mathbin{\cdot}\eta\gamma+V_1\varpi\omega\gamma,\tag{9}\end{equation}
where $\varpi$ is a covariant ${}_nV_2$ linity. The first term on the right of (9) gives the ratio of stretch $-\tfrac12V_0\omega\omega^\times$; the second term the rotation $\varpi\omega$, suffered by $\gamma$ by $d_\omega^T$.

\indent To find the properties of the linity $\varpi$ first transform (5) to the following,
\begin{equation}\Omega_\omega\gamma=\Gamma_\zeta\mathbin{\cdot}V_{1\omega}(\nabla_9-\lambda)'\gamma+(\Gamma_\alpha'\eta^{-1}\Gamma_\beta-\Gamma_\beta'\eta^{-1}\Gamma_\alpha)\gamma-\eta\gamma\mathbin{\cdot}V_0\omega\omega^\times.\tag{10}\end{equation}
I will indicate in outline the mode of doing this. The first term on the right of (5) is
\[
(\Gamma_\alpha'\eta^{-1})_9\eta\gamma V_0\beta\nabla_9-(\Gamma_\beta'\eta^{-1})_9\eta\gamma V_0\alpha\nabla_9
\]
and the second term is
\[
(-\Gamma_\alpha'\eta^{-1}\Gamma_\beta'+\Gamma_\beta'\eta^{-1}\Gamma_\alpha')\gamma=(\Gamma_\alpha'\eta^{-1}\Gamma_\beta-\Gamma_\beta'\eta^{-1}\Gamma_\alpha)\gamma+\text{two terms}
\]
of which the first is $-\Gamma_\alpha'\eta^{-1}(\Gamma_\beta+\Gamma_\beta')\gamma$ and the second is $+\Gamma_\beta'\eta^{-1}(\Gamma_\alpha+\Gamma_\alpha')\gamma$. It will now be found that the parts depending on $\nabla$ in $(\Gamma_\beta+\Gamma_\beta')$ and $(\Gamma_\alpha+\Gamma_\alpha')$ cancel with the following parts of the first term.
\[
\Gamma_\alpha'(\eta^{-1})_9\eta\gamma\mathbin{\cdot}V_0\beta\nabla_9-\Gamma_\beta'(\eta^{-1})_9\eta\gamma\mathbin{\cdot}V_0\alpha\nabla_9.
\]
Collecting the parts not cancelled we get (10).

\indent In the first term and the third on the right of (10) there are terms of the form $(\nabla_9,\lambda_9)$. These give
\[
-\tfrac12\eta\gamma\mathbin{\cdot}V_0\omega\omega^\times-\tfrac12V_1\mathbin{\cdot}V_2(\eta V_1\omega\nabla_9\mathbin{\cdot}\lambda_9)\gamma
\]
\[
=-\tfrac12\eta\gamma\mathbin{\cdot}V_0\omega\omega^\times-\tfrac14V_1\mathbin{\cdot}V_2(\eta V_1\omega\nabla_9\mathbin{\cdot}\lambda_9+\eta V_1\omega\lambda_9\mathbin{\cdot}\nabla_9)\gamma
\]
\[
-\tfrac14V_1\mathbin{\cdot}V_2(\eta V_1\omega\nabla_9\mathbin{\cdot}\lambda_9-\eta V_1\omega\lambda_9\mathbin{\cdot}\nabla_9)\gamma.
\]
The last of these three parts furnishes the skew part $\varpi_1$, of $\varpi$, given by
\begin{equation}\varpi_1\omega=\tfrac14V_2(V_1\zeta\omega^\times\mathbin{\cdot}\eta V_1\zeta\omega).\tag{11}\end{equation}
The neutral form of this is the suggestive expression $\tfrac14V_2\omega^\times\omega$. The remainder of $\varpi$, denoted by $\varpi_0$, the self-conjugate part, is
\begin{equation}
\left.\begin{aligned}
\varpi_0\omega={}&-\tfrac12V_2\bigl[(\nabla_9-\lambda)\eta_9V_1\omega(\nabla_9-\lambda)\bigr]+\tfrac12V_2\zeta_0\Gamma_\zeta'\eta^{-1}\Gamma_\zeta V_1\omega\zeta\\
&-\tfrac14V_2\bigl(\eta V_1\omega\nabla_9\mathbin{\cdot}\lambda_9+\eta V_1\omega\lambda_9\mathbin{\cdot}\nabla_9\bigr)
\end{aligned}\right\}.\tag{12}
\end{equation}
So far we have divided $\Omega_\omega\gamma$ into three parts, such that the corresponding parts of $C_\omega\gamma$ are \emph{truly} invariant, (1) the part self-conjugate in $\gamma$, namely, $-\tfrac12\eta\gamma V_0\omega\omega^\times$; (2) the part $V_1\varpi_1\omega\gamma$ in which $\varpi_1$ is skew and (apart from the fundamental Unity $\eta$) only involves $\omega^\times$ as with the first part; (3) the part $V_1\varpi_0\omega\gamma$ where $\varpi_0$ is self-conjugate. Part (3) obeys the cyclic rule; (1) and (2) together, but not separately, obey the same. The cyclic rule for $\varpi_0$ may be put, as will be shown later, in the form
\[
V_4\mathbin{\cdot}V_2\zeta_1\zeta_2\varpi_0V_2\zeta_1\zeta_2=0,
\]
so that $\varpi_0$ instead of having the full number
\[
\tfrac12\mathbin{\cdot}\tfrac12n(n-1)\bigl[\tfrac12n(n-1)+1\bigr]
\]
of scalars appropriate to a self-conjugate ${}_nV_2$ linity, has a smaller number because of the $n(n-1)(n-2)(n-3)/4!$ scalar equations of the cyclic condition just written. The number of scalars of $\varpi_0$ is thus $\tfrac1{12}n^2(n^2-1)$. For the parts (1) and (2) the $\tfrac12n(n-1)$ scalars of $\omega^\times$ are required. The total number of scalars involved in $C_\omega$ is thus
\[
\tfrac1{12}n(n-1)(n^2+n+6).
\]
When $n=4$ the number is 26.


\indent Putting $V_1\zeta\omega=\zeta_\omega$, $V_1\zeta\omega^\times=\zeta_{\omega\times}$ the two parts (1) and (2) together may be written as
\begin{equation}
\left.\begin{aligned}
-\tfrac12\eta\gamma V_0\zeta_\omega\zeta_{\omega\times}+\tfrac14V_1V_2(\zeta_{\omega\times}\eta\zeta_\omega)\gamma\\
=\tfrac14\bigl(-\eta\gamma\mathbin{\cdot}V_0\zeta_\omega\zeta_{\omega\times}-\eta\zeta_\omega\mathbin{\cdot}V_0\zeta_{\omega\times}\gamma+\zeta_{\omega\times}V_0\zeta_\omega\gamma\bigr)
\end{aligned}\right\}.\tag{13}
\end{equation}
It is now desirable to separate $\varpi_0$ into three parts $\varpi_a,\varpi_b,\varpi_c$, which are covariant as to change of co-ordinates though not as to change of gauge, namely, the parts of degrees zero, one, two in $\lambda$.

\indent $\varpi_a$, of course, gives precisely the Riemann terms, obtained by putting $\lambda=0$ in any expression for $\varpi$. Thus from (12)
\begin{equation}\varpi_a\omega=-\tfrac12V_2\mathbin{\cdot}\nabla_9\eta_9V_1\omega\nabla_9+\tfrac12V_2\zeta_0\Gamma_\zeta'\eta^{-1}\Gamma_\zeta V_1\omega\zeta.\tag{14}\end{equation}
Next write down the terms involving differentiations of $\lambda$ and make them invariantive by (13a) of \S15, using, however, $\bar E_\zeta'$ in place of $E_\zeta'$ as we desire invariance only for change of co-ordinates, and we desire \emph{not} to introduce quadratic terms in $\lambda$. Thus we obtain for the probable value of $\varpi_b$
\begin{equation}
\left.\begin{aligned}
\varpi_b\omega={}&-\tfrac14V_2\bigl(\eta V_1\omega\nabla_9\mathbin{\cdot}\lambda_9+\eta V_1\omega\lambda_9\mathbin{\cdot}\nabla_9\bigr)\\
&+\tfrac14V_2\bigl(\eta V_1\omega\zeta\mathbin{\cdot}\bar E_\zeta'\lambda+\eta V_1\omega\bar E_\zeta'\lambda\mathbin{\cdot}\zeta\bigr)
\end{aligned}\right\}.\tag{15}
\end{equation}
Before verifying that this contains all the terms linear in $\lambda$, not differentiated, find the quadratic terms, $\varpi_c\omega$. Put $\Lambda_a$ for the terms of $\Gamma_a$ which contain $\lambda$, that is, $\Lambda_a=\Gamma_a-\bar\Gamma_a$. Thus
\begin{equation}\varpi_c\omega=-\tfrac12V_2\lambda\eta V_1\omega\lambda+\tfrac12V_2\zeta_0\Lambda_\zeta'\eta^{-1}\Lambda_\zeta V_1\omega\zeta.\tag{16}\end{equation}
To simplify this put $\omega=V_2\alpha\beta$ and work in the neutral form, finally returning to invariant form. The result comes out that $4\varpi_c\omega$ is the normal expression of $\omega$ with reference to $\lambda$ (in neutral form $\lambda V_3\lambda\omega$). Hence
\begin{equation}\varpi_c\omega=\tfrac14\eta V_2\lambda V_3\eta^{-1}\lambda\omega=\tfrac14V_2\eta^{-1}\lambda V_3\lambda\eta\omega.\tag{17}\end{equation}
$\varpi_a,\varpi_b,\varpi_c$ are all separately self-conjugate, and all separately obey the cyclic rule.

\indent We will now show that (15) must contain all the terms of (12) which are linear in $\lambda$. Inspection of (12) shows that these terms form a function $(\nabla_9,\eta_9,\lambda,\omega)$ which is linear in all four constituents. If there is any term of this nature which has not been included in our three parts $\varpi_a,\varpi_b,\varpi_c$, such term must be covariant as to change of co-ordinates. Now $(\nabla_9,\eta_9,\lambda,\omega)$ can only be made covariant in this sense by the addition [first of (14a) of \S15] to it of
\[
(\zeta,-\bar E_\zeta'\eta-\eta\bar E_\zeta,\lambda,\omega)=-(\zeta,\bar\Gamma_\zeta'+\bar\Gamma_\zeta,\lambda,\omega)=-(\nabla_9,\eta_9,\lambda,\omega).
\]
This addition reduces the supposed absent term to zero. Hence there is no such term.

\indent We now proceed to the two successive ``contractions'' of $\Omega_\omega$ or $C_\omega$. The first contraction yields a vector linity denoted by $\Omega$ or $C$ in its covariant or contramixt form. The second contraction yields a scalar, an invariant of course, $F$.

\indent $\Omega_{V_2\zeta\beta}\eta^{-1}\zeta$ is a covariant vector because the subscript $\zeta$ occupies the place of a contravariant vector and the second $\zeta$ occupies the place of a covariant vector. Put then
\begin{equation}\Omega\beta=\Omega_{V_2\zeta\beta}\eta^{-1}\zeta,\qquad C\beta=\eta^{-1}\Omega\beta=C_{V_2\zeta\beta}\eta^{-1}\zeta.\tag{18}\end{equation}
\begin{equation}F=-V_0\zeta\eta^{-1}\Omega\zeta=-V_0\zeta C\zeta.\tag{19}\end{equation}
\indent Note, that of the four symbols $\Omega_\omega,C_\omega,\Omega,C$, \emph{the only two which are truly invariantive are $C_\omega$ and $\Omega$}; $\Omega_\omega$ is of power $1$ and $C$ of power $-1$; $F$ also is of power $-1$.

\indent Let us now obtain the terms contributed to $\Omega\beta$ and to $F$ by our five parts of $\Omega_\omega$, that is, by the first term in (9) and the four parts of $\varpi$, namely, $\varpi_1,\varpi_a,\varpi_b,\varpi_c$. For the latter four other than $\varpi_a$ work details in the neutral form, and at the last step put the result into invariant form by the rules of \S12. Where first differential coefficients of $\lambda$ occur, attend in the detailed working only to the terms in which the differential coefficients occur, and in
the last step use equation (13a), of \S15, modified by reading $\bar E_\zeta'$ instead of $E_\zeta'$ to obtain the invariant form.

\indent Let us denote the parts contributed by the first term of equation (9), to $\Omega$ and $F$ by $\Omega_2$ and $F_2$, and those contributed by $\varpi_1,\varpi_a,\varpi_b,\varpi_c$, by $\Omega_1,\Omega_a,\Omega_b,\Omega_c$ and by $F_1,F_a,F_b,F_c$, respectively. Thus we find
\begin{equation}
\left.\begin{aligned}
\Omega_2\beta&=-\tfrac12V_1\omega^\times\beta,\qquad \Omega_1\beta=\tfrac14(n-2)V_1\omega^\times\beta\\
(\Omega_2+\Omega_1)\beta&=\tfrac14(n-4)V_1\omega^\times\beta
\end{aligned}\right\}.\tag{20}
\end{equation}
These are the only parts of $\Omega\beta$ that are not self-conjugate, and they are purely skew. It is a remarkable fact that $\Omega$ is self-conjugate for the case $n=4$, and for no other value of $n$. Of course, this accidental circumstance, if we are to call it such, is fortunate for the theory of relativity. From (20), it at once follows that
\begin{equation}F_2=F_1=0.\tag{21}\end{equation}
It is easy to show that when a term $\varpi_x$ of $\varpi$ is self-conjugate the corresponding part $\Omega_x$ of $\Omega$ is self-conjugate. $\varpi_a,\varpi_b,\varpi_c$ are self-conjugate. We find
\begin{equation}
\left.\begin{aligned}
4\Omega_b\beta={}&-2\eta\beta\mathbin{\cdot}k^{-1}V_0\nabla(k\eta^{-1}\lambda)-(n-2)\bigl(\nabla_9V_0\lambda_9\beta+\lambda_9V_0\nabla_9\beta+2\bar E_\beta'\lambda\bigr)\\
F_b={}&-(n-1)k^{-1}V_0\nabla(k\eta^{-1}\lambda)
\end{aligned}\right\}.\tag{22}
\end{equation}
In $\Omega_c$ the normal expression of $\beta$ with respect to $\lambda$ occurs thus:
\begin{equation}
\left.\begin{aligned}
\Omega_c\beta&=\tfrac14(n-2)V_1\eta^{-1}\lambda V_2\lambda\eta\beta\\
F_c&=\tfrac14(n-1)(n-2)V_0\lambda\eta^{-1}\lambda
\end{aligned}\right\}.\tag{23}
\end{equation}
In a Riemann manifold $\varpi_a$ completely specifies $\Omega_\omega$, and therefore $\Omega,\Omega_{\omega a}$, and $\Omega_a$ are given by any expressions for $\Omega_\omega$ and $\Omega$ by putting $\lambda=0$. The most useful form of $\Omega_a$ is derived from (5). To get $\Omega\beta$ from (5) we have to put $\alpha=\zeta$, $\eta\gamma=\zeta$. Thus we obtain
\begin{equation}\Omega\beta=\bar E_{\beta9}'\nabla_9-\bar E_\zeta'\bar E_\beta'\zeta+\bar E_\beta'\bar E_\zeta'\zeta+V_0\beta\nabla\mathbin{\cdot}\bar E_\zeta'\zeta+V_1\beta\omega^\times.\tag{24}\end{equation}
Here note that the first three terms are just those, when we put $\bar E_\beta'=T$, that were under discussion in equations (16), (17), (18) of \S15, and (24) may be modified accordingly. Our immediate purpose, however, is to put $\lambda$, and therefore $\omega^\times$ zero, so as to obtain $\Omega_a$. Noting that $\bar E_\zeta'\zeta=k^{-1}\nabla k$ by (18) of \S15, we have
\begin{equation}\Omega_a\beta=\bar E_{\beta9}'\nabla_9-\bar E_\zeta'\bar E_\beta'\zeta+(\bar E_\beta'+V_0\beta\nabla\mathbin{\cdot})\nabla\log k.\tag{25}\end{equation}
This is the form (26$'$3) given in Eddington's Report on Relativity, that is to say, Eddington's $G_{bc}$ is our $V_0\iota_b\Omega_a\iota_c$ [see equation (7), \S17 below]. (25) gives
\begin{equation}F_a=-V_0\eta^{-1}\zeta\bigl(\bar E_{\zeta9}'\nabla_9-\bar E_\zeta'\bar E_\zeta'\zeta+[\bar E_\zeta'+V_0\zeta\nabla\mathbin{\cdot}]\nabla\log k\bigr).\tag{26}\end{equation}
\indent Returning to equation (24) let us find a useful expression for $C_a\beta$. It is convenient on occasion to use the following notations
\[
\eta^{-1}\Gamma_\beta\eta^{-1}=\Gamma_\beta^\times,\qquad \eta^{-1}\bar\Gamma_\beta\eta^{-1}=\bar\Gamma_\beta^\times,
\]
but this is a little cumbrous for our present purpose. Let us put
\begin{equation}\bar\Gamma_\beta^\times=\eta^{-1}\bar E_\beta'=\xi_\beta^\times=\mu_\beta^\times+\nu_\beta^\times,\tag{24a}\end{equation}
where $\mu_\beta^\times$ and $\nu_\beta^\times$ [see (24c) below] are put for the self-conjugate and skew parts of $\xi_\beta^\times$, and avoid the double suffix notation $\xi_{\beta0}^\times,\xi_{\beta1}^\times$. Putting $\lambda=0$ in (24) we have
\[
C_a\beta=\eta^{-1}(\eta\xi_\beta^\times)_9\nabla_9-\eta^{-1}\bar\Gamma_\zeta'\xi_\beta^\times\zeta+\xi_\beta^\times\nabla\log k+\eta^{-1}(V_0\beta\nabla\mathbin{\cdot}\nabla\log k)
\]
\[
\text{[Put }-\bar\Gamma_\zeta'=\bar\Gamma_\zeta-(\bar\Gamma_\zeta+\bar\Gamma_\zeta')=\bar\Gamma_\zeta+V_0\zeta\nabla\mathbin{\cdot}\eta\text{.]}
\]
\[
=\xi_{\beta9}^\times\nabla_9+\bar E_\zeta'\xi_\beta^\times\zeta+k^{-1}\xi_\beta^\times\nabla k+\eta^{-1}(V_0\beta\nabla\mathbin{\cdot}\nabla\log k)
\]
whence, using (17a) of \S15 for $\bar E_\zeta'\xi_\beta^\times\zeta$, which involves $\mu_\beta^\times$ but not $\nu_\beta^\times$,
\begin{equation}C_a\beta=k^{-1}(k\xi_\beta^\times)_9\nabla_9+\eta^{-1}\bigl(\eta\mu_\beta^\times+\tfrac12V_0\zeta\eta\mu_\beta^\times\zeta\bigr)\nabla+\eta^{-1}(V_0\beta\nabla\mathbin{\cdot}\nabla\log k).\tag{24b}\end{equation}
$\mu_\beta^\times$ and $\nu_\beta^\times$ being $\tfrac12\eta^{-1}(\pm\bar\Gamma_\beta+\bar\Gamma_\beta')\eta^{-1}$ are written down from (7) and (8) of \S14. Thus
\begin{equation}
\left.\begin{aligned}
\mu_\beta^\times\gamma^\times&=\tfrac12V_0\beta\nabla\zeta\mathbin{\cdot}(\eta^{-1})_9\gamma^\times\\
\nu_\beta^\times\gamma^\times&=-\tfrac12V_1\bigl[\eta^{-1}V_2\nabla\eta\beta\mathbin{\cdot}\gamma^\times\bigr]
\end{aligned}\right\}.\tag{24c}
\end{equation}
From (24b) we may write down $kF_a$ in the form $-kV_0\zeta C_a\zeta$ and we shall reach an expression useful for reducing, later, $\delta(kF_a)$ resulting from changing $\eta$ to $\eta+\delta\eta$.

\indent In what immediately follows let us work with the neutral forms, only occasionally stopping to point out the slightly more complicated invariant forms. In the case of a vector linity, $\phi$, we may say that there is a definite part, namely, $-n^{-1}V_0\zeta\phi\zeta$, which gives rise to the contraction $-V_0\zeta\phi\zeta$; the second part, $\phi+n^{-1}V_0\zeta\phi\zeta$, has zero contraction and involves only $n^2-1$ scalars as against the full number, $n^2$, for a vector linity. Similarly, it is possible to separate out from any vector expression $(\omega,\gamma)$ linear in each constituent a definite part $(n-1)^{-1}(V_2\zeta V_1\omega\gamma,\zeta)$ absorbing (generally) $n^2$ scalars from the full number, $\tfrac12n^3(n-1)$ required to specify $(\omega,\gamma)$. The remainder of $(\omega,\gamma)$, namely
\[
(\omega,\gamma)-(n-1)^{-1}(V_2\zeta V_1\omega\gamma,\zeta)
\]
has zero contraction. The part thus furnishing the contraction $C$ of $C_\omega$ is
\[
(n-1)^{-1}C_{V_2\zeta V_1\omega\gamma}\zeta=(n-1)^{-1}CV_1\omega\gamma
\]
and it obviously satisfies the cyclic rule (and would do so whether $C_\omega\gamma$ itself did so or not).

\indent The following form of the cyclic rule for $C_\omega\gamma$ may be noticed:---
\begin{equation}C_{V_2\zeta V_3\omega\gamma}\zeta=0.\tag{27}\end{equation}
This is true because by (1) of \S7
\[
V_2\zeta V_3\omega\gamma=V_2\beta\gamma V_0\alpha\zeta+V_2\gamma\alpha V_0\beta\zeta+V_2\alpha\beta V_0\gamma\zeta.
\]
(27) is closely connected with a generalisation of (1) and (2) above. The cyclic rule is an identity which is combinatorial in $\alpha,\beta,\gamma$; that is to say, the sign of the expression equated to zero alters by interchange of any two of the three $\alpha,\beta,\gamma$. If $(\alpha_1\alpha_2,\ldots,\alpha_a)$ is any expression linear in each vector and combinatorial in the vectors, it can always be expressed as a linear function of $\alpha_a^{(a)}$. It is actually equal to
\begin{equation}(a!)^{-1}(\zeta_1,\zeta_2,\ldots,\zeta_{a-1},V_1\alpha_a^{(a)}\zeta_{a-1}^{(a-1)}).\tag{28}\end{equation}
(28) may be proved by mathematical induction by proving the general connecting link in the following successive equalities, $a$ being, for example, taken as equal to 4,
\[
(\alpha,\beta,\gamma,\delta)=(2!)^{-1}(\zeta,V_1\mathbin{\cdot}V_2\alpha\beta\mathbin{\cdot}\zeta,\gamma,\delta)
\]
\[
=(3!)^{-1}(\zeta_1,\zeta_2,V_1\mathbin{\cdot}V_3\alpha\beta\gamma V_2\zeta_1\zeta_2,\delta)
\]
\[
=(4!)^{-1}(\zeta_1,\zeta_2,\zeta_3,V_1\mathbin{\cdot}V_4\alpha\beta\gamma\delta V_3\zeta_1\zeta_2\zeta_3).
\]
This is readily done when it is noticed that
\[
V_1\zeta_3^{'''}V_4\alpha\beta\gamma\delta=\sum\pm\alpha V_0\zeta_3^{'''}\beta\gamma\delta
\]
and the like $(\zeta_3^{'''}=V_3\zeta_1\zeta_2\zeta_3)$.

\indent Since $V_2\zeta V_1\omega\gamma-V_2(V_1\zeta\omega)\gamma=-\omega V_0\zeta\gamma$ we have
\begin{equation}C_\omega\gamma=CV_1\omega\gamma-C_{V_2(V_1\zeta\omega)\gamma}\zeta.\tag{29}\end{equation}
whether or not $C_\omega\gamma$ satisfies the cyclic rule. If the rule is satisfied we have further
\begin{equation}C_\omega\gamma=C_{V_2(V_1\zeta\omega)\gamma}\zeta.\tag{30}\end{equation}
This may be obtained from (29) or directly from the cyclic rule in its original form thus:
\[
C_\omega\gamma=C_{V_2\alpha\beta}\gamma=C_{V_2\alpha\gamma}\beta-C_{V_2\beta\gamma}\alpha=C_{V_2(V_1\zeta\omega)\gamma}\zeta
\]
by (1) above.

\indent It is useful to modify the cyclic rule for a self-conjugate $\varpi_x$, namely, $\varpi_0,\varpi_a,\varpi_b$, or $\varpi_c$. Thus, using $\varpi_a$ as an example, we have $\sum V_1\alpha\varpi_aV_2\beta\gamma=0$, or operating by $V_0\delta(\ )$, where $\delta$ is a fourth vector,
\[
V_0\mathbin{\cdot}V_2\alpha\beta\varpi_aV_2\gamma\delta+V_0\mathbin{\cdot}V_2\alpha\gamma\varpi_aV_2\delta\beta+V_0\mathbin{\cdot}V_2\alpha\delta\varpi_aV_2\beta\gamma=0.
\]
The expression on the left is combinatorial in $\alpha,\beta,\gamma,\delta$. Taking $\alpha,\beta,\gamma,\delta$ to be any four of the primitive vectors $\iota_1,\iota_2,\ldots,\iota_n$, say, $\iota_1,\iota_2,\iota_3,\iota_4$, we have successively the following three statements:---
\[
\sum\bigl[\iota_1^{-1}\iota_2^{-1}\iota_3^{-1}\iota_4^{-1}\sum V_0(\iota_1\iota_2)\varpi_a(\iota_3\iota_4)\bigr]=0,
\]
\[
V_4\zeta_1\zeta_2\zeta_3\zeta_4\mathbin{\cdot}V_0(V_2\zeta_1\zeta_2)\varpi_a(V_2\zeta_3\zeta_4)=0,
\]
\begin{equation}V_4(V_2\zeta_1\zeta_2)\varpi_a(V_2\zeta_3\zeta_4)=0,\tag{31}\end{equation}
as mentioned above between equations (12) and (13).
From equation (27) onwards we have been at no pains to put our results in invariant form. Nevertheless (27), (30), and (31) are already in invariant form, and (29) yields
\begin{equation}C_\omega\gamma=CV_1\omega\eta\gamma-C_{V_2(V_3\zeta\omega)\eta\gamma}\eta^{-1}\zeta.\tag{29a}\end{equation}
\indent \S17. \emph{Comparison with Theory of Tensors.}---Let $\phi(q_1,q_2,\ldots,q_a)$ be any contravariant or covariant function linear in each of the multenions $q_1,q_2,\ldots,q_a$, each one of which is also contravariant or covariant, not necessarily all of one type. Thus, for instance, $\Omega_{V_2\alpha\beta}\eta^{-1}\gamma^\times$ is such a covariant function of $\alpha,\beta,\gamma^\times$, two contravariant vectors, and one covariant vector. In the pages above we have defined various increments of $q$ due to changing $\rho$ to $\rho+\alpha$, where $\alpha$ is an arbitrary infinitesimal contravariant vector
\[
(d_\alpha^T,d_\alpha^I,d_\alpha^0,d_\alpha^A)q
\]
and we might add, when $q$ is of power $c$, \emph{the absolute increment proper}, also of power $c$, $d_\alpha^{AP}q=d_\alpha^Aq=(d_\alpha^A+cV_0\alpha\lambda)q$. The increments of $\phi$ of these types are all formed on the same principle. Thus take $d_\alpha^A$ as the type and use the abbreviated notation $d_\alpha^Aq=q_\alpha$. Then the increment $\phi_\alpha$ of $\phi$ is given by
\[
\phi_\alpha(q_1,q_2,\ldots,q_a)
\]
\[
=[\phi(q_1,q_2,\ldots,q_a)]_\alpha-\sum_{c=1}^a\phi(q_1,q_2,\ldots,q_{c\alpha},\ldots,q_a)
\]
where $q_{c\alpha}$ occupies the position of $q_c$ in $\phi(q_1,q_2,\ldots,q_a)$. Once having obtained our increment, we may cease to regard $\alpha$ as infinitesimal, so that $d_\alpha^A$ becomes, not so much an \emph{increment} as a \emph{rate of increase} multiplied by the magnitude of $\alpha$, though above it has always been called an increment.

\indent All this is on the model of the Theory of Tensors. In that theory, however, the only linear form $\phi(q_1,q_2,\ldots,q_a)$ recognised is the scalar form $\phi(\alpha_1,\alpha_2,\ldots,\alpha_a)$ in several vectors, some of which are covariant and the rest contravariant. The present writer marvels at the great results obtained by expounders of the Theory of Tensors, since the basis is such a limited concept.

\indent The precise translation from our presentation to that of the Theory of Tensors and \emph{vice vers\^a} is to a certain extent arbitrary, on account of vagueness, especially in the matter of sign, that writers on the Theory of Tensors seem to permit themselves. The following will be found to form one consistent system. Let $A^a,B^b$ be two contravariant vectors, $\alpha$ and $\beta$, and let $X_{ab}$ be a covariant tensor of the second rank, with the same intrinsic meaning as $\xi$, a covariant vector linity. $A_aB_b$ seems to be taken to mean the same as $B_bA_a$ in the Theory of Tensors, but I cannot see how this can be consistently done, and in the system here given $A_aB_b$ and $B_bA_a$, when interpreted as vector linities, must have different meanings.
With the basis as given in \S11 above, we have as follows. Let
\begin{equation}\alpha^\times=\eta\alpha,\qquad \beta^\times=\eta\beta.\tag{1}\end{equation}
Then
\begin{equation}\xi^\times=\eta^{-1}\xi\eta^{-1},\qquad X=\xi\eta^{-1},\qquad X^\times=\eta^{-1}\xi.\tag{2}\end{equation}
\begin{equation}\nabla=-\sum_{a=1}^n\iota_a^{-1}\frac{\partial}{\partial x_a},\qquad \frac{\partial}{\partial x_a}=-V_0\iota_a\nabla.\tag{3}\end{equation}
\begin{equation}A^a=V_0\iota_a^{-1}\alpha,\qquad A_a=V_0\iota_a\alpha^\times,\tag{4}\end{equation}
\begin{equation}g_{ab}=V_0\iota_a\eta\iota_b,\qquad g^{ab}=V_0\iota_a^{-1}\eta^{-1}\iota_b^{-1},\qquad g_b^a=V_0\iota_a\iota_b^{-1},\tag{5}\end{equation}
\begin{equation}[bc,a]=V_0\iota_a\bar\Gamma_b\iota_c,\qquad \{bc,a\}=V_0\iota_a^{-1}\bar E_b\iota_c,\tag{6}\end{equation}
\begin{equation}X_{ab}=V_0\iota_a\xi\iota_b,\qquad X_a^b=V_0\iota_aX\iota_b^{-1},\qquad X^{ab}=V_0\iota_a^{-1}\xi^\times\iota_b^{-1}.\tag{7}\end{equation}
The equation $X_{ab}=A_aB_b$ consistently with the above and the related equations may be interpreted to mean as follows
\begin{equation}X_{ab}=A_aB_b,\qquad \xi=\eta\alpha V_0(\ )\eta\beta=\alpha^\times V_0(\ )\beta^\times,\tag{8}\end{equation}
\begin{equation}X_a^b=A_aB^b,\qquad X=\eta\alpha V_0(\ )\beta=\alpha^\times V_0(\ )\beta.\tag{9}\end{equation}
\begin{equation}X^{ab}=A^aB^b,\qquad \xi^\times=\alpha V_0(\ )\beta.\tag{10}\end{equation}
Now (8) and (10) definitely tell us what $B_bA_a$ and $B^bA^a$ must mean; instead of representing $\xi$ and $\xi^\times$ respectively, they represent $\xi'$ and $\xi^{\times\prime}$ and $B^bA_a$ is similarly naturally interpreted as $\eta\beta V_0(\ )\alpha$, that is as $\eta X'\eta^{-1}$.

\indent Lastly we have
\begin{equation}(\xi,\zeta)=-\sum_{a=1}^n(\iota_a,\iota_a^{-1}).\tag{11}\end{equation}
It is with these interpretations that equation (25), \S16, was stated to be equivalent to Eddington's (26$'$3).

\indent I may add here that Eddington's equation (23) for $B_{abc}^h$ gives
\begin{equation}B_{abc}^h=-V_0\iota_h^{-1}\bar C_{V_2\iota_b\iota_c}\iota_a.\tag{12}\end{equation}
where $\bar C_\omega$ means $C_{\omega a}$, that is the part of $C_\omega$ obtained by putting $\lambda=0$.

\indent It will be seen that I interpret the mixed tensor to mean a comixt, not contramixt, vector linity.

\indent \S18. \emph{Prof. Eddington's recent work.}{\renewcommand{\thefootnote}{*}\footnote{‘Roy. Soc. Proc.,’ A, vol. 99, p. 104 (1921).}} [Re-written September, 1921.]

\indent Suppose $ABC$ is a triangle whose sides $AB,BC,CA$ are geodesics. Let the point $O$ of an infinitesimal lattice-work pass round the triangle, starting and ending at $A$, and in each stage of the journey let the lattice-work suffer ``parallel displacement'' only. When the circuit has been completed, then:---according to Euclid and Galileo the lattice-work will be found occupying its initial position exactly; according to Riemann it will be found turned through an angle in a plane containing $O$; according to Weyl it will
be turned and changed in size, but not in shape; according to Eddington it will be found turned, changed in size and also in shape. If Prof. Eddington meant no more than this, his manifold, as will be shown below, would be mathematically equivalent to Weyl's, but he does mean more. In his primary manifold \emph{there is no criterion for shape at all}, none to distinguish a square from any other parallelogram, or to distinguish a cube from any other parallelepiped. In Riemann's manifold the fundamental idea is that of the invariant quadratic form (squared interval), and ``parallel displacement'' is a mathematically \emph{derived} idea. In Eddington's manifold the \emph{fundamental} idea is that of parallel displacement, and an invariant quadratic form is a mathematically derived idea. When $n=4$, Riemann's, Weyl's and Eddington's foundations require ten, thirteen (nine \emph{ratios} of coefficients of quadratic form and four coefficients of a linear form) and forty scalars at a point to be given, respectively.

\indent Eddington then starts with $E_\alpha\beta(=E_\beta\alpha,-E_\alpha\beta$ meaning increment of $\beta$ due to the parallel displacement $\alpha;\ \alpha,\beta$ being contravariant vectors) but without any fundamental $\eta$ specifying the intrinsic invariant $V_0d\rho\eta d\rho$. \emph{All our previous work and results, which are independent of $\eta$, or can be made so, continue to be valid in Eddington's manifold.} On looking back I only see one case where our old reasoning ceases to hold in the new circumstances. This is equation (4), \S15, which was proved by aid of $\eta$. But the equation is still true because the condition that $\tau$, a contravariant vector, suffers the parallel displacement $\beta$ is
\[
-V_0\beta\nabla\mathbin{\cdot}\tau=-E_\beta\tau,
\]
and this is an intrinsic condition independent of choice of coordinates. Infinitesimally change the coordinates as in \S15 [$\delta\tau=\chi_0\tau,\ \delta\beta=\chi_0\beta,\ \delta(V_0\beta\nabla)=0$]. Thus
\[
-V_0\beta\nabla\mathbin{\cdot}(\chi_0\tau)=-\delta(E_\beta\tau);
\]
and here the left transforms to
\[
-\chi_0(V_0\beta\nabla\mathbin{\cdot}\tau)-V_0\beta\nabla_9\mathbin{\cdot}\chi_{09}\tau=-\chi_0E_\beta\tau-V_0\beta\nabla_9\mathbin{\cdot}\chi_{09}\tau
\]
so that
\[
\delta(E_\beta\tau)=\chi_0E_\beta\tau+V_0\beta\nabla_9\mathbin{\cdot}\chi_{09}\tau.
\]
This is equation (4) \S15 with $\tau$ in place of $\gamma$. We may then make the very important use of that equation which we made before, namely, if $\tau$ is \emph{any} contravariant vector,
\[
d_\alpha^A\tau=-V_0\alpha\nabla\mathbin{\cdot}\tau+E_a\tau
\]
is contravariant and may therefore be appropriately called the \emph{absolute increment} of $\tau$. [Prof. Eddington does not show in his paper that absolute differentiation thus holds good in the ``in-'' sense in his manifold.]
\indent Our chief disability from absence of a fundamental $\eta$ seems to be that there is now found no place for the ideas associated with normal and incident components (and expressions) and with rotations. $V_{\alpha+\beta}\omega$ and $V_{\alpha-\beta}\omega$ have their invariantive equivalents, but apparently $V_{\alpha+\beta-2c}\omega$ in general has not. We can still speak of $\alpha$ and $\beta^\times$ being normal when $V_0\alpha\beta^\times=0$ and this leaves us with (what we met with in \S10 before $\eta$ was introduced) the ideas of normal and incident components with reference not to a single $\tau$ but with reference to $\tau_c$, some one of $n$ given independent contravariant vectors $\tau_1,\tau_2,\ldots\tau_n$. In a word the normal reciprocals of $\tau_1,\tau_2,\ldots\tau_n$ continue to have invariantive meanings.

\indent For present purposes the contraction we selected in \S16 to pass from $C_\omega$ to $\Omega$ was unhappily chosen as it depended on $\eta$. We will here, therefore, change the contraction to another. The previous definition of $\Omega\beta$ was that it was obtained by setting $\alpha,\gamma$ equal to $\zeta,\eta^{-1}\zeta$ in $\eta C_{V_2\alpha\beta}\gamma$; our new meaning of $\Omega'\alpha$ is obtained by setting $\beta,\gamma^\times$ equal to $\zeta,\zeta$ in $C_{V_2\alpha\beta}'\gamma^\times$. Thus
\begin{equation}
\left.\begin{aligned}
\text{Previous}\quad \Omega\beta&=\eta C_{V_2\zeta\beta}\eta^{-1}\zeta\\
\text{Present}\quad \Omega'\alpha&=C_{V_2\alpha\zeta}'\zeta
\end{aligned}\right\}.\tag{1}
\end{equation}
\indent In Weyl's manifold the change in meaning is slight. The old and new self-conjugate parts are the same, but whereas the previous skew part of $\Omega\beta$ was $\tfrac14(n-4)V_1\omega^\times\beta$, the present is $-\tfrac14nV_1\omega^\times\beta$. This simplicity of connections between old and new meanings may be proved from the identity [equation (8), \S16].

\[
\eta C_{V_2\alpha\beta}\gamma+C_{V_2\alpha\beta}'\eta\gamma=-V_0\omega^\times V_2\alpha\beta\mathbin{\cdot}\eta\gamma,
\]
by setting $\alpha,\gamma$ equal to $\zeta,\eta^{-1}\zeta$. The identity is the one which expresses that our lattice-work, after its journey round the circuit $ABC$, is changed in size but not in shape.

\indent Our equation for $C_\omega$ (4) of \S16 still holds while $\eta$ is unassigned, and from it
\begin{equation}C_\omega'\gamma^\times=-E_\zeta\mathbin{\cdot}V_{1\omega}\nabla_\zeta'\gamma^\times+(-E_\beta'E_\alpha'+E_\alpha'E_\beta')\gamma^\times.\tag{2}\end{equation}
Putting here $V_{1\omega}\nabla=\alpha V_0\beta\nabla-\beta V_0\alpha\nabla$ and using (1) we get
\begin{equation}\Omega'\alpha=E_{\zeta,\alpha}'\nabla_\zeta+V_0\alpha\nabla\mathbin{\cdot}E_\zeta'\zeta-E_\zeta E_\alpha'\zeta+E_\alpha'E_\zeta'\zeta.\tag{3}\end{equation}
The first, third and fourth terms on the right are all self-conjugate because $E_\alpha\beta=E_\beta\alpha$. Hence the skew part of $\Omega\alpha$ is $-\tfrac12V_1\alpha V_2\nabla E_\zeta'\zeta$.

\indent By its original meaning $C_\omega\gamma$ is the difference of two truly contravariant, or as Eddington says in-contra variant (would not \emph{preter}-contravariant, be better?) vectors at the same point. Hence $C_\omega$, $C_\omega'$, $\Omega$ and $V_2\nabla E_\zeta'\zeta$ all have the ``in-'' property and $\tfrac12(\Omega+\Omega')$ is fitted for serving as the fundamental $\eta$ of a Riemann manifold. Prof. Eddington, therefore, identifies
$\tfrac12cV_0d\rho(\Omega+\Omega')d\rho$ with Einstein's quadratic form, the squared interval; $c$ being an absolute constant. I do not propose to follow him in this respect, but to show instead, that as soon as it is recognized that ``lengths'' can be compared (even in Weyl's limited sense that equality of lengths can be postulated \emph{at} a point only), then Eddington's theory is mathematically equivalent to Weyl's; and instead of merely one simple quadratic form there are several available to identify with Einstein's.

\indent It is now necessary to distinguish between Weyl's and Eddington's meanings of $E_a,C_\omega,\Omega$ and related symbols. For Eddington's meanings we shall use ${}^E E_a,{}^E C_\omega,{}^E\Omega$. For Weyl's we shall use ${}^W E_a,{}^W C_\omega,{}^W\Omega$.

\indent The following three modifications of notations occurring in \S\S14 to 16 above are desirable; (1) The letter $G$ will be used in place of the letter $F$, because $F$ is wanted below for another purpose; (2) the ``contraction'' to pass from $C_\omega$ to $\Omega$ (both of power zero), will be altered in meaning as already described; (3) the symbols $2,1,0,a,b,c$, used as suffixes to $C_\omega,\Omega_\omega,C,\Omega,\omega,G$ (former $F$), will be shifted from the right-hand bottom corner of the letter affected to the left-hand top corner. The following scheme will probably be helpful:---
\[
\begin{tikzpicture}[baseline=(current bounding box.center),x=1cm,y=.62cm,
  every node/.style={font=\small,inner sep=1pt}]
  \node (e) at (1.15,4) {${}^E C_\omega$};
  \node (four) at (-3.45,2.35) {${}^4 C_\omega$};
  \node at (-3.45,1.5) {$(144)$};
  \node (three) at (0,2.35) {${}^3 C_\omega$};
  \node at (0,1.5) {$(36)$};
  \node (w) at (3.45,2.35) {${}^W C_\omega$};
  \node (two) at (2,0.7) {${}^2 C_\omega$};
  \node (one) at (4.175,0.7) {${}^1 C_\omega$};
  \node (zero) at (6.35,0.7) {${}^0 C_\omega$};
  \node (a) at (4.1,-1.28) {${}^a C_\omega$};
  \node (b) at (7,-1.28) {${}^b C_\omega$};
  \node (c) at (10.3,-1.28) {${}^c C_\omega$};
  \draw (e.south) -- (1.15,3.25);
  \draw (-3.45,3.25) -- (3.45,3.25);
  \draw (-3.45,3.25) -- (-3.45,2.8);
  \draw (3.45,3.25) -- (3.45,2.8);
  \draw (0,3.25) -- (0,2.8);
  \draw (w.south) -- (3.45,1.6);
  \draw (2,1.6) -- (6.35,1.6);
  \draw (2,1.6) -- (2,1.15);
  \draw (6.35,1.6) -- (6.35,1.15);
  \draw (3.45,1.6) -- (3.45,1.15);
  \draw (4.175,1.6) -- (4.175,1.15);
  \draw[decorate,decoration={brace,amplitude=4pt,mirror}] (two.south) -- (one.south)
    node[midway,below=7pt] {$(6)$};
  \node (twenty) at (6.35,-.05) {$(20)$};
  \draw (twenty.south) -- (6.35,-.5);
  \draw (4.1,-.5) -- (10.3,-.5);
  \draw (4.1,-.5) -- (4.1,-.95);
  \draw (10.3,-.5) -- (10.3,-.95);
  \draw (7,-.5) -- (7,-.95);
\end{tikzpicture}
\]

\indent [In each of these four lines there is one symbol equivalent to a tensor of Eddington's, or to a tensor of Weyl's. ${}^E C_\omega$ is equivalent to Eddington's ${}^*B_{abc}^{\ h}$, ${}^W C_\omega$ to Weyl's $F_{abc}^{\ h}$, ${}^0 C_\omega$ to Weyl's ${}^*F_{abc}^{\ h}$, and ${}^a C_\omega$ to both Eddington's $B_{abc}^{\ h}$ and Weyl's $R_{abc}^{\ h}$.] In the scheme the $E$ of ${}^E C_\omega$ implies ``in Eddington's sense''; the $W$ of ${}^W C_\omega$ implies ``in Weyl's sense''; further, the scheme implies the following three defining equations:---
\[
{}^E C_\omega={}^4 C_\omega+{}^3 C_\omega+{}^W C_\omega,
\]
\[
{}^W C_\omega={}^2 C_\omega+{}^1 C_\omega+{}^0 C_\omega,
\]
\[
{}^0 C_\omega={}^a C_\omega+{}^b C_\omega+{}^c C_\omega.
\]
We may write $x$ for any one of the ten symbols $E,W,4,3,2,1,0,a,b,c$. For seven out of the ten meanings of ${}^xC_\omega$, namely, for all except ${}^aC_\omega,{}^bC_\omega,{}^cC_\omega$, it is true that ${}^xC_\omega$ is of zero power as meant by Weyl. The three ${}^aC_\omega,{}^bC_\omega,{}^cC_\omega$, are invariantive for change of co-ordinates but not for change of gauge.

\indent The numbers in brackets $(144)$, etc., imply that ${}^E C_\omega$ involves, when $n=4$, $206(=144+36+6+20)$, independent scalars, that ${}^W C_\omega$ involves $26(=6+20)$, independent scalars, and so on. [It is explained below that
the statement that ${}^3C_\omega$ involves 36 scalars, merely means that ${}^3C_\omega$ is fully known when the 36 scalars of $Y_{\alpha\beta}$ below are known; and, again, ${}^4C_\omega$ is fully known when the 144 first derivatives of the 36 scalars are known. In a purely algebraic sense ${}^4C_\omega$ cannot involve more than 96 scalars.]

\indent In what now follows Eddington's symbols
\[
\Gamma_{ab}^{\ c},\quad {}^*B_{abc}^{\ h},\quad {}^*G_{ab},\quad R_{ab},\quad F_{ab},\quad K_{ab;c},\quad S_{ab;c},\quad 2\kappa_a,
\]
are the equivalents of our
\[
\mathbf{E}_a,\quad {}^E C_\omega,\quad {}^E\Omega,\quad \tfrac12({}^E\Omega+{}^E\Omega'),\quad \tfrac12({}^E\Omega-{}^E\Omega'),\quad \Psi_{\beta\gamma},\quad \Lambda_{\beta\gamma},\quad \mathbf{E}_\zeta'\zeta,
\]
respectively.

\indent In our interpretations, Eddington's ``increment due to parallel displacement'' has to be distinguished from our (for Weyl's manifold), ``increment due to translation.'' We use $d_\alpha^{\mathrm{PD}}$ and $d_\alpha^{\mathrm{T}}$ for these phrases. As with (5) of \S14 above, we have
\[
d_\alpha^{\mathrm{PD}}V_0\beta\eta\gamma=-V_0\alpha\nabla_9\mathbin{\cdot}V_0\beta\eta\gamma-V_0\mathbf{E}_\alpha\beta\eta\gamma-V_0\beta\eta\mathbf{E}_\alpha\gamma.
\]
Prof. Eddington remarks that as this is the difference of two invariants it must be invariant, and therefore must be of the form $2V_0\alpha\Psi_{\beta\gamma}$ (a more general form than ours, after Weyl, $-V_0\beta\eta\gamma\mathbin{\cdot}V_0\lambda\alpha$ in equation (4) of \S14), where $\Psi_{\beta\gamma}$ is a covariant vector linear and symmetrical in $\beta$ and $\gamma$ (contravariant vector dummies), given by
\begin{equation}2\Psi_{\beta\gamma}=-\nabla_9V_0\beta\eta\gamma-\mathbf{E}_\beta'\eta\gamma-\mathbf{E}_\gamma'\eta\beta.\tag{4}\end{equation}
\indent $\Psi_{\beta\gamma}$ can easily be expressed uniquely as the sum of two parts, a contractile part involving $n$ scalars and a non-contractile part involving $\tfrac12(n+2)n(n-1)$ scalars. The contractile part is $-n^{-1}\Psi_{\zeta\eta^{-1}\zeta}\mathbin{\cdot}V_0\gamma\eta\beta$, and is so named because it is the sole part of $\Psi_{\beta\gamma}$ which is responsible for the contraction $\Psi_{\zeta\eta^{-1}\zeta}$, a vector. Thus
\begin{equation}2\Psi_{\beta\gamma}=-\lambda V_0\beta\eta\gamma+2\Upsilon{T}_{\beta\gamma}.\tag{5}\end{equation}
where $\lambda$ is a covariant vector (namely, $2n^{-1}\Psi_{\zeta\eta^{-1}\zeta}$), and $\Upsilon{T}_{\beta\gamma}$ is a covariant vector, linear and symmetrical in $\beta$ and $\gamma$, and satisfying the vector equation
\begin{equation}\Upsilon{T}_{\zeta\eta^{-1}\zeta}=0.\tag{6}\end{equation}
\indent [Digression.---If $Z_{a\beta}$ is contravariant (and similar remarks hold when it is covariant), and does not satisfy the equation of symmetry $Z_{a\beta}=Z_{\beta a}$, there are three kinds of contraction (derived from the scalar form $V_0\gamma^\times Z_{a\beta}$). The corresponding contractile parts of $Z_{a\beta}$ are
\[
-\eta^{-1}\zeta_1\mathbin{\cdot}V_0\zeta_2Z_{\zeta_1\zeta_2}\mathbin{\cdot}V_0\alpha\eta\beta,\qquad \eta^{-1}Z_\zeta'\zeta\mathbin{\cdot}V_0\alpha\eta\beta,\qquad -\eta^{-1}Z_{\zeta\eta^{-1}\zeta}\mathbin{\cdot}V_0\alpha\eta\beta,
\]
the three contractions themselves being
\[
-\zeta_1\mathbin{\cdot}V_0\zeta_2Z_{\zeta_1\zeta_2},\qquad Z_\zeta'\zeta,\qquad Z_{\zeta\eta^{-1}\zeta}.
\]
The first and second of these contractions are equal when $Z_{a\beta}=Z_{\beta a}$.]
\indent The $\lambda$ appearing in equation (5) is the $\lambda$ we met with in Weyl's theory. If we examine what change occurs in $\Psi_{\beta\gamma}$, or better in $\eta^{-1}\Psi_{\beta\gamma}$, due to change of gauge, we at once find that in $\eta^{-1}\Psi_{\beta\gamma}$ the whole change falls on the $\lambda$ part, no change at all occurring in $\eta^{-1}\mathfrak{T}_{\beta\gamma}$; that is $\eta^{-1}\mathfrak{T}_{\beta\gamma}$ is of zero power.

\indent Calculating from (4) the value of
\[
V_0\alpha\Lambda_{\beta\gamma}=V_0\alpha\Psi_{\beta\gamma}-V_0\beta\Psi_\gamma'\alpha-V_0\gamma\Psi_\alpha'\beta
\]
we at once find that
\begin{equation}\Lambda_{\beta\gamma}=(\Psi_{\beta\gamma}-\Psi_\beta'\gamma)-\Psi_\gamma'\beta\tag{7}\end{equation}
and
\begin{equation}\mathbf{E}_{\beta\gamma}=E_{\beta\gamma}+Y_{\beta\gamma}\tag{8}\end{equation}
where
\begin{equation}Y_{\beta\gamma}=\eta^{-1}(\mathfrak{T}_{\beta\gamma}-\mathfrak{T}_\beta'\gamma-\mathfrak{T}_\gamma'\beta)\tag{9}\end{equation}
and $E_{\beta\gamma}$ is exactly as in \S14. On the right of (7) the two bracketed terms furnish the skew part of $\Lambda_\beta$ and the third term the self-conjugate part. It will be easily seen that
\begin{equation}Y_{\zeta'\zeta}=-\mathfrak{T}_{\zeta\eta^{-1}\zeta}=0.\tag{10}\end{equation}

\indent From (9) $Y_{\beta\gamma}$ is of zero power. It follows that it is just as legitimate to name $-E_{\beta\gamma}$ the increment of $\gamma$ due to the translation $\beta$ ($d_\beta^{\mathrm T}\gamma$) as it is to name $-{}^E\mathbf{E}_{\beta\gamma}$ the increment of $\gamma$ due to the parallel displacement $\beta$ ($d_\beta^{\mathrm{PD}}\gamma$). In other words Eddington's manifold is simply Weyl's in which there is supposed to exist a given $Y_{\beta\gamma}$ of zero power and satisfying the equation $Y_{\zeta'\zeta}=0$.

\indent With our present notation [see (4) of \S16]
\begin{equation}{}^E C_\omega\gamma=d_\omega^{\mathrm{PD}}\gamma=-\mathbf{E}_9\mathbin{\cdot}V_{1\omega}\nabla_9\gamma+(-\mathbf{E}_\alpha\mathbf{E}_\beta+\mathbf{E}_\beta\mathbf{E}_\alpha)\gamma.\tag{11}\end{equation}

\indent Putting $\mathbf{E}_\alpha=E_\alpha+Y_\alpha$ we get a part involving $E_\alpha,E_\beta$ only. This is, of course, our former
\begin{equation}{}^W C_\omega\gamma=d_\omega^{\mathrm T}\gamma=-E_9\mathbin{\cdot}V_{1\omega}\nabla_9\gamma+(-E_\alpha E_\beta+E_\beta E_\alpha)\gamma.\tag{12}\end{equation}
\indent Next we get a second part
\begin{equation}{}^3C_\omega\gamma=(-Y_\alpha Y_\beta+Y_\beta Y_\alpha)\gamma\tag{13}\end{equation}
and a third part which may be written
\begin{equation}{}^4C_\omega\gamma=-(Y_\alpha\gamma)_\beta+(Y_\beta\gamma)_\alpha\tag{14}\end{equation}
where $(\ )_\beta$ and $(\ )_\alpha$ denote absolute increments in the sense developed in \S15 above for Weyl's manifold. (14) follows from the beginning of \S17 from which is deducible
\[
(Y_\alpha\gamma)_\beta=-Y_{\vartheta\alpha}\gamma\mathbin{\cdot}V_0\beta\nabla_\vartheta+E_\beta Y_\alpha\gamma-Y_\alpha E_{\beta\gamma}-Y_\gamma E_{\beta\alpha}.
\]
\indent With $n=4$ (13) involves the thirty-six scalars of $Y_{\alpha\beta}$ and (14) involves their 144 first derivatives $(\nabla)$. The full meaning of our scheme for ${}^E C_\omega$, given above, will now be evident.

\indent There are numerous simple quadratic forms available to identify, in Prof.\ Eddington's manner, with the square of Einstein's interval. Not only
may we contract (12) and select its self-conjugate part ${}^0\Omega$, but we may similarly contract each of (13) and (14) in two different ways and this by no means exhausts our available simple forms. Besides (13) and (14) we may similarly use $(Y_\alpha Y_\beta+Y_\beta Y_\alpha)\gamma$ and $(Y_\alpha\gamma)_\beta+(Y_\beta\gamma)_\alpha$ where yet another kind of contraction is available (put $\alpha,\beta=\zeta,\eta^{-1}\zeta$).

\indent Returning to the use we made of (3) above, we see that the skew part of ${}^E\Omega\alpha$ is $-\tfrac12V_1\alpha V_2\nabla\mathbf{E}_\zeta'\zeta$. Since $Y_{\zeta'\zeta}=0$,
\begin{equation}\mathbf{E}_\zeta'\zeta=E_\zeta'\zeta=\nabla\log k-\tfrac12n\lambda,\tag{15}\end{equation}
by (18) of \S15 above. Hence Prof.\ Eddington's identification of $V_2\nabla\mathbf{E}_\zeta'\zeta$ in the electromagnetic field is virtually the same as Weyl's identification of $V_2\nabla\lambda$. Also note that ${}^E\Omega-{}^W\Omega(={}^3\Omega+{}^4\Omega)$ is self-conjugate.

\indent One consequence of Eddington's suggestions is that we may quite likely desire to identify Einstein's squared interval with something very different from our fundamental form $V_0d\rho\eta d\rho$. Let Einstein's squared interval be $V_0d\rho\theta d\rho$. We may wish to put $\theta=\tfrac12({}^E\Omega+{}^E\Omega')$ or $\theta={}^0\Omega=\tfrac12({}^W\Omega+{}^W\Omega')$ or $\theta=\tfrac12({}^3\Omega+{}^3\Omega')$, the last making $\theta$ depend solely on $Y_{\alpha\beta}$. With any such identification we require symbols dependent on $\theta$, as are $k,{}^W\Omega$, etc., on $\eta$. We shall, therefore, henceforth use
\[
\theta,\quad l,\quad F_\beta,\quad \Theta_\beta,\quad C_\omega,\quad \Omega,\quad \bar\omega,\quad G,
\]
in the $\theta$ manifold (Riemann manifold) in the same way as above
\[
\eta,\quad k,\quad E_\beta,\quad \Gamma_\beta,\quad {}^W C_\omega,\quad {}^W\Omega,\quad {}^W\bar\omega,\quad {}^W G,
\]
have been used in the $\eta$ manifold (Weyl manifold). The relations between $\theta,l,F_\beta$, etc. are precisely the same as those between
\[
\eta,\quad k,\quad {}^aE_\beta,\quad {}^a\Gamma_\beta,\quad {}^aC_\omega,\quad {}^a\Omega,\quad {}^a\bar\omega,\quad {}^aG,
\]
(${}^aE_\beta,{}^a\Gamma_\beta$ here replacing our earlier $\bar E_\beta,\bar\Gamma_\beta$).

\indent If $\theta$ is identified in some such manner as just suggested the determinant $|\theta|$ may become zero. $|\theta|=0$ and $|\theta^{-1}|=0$ represent loci. Do the loci if really existent, constitute the boundaries of electrons and atomic nuclei? On crossing the boundary do we pass to a region with four real dimensions or to a region of two real and two imaginary dimensions (two-dimensional time and two-dimensional space)? Is it on the other hand meaningless to speak about crossing the boundary, much as to an inhabitant of hyperbolic space it would be meaningless to speak of passing across his space's boundary, the absolute?

\noindent\textsc{§ 19. \emph{Suggested Improvement of Some Notations and Terms.}}---The reader has probably found that the terminology for multenion linities, adopted hitherto because of the writer's timid reluctance to depart from accepted meanings, is too inexpressive of the intended signification. A covariant linity fabricates a covariant multenion out of a contravariant. Let us call it a $c/\invc$ or a coexcontra (co-ex-contra); and similarly let us now denote a contravariant, comixt or contramixt linity by the name contraexco, coexco or contraexcontra, and by the symbol $\invc/c$, $c/c$ or $\invc/\invc$ respectively. Let us also use parenthesis marks\footnote{Intended by author to be \emph{commas}, upright and inverted. In future printing the author still thinks commas would be preferable to the bracket substitute.} before a letter such as $\phi$, of the same general shapes as $c$ and $\invc$ (omitting the solidus), to indicate the type of invariance, thus---let us say that
\[
  \phi=\markpair{(}{)}{\phi},\ \markpair{)}{(}{\phi},\ \markpair{(}{(}{\phi}\ \text{or}\ \markpair{)}{)}{\phi}
\]
indicates that $\phi$ is a
\[
  c/\invc,\quad \invc/c,\quad c/c\ \text{or}\ \invc/\invc
\]
respectively. Similarly, $p={}^{(}p$ indicates that $p$ is a covariant multenion and $q={}^{)}q$ indicates that $q$ is a contravariant multenion. If we actually insert all these indicating parenthesis marks, which we shall do but rarely, their mode of combination is simple. Thus
\[
  {}^{)}p=\markpair{)}{)}{\phi}\ \markpair{)}{(}{\psi}\ {}^{(}q
\]
is an invariant equation. The features that render it so are (1) that the second mark of one letter, $\phi$ or $\psi$, is identical with the first mark of the next; these two marks virtually cancel one another; and (2) the proper mark of the resultant symbol $p$, of the whole operation, is the first written mark of the compound symbol $\phi\psi q$.

\indent When with these notations we use the dash which indicates conjugacy, there is, unless we insert brackets, ambiguity in the meanings of $\markpair{(}{(}{F'}$ and $\markpair{)}{)}{F'}$ (the two mixed types as we have hitherto called them), because $(\markpair{(}{(}{F})'$ is not of the same type $c/c$ as $\markpair{(}{(}{F}$, but of the type $\invc/\invc$. There is no such ambiguity with $\markpair{(}{)}{\phi'}$ and $\markpair{)}{(}{\phi'}$ nor with $f\markpair{(}{(}{F}$ and $f\markpair{)}{)}{F}$, where $fx$ is a rational function of $x$. [Remember that $f^{(}\phi$ and $f^{)}\phi$ have no invariantive meaning. I think it best to adhere to the general rule, not rigidly prescribed, that $c/\invc$ and $\invc/c$ should be denoted by Greek letters and $c/c$ and $\invc/\invc$ by Roman capitals.]

\indent Our new notations and terms apply to other symbols than invariantive; for instance, to $\Gamma_\beta$. Thus
\[
  \Gamma_\beta={}^{)}\Gamma_\beta,\qquad E_\beta'=\markpair{(}{(}{(E_\beta')}.
\]
We shall retain the ``cross'' notation to indicate different symbols, and generally restrict it, without rigid prescription, to the same meanings as hitherto.

\indent Our former $\bar\Gamma_\beta$ and $\bar E_\beta$ will henceforth not be used, because we have in §18 a new notation, according to which ${}^aE_\beta$ would have reference to a Weyl manifold, and would indicate that ${}^aE_\beta$ was the part of $E_\beta$ not containing terms in $\lambda$; and if $\lambda$ is actually everywhere zero so that the manifold becomes a Riemann manifold, §18 provides the use of $F_\beta$ in place of $E_\beta$. In the present paper we occupy ourselves exclusively with a Riemann manifold, and use the letters $\theta$, $l$, $F_\beta$, $\Theta_\beta$, $\Omega$, $C\ (=\markpair{(}{(}{C}=\Omega\theta^{-1})$ provided for that case.

\indent A thick dot\footnote{It is necessary, because of papers despatched to England in the early months of 1922, to state that this dot replaced a \emph{bar} placed over the symbol affected. The replacement was made in August, 1922, at the suggestion of the printers. It seems now (August) too late to replace the bar in the papers just mentioned.} will be used to denote that a symbol is a ``density,'' though occasionally, as with $W$ (action density), and $f$ a part of $W$ below, we shall not consider it necessary to use the dot. Once for all, we define $\dot\eta'$ and $\dot\theta$ by the equations
\begin{equation}
  \dot\eta'=k\eta^{-1},\qquad \dot\theta=l\theta^{-1}.
\tag{1}
\end{equation}
but we may regard it as a rule, unless the contrary is definitely stated, as here with $\dot\eta'$ and $\dot\theta$, that the dot attached to a symbol means that the undotted symbol has been multiplied by $l$ (by $k$ if we were dealing with a Weyl manifold). Thus $\dot C$, $\dot G$ stand for $lC$, $lG$, and again, if we introduce $\dot\kappa$ as a contravariant vector density, then $\kappa$ will mean $l^{-1}\dot\kappa$.

\indent With regard to $\dot\theta$ in (1), we shall use it as applied to a general multenion, say, to $w^\times=V_aq^\times$, but not exactly with the extensive meaning. Thus $\dot\theta u^\times$ means $l\theta^{-1}u^\times$ where $\theta^{-1}$ but not $l$ is used in the extensive meaning. Thus, for instance, $\dot\theta V_3\alpha^\times\beta^\times\gamma^\times$ means {\small$lV_3\theta^{-1}\alpha^\times\theta^{-1}\beta^\times\theta^{-1}\gamma^\times$}, and not {\small$l^3V_3\theta^{-1}\alpha^\times\theta^{-1}\beta^\times\theta^{-1}\gamma^\times$}.

\noindent\textsc{§ 20. \emph{Maxwell's Equations for Matter in Bulk.}}---It is fully time to show how the special requirements of Relativity are met in detail in our method. In now introducing Maxwell's equations, let us throughout keep in mind the application to matter in bulk. If for no other reason this is desirable because thereby we are prepared to accept the beautiful explanation that Maxwell's theory affords of the known behaviour of light in uniaxal and biaxal crystals. When we compare results with Maxwell's we shall speak as if $n=4$. Nevertheless, in the general argument, we shall continue to suppose $n$ to have any positive integral value. When with $n=4$ we put
\[
  \rho=xl_1+yl_2+zl_3+tl_4
\]
we shall speak of $x$, $y$, $z$ as the spatial co-ordinates and $t$ as the time, even when by transformation of all four co-ordinates $t$ may mean something very different from the time as ordinarily understood.

\indent Let Maxwell's vector potential be $\mathbf{A}=(A_1,A_2,A_3)$, electrostatic potential $=A_4$, magnetic induction $\mathbf{B}=(B_1,B_2,B_3)$, electric intensity $\mathbf{E}=(E_1,E_2,E_3)$, magnetic field strength $\mathbf{H}=(H_1,H_2,H_3)$, dielectric displacement $\mathbf{D}=(D_1,D_2,D_3)$, conduction current $\mathbf{K}=(K_1,K_2,K_3)$, electrostatic density $=K_4$, force per unit volume $\mathbf{F}=(F_1,F_2,F_3)$, power $(-S\mathbf{E}\mathbf{K})$ per unit volume $=F_4$.

\indent Also let $\lambda$ be a covariant vector (better use $\lambda^\times$ in the future) given by
\[
  \lambda=\sum A_cl_c\quad(c=1,2,3,4),
\]
$\omega^\times$ a covariant ${}_nV_2$ given by
\[
  \omega^\times=\sum B_1l_2l_3+\sum E_1l_4l_1
\]
(the meaning of each summation, of three terms, being sufficiently obvious), $\dot\omega$ a contravariant ${}_nV_2$ density given by
\[
  \dot\omega=\sum H_1l_2l_3+\sum D_1l_4l_1,
\]
$\dot\kappa$ a contravariant vector density given by
\[
  \dot\kappa=\sum K_cl_c\quad(c=1,2,3,4),
\]
and $\dot v^\times$ a covariant vector density given by
\[
  \dot v^\times=\sum F_cl_c\quad(c=1,2,3,4).
\]
Maxwell's quaternion equations
\begin{alignat}{3}
  \mathbf{B}&=V\nabla\mathbf{A}, &\qquad \mathbf{E}&=-\nabla A_4-\partial\mathbf{A}/c\partial t, &\tag{1}\\
  \mathbf{K}&=V\nabla\mathbf{H}-\partial\mathbf{D}/\partial t, & \mathbf{K}_4&=-S\nabla\mathbf{D}, &\tag{2}\\
  \mathbf{F}&=V\mathbf{K}\mathbf{B}+K_4\mathbf{E}, & \mathbf{F}_4&=-S\mathbf{K}\mathbf{E}, &\tag{3}
\end{alignat}
are identical with the multenion equations\footnote{For a convenient, simple, systematic method of converting all $(n=4)$ multenion formulae to corresponding quaternion formulae, see a paper by the author (despatched a month ago, viz., November, 1922, to England), ``The Mechanical Forces \ldots\ in an Electromagnetic Field.''}
\begin{alignat}{2}
  \omega^\times&=V_2\nabla\lambda, &\tag{1a}\\
  \dot\kappa&=V_1\nabla\dot\omega, &\tag{2a}\\
  \dot v^\times&=V_1\dot\kappa\omega^\times. &\tag{3a}
\end{alignat}
The equations (1a), (2a), (3a) are all invariantive.


\noindent [A digression. In Maxwell's presentation of the fundamental aspects of the electromagnetic field $V\nabla\mathbf{H}$ in equation (2) is looked upon as \emph{the} electric current, consisting of two parts, conduction current $\mathbf{K}$ and dielectric current $\partial\mathbf{D}/\partial t$. In the new Relativity scheme, as given by (2a), we naturally interchange the \emph{r\^oles} of $V\nabla\mathbf{H}$ and $\mathbf{K}$; that is to say, we look upon $\mathbf{K}$ as \emph{the} current, consisting of two parts, a magnetic part, $V\nabla\mathbf{H}$, and an electric part, $-\partial\mathbf{D}/\partial t$. In the multenion presentation (see below for $N$ and $I$) $V\nabla\mathbf{H}$ is $V_1\nabla(N\dot\omega)$, $-\partial\mathbf{D}/\partial t$ is $NV_1\nabla(I\dot\omega)$, and $K_4$ (multiplied into $l_4$) is $IV_1\nabla(I\dot\omega)$. [In the last sentence $N$ and $I$ are used as if $N=\markpair{)}{)}{N}$, $I=\markpair{)}{)}{I}$; if, as usual, we take $N=\markpair{(}{(}{N}$, $I=\markpair{(}{(}{I}$, then in the last sentence each $N$ should be $N'$ and each $I$ should be $I'$ .] In the future, then, when we use the term current, Maxwell's conduction current will be meant. Maxwell does not put $\mathbf{F}=V\mathbf{K}\mathbf{B}$ as we have asserted in (3), but instead $\mathbf{F}=V(V\nabla\mathbf{H})\mathbf{B}$.]

\indent Before speaking of the relations between Maxwell's $(\mathbf{B},\mathbf{E})$ and $(\mathbf{H},\mathbf{D})$, that is, between $\omega^\times$ and $\dot\omega$, let us derive Lorentz's generalizations of Maxwell's equations to bodies in motion from the general principle of Relativity. The generalisations in nowise alter (1a), (2a), (3a). Put $N_4$ and $I_4$ for the normal and incident components with reference to $l_4$, when the co-ordinates are Galilean. In reality we have
\begin{equation}
  \left.\begin{aligned}
    \mathbf{A}&\text{ is }N_4\lambda, &\qquad l_4A_4&\text{ is }I_4\lambda\\
    \mathbf{B}&\text{ is }N_4\omega^\times, & \mathbf{E}&\text{ is }I_4\omega^\times
  \end{aligned}\right\}
  \tag{1b}
\end{equation}
\begin{equation}
  \left.\begin{aligned}
    \mathbf{H}&\text{ is }N_4\dot\omega, &\qquad \mathbf{D}&\text{ is }I_4\dot\omega\\
    \mathbf{K}&\text{ is }N_4\dot\kappa, & l_4K_4&\text{ is }I_4\dot\kappa
  \end{aligned}\right\}
  \tag{2b}
\end{equation}
\begin{equation}
  \mathbf{F}\text{ is }N_4\dot v^\times,\qquad l_4F_4\text{ is }I_4\dot v^\times.
  \tag{3b}
\end{equation}
In the general $\theta$ (Riemann) manifold
\[
  \iota=dp/ds=dp/\sqrt{(V_0dp\theta dp)}
\]
degenerates when Galilean co-ordinates are used $(\theta=1)$ to
\[
  (l_1dx+l_2dy+l_3dz+l_4dt)/\sqrt{(-dx^2-dy^2-dz^2+dt^2)}
\]

\[
  =(1-\sum\dot x^2)^{-\frac12}\cdot(l_1\dot x+l_2\dot y+l_3\dot z+l_4),
\]
of which we say nothing more here than that when the co-ordinates are so chosen that the particle is at rest $\iota$ (the contravariant world unit tangent vector to the world locus of a particle), degenerates to $l_4$. Maxwell's equations, as given above, are to be supposed true only when the medium is at rest. The generalisations of $N_4$ and $I_4$ are thus $N$ and $I$ (or $N'$ and $I'$), the normal and incident components with reference to $\iota$; normal \emph{component} and normal \emph{expression}, in this case, meaning the same, because

$V_0\iota\theta\iota=1$. \emph{Below $N$ and $I$ will always mean what has just been described with reference to $\iota$,} unless explicit exception is made; (1b), (2b), (3b), become like (1a), (2a), (3a), all invariantive when we read $\iota$, $N$ and $I$ in place of $l_4$, $N_4$, and $I_4$, in (1a) and (3a); and in (2a), $\iota$, $N'$ and $I'$ in place of $l_4$, $N_4$, and $I_4$. [Remember that $N=\markpair{(}{(}{N}$, $N'=\markpair{)}{(}{N'}$, and similarly for $I$.]

\indent \emph{The equations this process gives when a body is moving, and we are using Galilean co-ordinates, are Lorentz's well-known generalizations of Maxwell's quaternion equations} (not of our multenion equations, which have been in no way altered by what has been said).

\indent From (1a) and (2a) we have of course
\begin{alignat}{2}
  V_3\nabla\omega^\times&=0, &\tag{4}\\
  V_0\nabla\dot\kappa&=0. &\tag{5}
\end{alignat}
These must be implicitly contained in (1) and (2) and indeed are found to be
\begin{alignat}{3}
  S\nabla\mathbf{B}&=0, &\qquad V\nabla\mathbf{E}+\partial\mathbf{B}/\partial t&=0, &\tag{4a}\\
  S\nabla\mathbf{K}&=\partial K_4/\partial t. & & &\tag{5a}
\end{alignat}

\indent Maxwell's relations between $(\mathbf{B},\mathbf{E})$ and $(\mathbf{H},\mathbf{D})$ are that $\mathbf{H}$ is a self-conjugate linity of $\mathbf{B}$, (permeability)$^{-1}$; and $\mathbf{D}$ a similar function of $\mathbf{E}$, capacity. In our case
\begin{equation}
  \left.\begin{aligned}
    \dot\omega&=\dot\psi\omega^\times, &\qquad \dot\psi&=\dot\psi'=\markpair{)}{(}{\dot\psi}\\
    I'\dot\psi N&=0=N'\dot\psi I
  \end{aligned}\right\}
  \tag{6}
\end{equation}
Here $\dot\psi$ is a ${}_nV_2$ linity concerning which, $\dot\psi=\dot\psi'$ asserts that $\dot\psi$ is self-conjugate, $\dot\psi=\markpair{)}{(}{\dot\psi}$ asserts that it is contraexco, the dot implies that it is a density, and the two equations in the second line assert that $\dot\psi$ consists of two parts $N'\dot\psi N+I'\dot\psi I$, of which the first is confined to normal components with reference to $\iota$, and the other to incident components.

\indent [Any $c/c$, say $\phi$, may be expressed uniquely as the sum of four parts thus
\[
  \phi=(N+I)\phi(N+I)=N\phi N+I\phi I+N\phi I+I\phi N.
\]
With a $\invc/c$ the four parts are obtained from $(N+I)'\ (\ )\ (N+I)$, and in the case of $\dot\psi$ two of these parts are asserted in (6) to be zero.

\indent \emph{Digression on I and N.}---I have forgotten to point out earlier that $I'$, $I$, $N'$ and $N$ all have subextensive significance, and that $(I-N)$ also has a close connection with extensive interpretation. In the neutral form we virtually deal with a Galilean manifold. If in such a manifold $I$, $N$ are the incident and normal components with reference to $a$, we can, by a rotation, bring $a$ to coincidence with $l_1$, or some $l_a$, and then $(I-N)$ becomes the capital linity $P_1$ (which reverses the sign of $l_1$ in every term of $q$), introduced at the outset of our subject. The properties of $I$, $N$, $(I-N)$, may be based on this connection

with $P_1$, through neutral to invariant forms. However, it is perfectly easy to prove directly that, say, $I'$, has subextensive significance, for we have defined thus
\[
  I'\beta=\alpha V_0\eta\alpha\beta,\qquad I'u=V_{a\alpha}V_{a-1}\eta\alpha u,
\]
and to show that $I'$ has subextensive significance we have to prove that
\[
  I'u=-V_aI'\xi V_{a-1}\xi u,
\]
which is obvious as soon as stated.]

\indent Maxwell's scheme also connects $\mathbf{K}$ with $\mathbf{E}$ by the self-conjugate linity conductivity. In our case express this first with Galilean co-ordinates. We have
\[
  \sum_{a=1}^{3}E_al_a=V_1l_4\omega^\times=RN_4\dot\kappa
\]
where $R$ is put for resistivity. This generalizes to
\begin{equation}
  N'\dot\kappa=\dot\phi V_1I\omega^\times,\qquad
  \dot\phi=N'\dot\phi N=\dot\phi'=\markpair{)}{(}{\dot\phi}.
  \tag{7}
\end{equation}
where $\dot\phi$ is conductivity, a vector linity.

\indent It appears to me an impressive fact that Maxwell's analysis of the electromagnetic field was so wonderfully true to the nature of things, that not one single feature of his scheme, so far considered, requires the slightest revision when examined with the powerful searchlight of Relativity. Only when the structure perforce comes into contact with Newtonian mechanics, as is the case in what we are about to consider, does it require, along with classical mechanics, to be amended.

\indent For stress put
\begin{equation}
  \dot M\alpha^\times=-\tfrac12\alpha^\times V_0\dot\omega\omega^\times+V_1(V_1\alpha^\times\dot\omega')\omega^\times.
  \tag{8}
\end{equation}
We have
\[
  \dot M_9\nabla_9=-\tfrac12\nabla V_0\dot\omega'\omega^\times+V_1(V_1\nabla_9\dot\omega')\omega_9^\times+v^\times
\]
by (2a) and (3a). Now
\[
  V_1(V_1\nabla_9\dot\omega')\omega_9^\times
  =\nabla_9V_0\dot\omega'\omega_9^\times-V_1\dot\omega'V_3\nabla\omega^\times
\]
of which the last term is zero by (4). Hence
\begin{equation}
  \left.\begin{aligned}
    v^\times&=\dot M_9\nabla_9+\tfrac12(\nabla_8-\nabla_9)V_0\dot\omega_8'\omega_9^\times\\
    &=\dot M_9\nabla_9+\tfrac12\nabla_9V_0\omega^\times\dot\psi_9^{-1}\omega^\times
  \end{aligned}\right\}
  \tag{9}
\end{equation}
by (5). By using $\dot\omega$ instead of $\omega^\times$ the last term on the right becomes $-\tfrac12\nabla_9V_0\omega^\times\dot\psi_9^{-1}\omega^\times$. Because of the factor $l^{-1}$ in $\dot\psi^{-1}=l^{-1}\psi^{-1}$ we may speak of $\dot\psi^{-1}$ as an \emph{anti}-density.

\noindent\textsc{§ 21. \emph{Stationary Action in a Riemann Manifold.}}---In our theorems of integration $d\dot p_c$ being equal to $V_cd p_1d p_2\ldots d p_c$ is contravariant. But since
\[
  ldb=lV_0d\alpha d p_n
\]

is an invariant, $d\alpha$, must be a covariant vector anti-density, and, similarly, $d\dot p_c$ in general (including $db$) is a covariant multenion anti-density.

\indent (8) and (9) of §20 are unsatisfactory in that $\dot M\theta$ is not apparently self-conjugate, and from the second term on the right of (9) the force (cum power), per unit volume, cannot be expressed as a stress even of the non self-conjugate kind. But the pioneers of Relativity have developed a principle of stationary action which has for consequence that all of the equations of §20 may be retained, and that $\dot M$ of (8) \emph{is} self-conjugate; and by what seems little more than a bold assumption as to the nature of things, they conclude that the intractable last term of (9) need not trouble us, because force generated by the field need not be begotten by stress; they give us something else instead, as will appear below. [I think I have here credited the pioneers with just a little more than they have set forth, but no more than such as, by small additions, can be included.]

\indent Although in the main we copy their procedure, yet in two important respects we depart from it, and in doing so return much more nearly than they do to Maxwell's outlook. First, we keep in mind the application to matter in bulk; and, secondly, we abjure the use of $\lambda$ \emph{in toto}. Probably Heaviside, among those who studied these matters in the days before the advent of Relativity, as we know it to-day, emphasised as much as any writer the view that the potentials of $\lambda$ have nothing to do with actualities, unless among such are to be reckoned the imperfections of our mathematics. In returning to the procedure I adopted in 1892,\footnote{``On the Mathematical Theory of Electromagnetism,'' `Phil. Trans.,' 1893.} of taking Maxwell's $\mathbf{H}$ and $\mathbf{D}$ as the fundamental (three-dimensional) vectors of the field, instead of any of $\mathbf{A}$, $\mathbf{B}$, $\mathbf{C}$, $\mathbf{K}$, or $\mathbf{E}$, I find that the potentials can be completely ignored, that they are not wanted even as a mathematical crutch, and that the ignoring of them considerably simplifies the preliminary organising plans of our exploration. The electro-magnetic terms are all dependent on a single arbitrary scalar function of the ${}_nV_2$ called $\dot\omega$ in §20.

\indent Let us recall that a particle of matter, with three dimensional extension, occupies a four-dimensional tube in the space-time continuum. The strength of this tube is the invariant constant (static) mass $d\dot M=ldM$ of the particle. The three dimensional anti-density covariant cross-section being $d\alpha$, an element of ``length'' of the tube being the proper time $ds$, corresponding to the contravariant vector $dp$, where $d\alpha$ and $dp$ are both time-like, $\iota$ being the unit tangent vector $dp/ds$, $db$ the four-dimensional volume of the element, we have
\begin{alignat}{3}
  ds^2&=V_0dp\theta dp, &\qquad \iota&=dp/ds, &\qquad db&=V_0d\alpha dp. \tag{1}
\end{alignat}
and
\begin{equation}
  d\dot M\mathbin{\cdot}ds=\dot m\,db,
  \tag{2}
\end{equation}

where $\dot m$ or sometimes $m$ is called the material density. The momentum is
\begin{equation}
  \iota d\dot M=\iota\dot m\,(db/ds),
  \tag{3}
\end{equation}
and
\begin{equation}
  \iota\dot m=\dot\mu,
  \tag{4}
\end{equation}
is called the momentum per unit volume. It is this vector whose magnitude multiplied by the cross-section yields the constant invariant strength, the quantity $d\dot M$. [In invariants we should say $m\times(ldb\div ds)=d\dot M$; so that perhaps $m\iota=\dot\mu$ rather than $\dot\mu$ should be called the momentum per unit volume.] We may note from the integration theorems, first for a finite portion of a tube, and next for a finite region of four dimensions, that
\begin{equation}
  V_0\nabla\dot\mu=0.
  \tag{5}
\end{equation}
This ``equation of continuity'' or ``equation of conservation'' of matter is given by Weyl, but seems nowhere to appear in Eddington's Report.\footnote{For a more critical examination of equation (5) see the later paper by the author: ``The Mechanical Forces \ldots\ Field'' referred to above, footnote p. 164.}

\indent Our fundamental view of electricity is the same in principle, but we do not need to elaborate a notation such as $d\mathbf{E}^\bullet$, $e^\bullet$, etc., in connection with $\dot\kappa$ to serve the purposes of $d\dot M$, $\dot m$, etc., in connection with $\dot\mu$. The tubes in the case of electricity (strength $V_0\dot\kappa' d\alpha$, just as strength of momentum tube $=V_0\dot\mu' d\alpha$) are called current tubes. If we were dealing with the smallest portions of matter it would accord with the prevalent views of to-day to suppose the two sets of tubes to be identical, but for matter in bulk we must consider them to be independent. Fortunately this is without any influence on the deductions we have to make from our fundamental principles.

\indent Our principle of stationary action is that
\begin{equation}
  \delta\int^n\!\int Wdb=0,
  \tag{6}
\end{equation}
where $W$ is an invariant scalar density (meaning that $l^{-1}W$ is invariant) function of position consisting of three parts
\begin{equation}
  W=\dot m+f(\dot\omega,\theta)+l\bigl(c-\tfrac12G\bigr),
  \tag{7}
\end{equation}
the first, material, being $\dot m$ already described and equal to $\sqrt{(V_0\dot\mu\theta\dot\mu)}$; the second, electrical, a function $f(\dot\omega,\theta)$ of a ${}_nV_2$, $\dot\omega$, identified with Maxwell's $\mathbf{H}$ and $\mathbf{D}$ as in §20, and of $\theta$ the fundamental linity of the manifold; and the third, gravitational, $l(c-\tfrac12G)$ where $G$ is the scalar curvature $-V_0\zeta C\zeta$, the cosmological constant $c$ possibly being zero, though Einstein teaches that it \emph{must} differ from zero.

\indent The variations implied by $\delta$ in (6) are of three independent kinds; (1) an arbitrary increment in $\theta$
\begin{equation}
  \delta\theta=\theta_0,
  \tag{8}
\end{equation}

at every point of the manifold; (2) an arbitrary increment in $\dot\omega$
\begin{equation}
  \delta\dot\omega=V_2\nabla v^\bullet\qquad(v^\bullet=V_3w^\bullet)
  \tag{9}
\end{equation}
at every point, of such a kind that the current tubes, that is to say $\dot\kappa$, which last is \emph{defined} by
\begin{equation}
  \dot\kappa=V_1\nabla\dot\omega
  \tag{10}
\end{equation}
remain unvaried; and (3) an arbitrary displacement which shifts the tubes, momentum and current, together,
\begin{equation}
  \delta p=\epsilon
  \tag{11}
\end{equation}
without alteration of their strengths, $V_0\dot\mu'd\alpha$, $V_0\dot\kappa'd\alpha$.

\indent [\emph{Added October, 1921.}---A consequence of remarks which occur below is so obvious that it is singular that only now have I happened to recognise it. If it is assumed that for the first two variations $(\delta\theta=\theta_0$ and $\delta\dot\omega=V_2\nabla v^\bullet)$ the stationary condition (6) holds; then \emph{it follows as a mathematical deduction} that the stationary condition also holds for the third variation $(\delta p=\epsilon)$ which shifts the tubes. Clearly then our postulated principle of stationary action should not include the \emph{assumption} of the stationary condition for the third variation. The reader will easily justify the assertions that have just been made, if he will take as the proof of the equation $T_{\xi}^{(t)}\zeta=0$, the indication of a proof, immediately following the identity (35) below. This equation $T_{\xi}^{(t)}\zeta=0$ is virtually the necessary and sufficient condition for the existence of the stationary condition under the variation $\delta p=\epsilon$.]

\indent The integral of (6) extends over any arbitrary finite region at the boundary of which $\theta_0$, $\dot\omega$ and $\epsilon$ all vanish.

\indent $\omega^\times$, eventually to be identified with Maxwell's $\mathbf{B}$ and $\mathbf{E}$ as in §20, is \emph{defined} by
\begin{equation}
  \omega^\times=-\dot\omega\nabla
  \tag{12}
\end{equation}
where $\dot\omega\nabla$, as in Quaternions, means the differential operator for which $df=-V_0d\dot\omega\,\dot\omega\nabla f$, $df$ being the increment in $f$ due to an arbitrary increment $d\dot\omega$ in $\dot\omega$.

\indent It is necessary to point out that if $W$ consisted of $\dot m$ alone, \emph{the condition $\delta W=0$ would mean that a momentum tube is geodesic} because by (2) $\dot m\,db=d\dot M\,ds$, and by hypothesis $d\dot M$, the strength, is unaffected by the displacement $\epsilon$ of the tube. Indeed, following Weyl, we might use this identification to find the effects of $\epsilon$ and $\theta_0$ on the variation of $\dot m$ in (6), but we pursue another course. It seems almost unnecessary to say that for a full treatment of matter in bulk $W$ as above described is by no means adequate;

there are several scalar functions of position such as temperature and strain co-ordinates (adapted to the method of Relativity) that should be included in $W$ for such a treatment.\footnote{Treated in a series of three papers by the author; despatched within the past month (December, 1922) to England.} We may accept the inclusion of fluid and elastic solid stresses that Eddington so lucidly justifies in §37 of his report, despite the fact that his appeal to the kinetic theory of matter would by a precisionist be described as a departure from our elected task of dealing with matter in bulk.

\indent We shall consider the effects of $w^\bullet$, $\epsilon$ and $\theta_0$ in succession because that is the order of simplicity; $\theta_0$, $w^\bullet$ and $\epsilon$ all take part in the variation of $f$; $\theta_0$ and $\epsilon$ in that of $\dot m$; and $\theta_0$ alone in that of $G$.

\indent The variation due to $\dot\omega$ gives by (9) and (12)
\[
  0=-\int\!\int V_0\omega^\times\nabla w^\bullet\,db
  =-\int^{n-1}\!\int V_0w^\bullet\omega^\times d\alpha
    +\int\!\int V_0w^\bullet V_3\nabla\omega^\times\,db,
\]
whence, since $w^\bullet$ is zero at the boundary,
\begin{equation}
  V_3\nabla\omega^\times=0.
  \tag{13}
\end{equation}
The process, here exhibited, of making a $\nabla$, explicitly occurring in the integrand, apply to the whole integrand, and the consequent removal of the resulting expression (or the potentiality of such removal) to the boundary, is known as integration by parts.

\indent The effect of $\epsilon$ we will first consider with reference to the current tubes. For the strength of every tube to remain unchanged by the displacement it suffices that $\dot\omega$ changes in precisely the manner that it changed in §15 when $\epsilon$ indicated a change of co-ordinates. Quoting then from §15, we first have
\[
  l^{-1}\delta'l=-\tfrac12V_0\xi\delta\theta\theta^{-1}\xi
  =V_0\xi\chi_0\xi=V_0\nabla\epsilon,
\]
and next
\[
  \begin{aligned}
    \delta\dot\omega&=\dot\omega V_0\nabla\epsilon-V_2\chi_0\xi V_1\xi\dot\omega\\
    &=\dot\omega V_0\nabla\epsilon-V_2\epsilon_9V_1\nabla_9\dot\omega
      =V_2\nabla_9V_3\dot\omega'\epsilon_9.
  \end{aligned}
\]
For the future we shall write $\delta'$ for these changes to be quoted from §15 leaving $\delta$ for a second meaning thus (Weyl's notation)
\begin{equation}
  \left.\begin{aligned}
    \delta'(\ )&=(\text{value at }p+\epsilon\text{ after displacement})-(\text{value at }p\text{ before})\\
    \delta(\ )&=(\text{value at }p\text{ after displacement})-(\text{value at }p\text{ before},)
  \end{aligned}\right\}
  \tag{14}
\end{equation}
so that
\[
  \delta(\ )=\delta'(\ )+V_0\epsilon\nabla\mathbin{\cdot}(\ ),
\]
\[
  \delta\dot\omega=V_2\nabla_9V_3\dot\omega'\epsilon_9+V_0\epsilon\nabla\mathbin{\cdot}\dot\omega,
\]
\[
  \delta f=V_0\omega_9^\times\delta\dot\omega
  =V_0\omega_9^\times\nabla_9V_3\dot\omega'\epsilon_9
   -V_0\mathbin{\cdot}V_3\dot\omega^\times\nabla_9V_3\dot\omega'\epsilon
   +V_0\epsilon\nabla_9V_0\omega^\times\omega_9',
\]
because $V_3\nabla\omega^\times=0$. Thus $\delta f-$ (first term on right)
\[
  =V_0\epsilon V_1\dot\kappa'\omega^\times.
\]
Hence
\begin{equation}
  (\epsilon\text{ part})\quad\int\!\int\delta f\,db
  =\int\!\int V_0\epsilon V_1\dot\kappa'\omega^\times\,db.
  \tag{15}
\end{equation}

Next calculate the corresponding part from $\delta\dot m$, in which $\delta\dot\mu$ is obtained, as with $\delta\dot\omega$ above, but $\delta\theta=0$. We shall for brevity use
\[
  \dot m^2=V_0\dot\mu\theta\dot\mu,\qquad \dot m\iota=\dot\mu,
\]
but the reader would do well at the beginning to use only $\dot m=\sqrt{(V_0\dot\mu\theta\dot\mu)}$ and not $\iota$.
\[
  \delta\dot m=\tfrac12\delta(\dot m^2)/\dot m
  =V_0\iota\theta\bigl(-\epsilon_9V_0\dot\mu\nabla_9+\dot\mu V_0\nabla\epsilon+V_0\epsilon\nabla\mathbin{\cdot}\dot\mu\bigr).
\]
Here prepare for integration by parts by altering the incidence of $\nabla$ in the first two terms on the right, so as to apply to the whole term, and reject this part of each term. Thus $\delta\dot m-$ (the rejected parts)
\[
  \begin{aligned}
    &=V_0\epsilon\bigl[(\theta\iota)_9V_0\dot\mu_9\nabla_9-\nabla\dot m+\nabla_9V_0\dot\mu_9\theta\iota\bigr]\\
    &=V_0\epsilon\bigl(T_9\nabla_9+\tfrac12\nabla_9V_0\zeta\theta_9\theta^{-1}T\zeta\bigr)
  \end{aligned}
\]
where
\begin{equation}
  T^\bullet\alpha^\times=lT\alpha^\times=\dot m\theta\iota V_0\alpha^\times\iota\quad[+\ \text{stress terms}]
  \tag{16}
\end{equation}
(the stress terms meaning those inserted here from Eddington's argument in §37 of the Report). Thus
\begin{equation}
  (\epsilon\text{ part})\quad\int\!\int\delta\dot m\,db
  =\int\!\int lV_0\epsilon T_{(\zeta)}\zeta\,db
  \tag{17}
\end{equation}
by equation (19) of §15.

\indent $T$ or $T^\bullet$ is called the material kinetic energy linity, or if the stress terms (say $T^{(s)}$) are included, it may be called the material stress-energy linity. The stress terms are self-conjugate and satisfy the equation $T^{(s)}=NT^{(s)}N$.

\indent (15) and (17) are the only contributions of $\epsilon$ to (6), so that we have
\begin{equation}
  lT_{(\zeta)}\zeta+V_1\dot\kappa'\omega^\times=0.
  \tag{18}
\end{equation}
We have now fully identified $\omega^\times$ with Maxwell's $\mathbf{B}$ and $\mathbf{E}$.

\indent Of equations (1a), (2a), (3a), (4), (5), (6) of §20 we have decided that (1a) is quite unnecessary, is in fact better neglected than cared for (though, of course, it is a mere consequence of our present equation (13)). The other five are replaced in the present section by (10), (18), (13), (10), (12) respectively. [(10) does duty twice, since it has the obvious consequence $V_0\nabla\dot\kappa=0$.]

\indent We are now about to trace the consequences flowing from the fact that $l^{-1}f$ is an invariant. These are wholly satisfactory for empty space, but not for matter in bulk, for which, in accordance with Maxwell's principles, we must modify our view as to the nature of $f$, the electrical term of the action density. It is well to show, however, that it is this invariance of $l^{-1}f$ which causes the inadequacy of our present supposition that $f$ is a function of $\dot\omega$ and $\theta$ only.

\indent With our present (vacuum) assumption, we shall show that (18) can be added to, thus ($\kappa=l^{-1}\dot\kappa$)
\begin{equation}
  -T_{(\zeta)}\zeta=V_1\kappa\omega^\times
  =2\bigl(l^{-1}(\thetaPi f)_{(\zeta)}\zeta\bigr).
  \tag{19}
\end{equation}

where, as in Quaternions, $\thetaPi$ means the self-conjugate linity differential operator given by
\begin{equation}
  df=-V_0d\theta\zeta\,\thetaPi f\zeta
  \tag{20}
\end{equation}
where $df$ is the increment in $f$ due to the arbitrary increment $d\theta$ in $\theta$. It is, therefore, not the linity $T^\bullet$, but, instead, the linity $T^{\bullet(t)}$ (= \emph{total energy linity} $=lT^{(t)}$) given by
\begin{equation}
  T^{\bullet(t)}=T^\bullet+2\thetaPi f
  \tag{21}
\end{equation}
which may be supposed to be handed on through space by some kind of contact action, because from (19)
\begin{equation}
  T_{(\zeta)}^{(t)}\zeta=0.
  \tag{22}
\end{equation}
As a matter of fact, we shall later, when considering gravitation, show that this equation can be transformed to the shape
\[
  \phi_9\nabla_9=0.
\]

\indent We will now prove (19) on our present hypothesis [$f$ as in (7)]. Let us, for brevity, use $v^\times$ as in (3a), §20; that is, we define $v^\times$ by
\begin{equation}
  v^\times=V_1\dot\kappa'\omega^\times=V_1(V_1\nabla\dot\omega')\omega^\times.
  \tag{23}
\end{equation}
Since $l^{-1}f$ is invariant, and since, as we saw above, before equation (14), $l^{-1}\delta'l=V_0\nabla\epsilon$, where $\epsilon$ now, as in §15, expresses an infinitesimal change of co-ordinates, we have as with $\delta\dot\omega$ above in the equation following (14)
\[
  fV_0\nabla\epsilon_9+V_0\epsilon\nabla f=\delta f
  =V_0\omega^\times\delta\dot\omega'-V_0\delta\theta\zeta\,\thetaPi f\zeta.
\]
$\delta\dot\omega'$ is given in the equation following (14) and $\delta'\theta$ is given in §15 as $-\chi_0'\theta-\theta\chi_0'$. Hence
\begin{equation}
  \begin{aligned}
    fV_0\nabla\epsilon_9+V_0\epsilon\nabla f
    &=V_0\omega_9^\times\bigl(V_2\nabla_9V_3\dot\omega'\epsilon_9+V_0\epsilon\nabla\mathbin{\cdot}\dot\omega'\bigr)\\
    &\quad+V_0\bigl(2\epsilon_9\thetaPi f\nabla_9-\epsilon\thetaPi f\zeta\theta_9\zeta\mathbin{\cdot}V_0\epsilon\nabla_9\bigr).
  \end{aligned}
  \tag{24}
\end{equation}
Here $\epsilon$ and all the $n^2$ derivatives $(\nabla_9,\epsilon_9)$ of $\epsilon$ are arbitrary, so that we may put arbitrary dummies $\gamma$, $\alpha^\times$, $\beta$ for $\epsilon$, $\nabla_9$, $\epsilon_9$ respectively. From the dummy $\gamma$ we get
\[
  \nabla f=\nabla_9V_0\dot\omega'_9\omega^\times-\nabla_9V_0\theta_9\zeta\,\thetaPi f\zeta.
\]
This identity, satisfied by $f$, does not depend on the invariance of $f$. It is easily seen to be a direct consequence of the assumption that $f$ is a function of $\dot\omega$ and $\theta$. A useful result is obtained from it by modifying the first term on the right, thus
\[
  \begin{aligned}
    \nabla_9V_0\dot\omega'_9\omega^\times
    &=V_1(V_3\omega^\times\nabla_9)\dot\omega'_9+v^\times &&\text{[eq. (23)]}\\
    &=V_1(V_3\omega_9^\times\nabla_9)\dot\omega'_9+v^\times &&\text{[eq. (13)]}\\
    &=\bigl(-\dot M_9\nabla_9+\tfrac12\nabla V_0\dot\omega'\omega^\times\bigr)+v^\times &&\text{[eq. (8) §20]},
  \end{aligned}
\]
whence
\begin{equation}
  \nabla f=-\dot M_9\nabla_9+\tfrac12\nabla V_0\dot\omega'\omega^\times+v^\times-\nabla_9V_0\theta_9\zeta\,\thetaPi f\zeta.
  \tag{25}
\end{equation}

From the dummy $\beta$ above in (24)
\begin{equation}
  \left.\begin{aligned}
    f\alpha^\times&=V_1(V_3\omega^\times\alpha^\times)\dot\omega+2\thetaPi f\alpha^\times,\\
    f&=-\dot M+\tfrac12V_0\dot\omega\omega^\times+2\thetaPi f
  \end{aligned}\right\}
  \tag{26}
\end{equation}

This is an identity (remember $\omega^\times$ is defined as $-\omega\mathbin{\cdot}\nabla f$) which we have proved $f$ to satisfy by reason of the invariance of $l^{-1}f$. Putting $\alpha^\times$ in (26) $=\nabla_9$ and subtracting (25), we get (19).

\indent (26) shows that $\theta^{-1}\dot M$ is self-conjugate, which the reader will at first be inclined to claim as a highly satisfactory result. I am not inclined to dispute this so far as empty space is concerned, but for matter in bulk it is precisely what we do \emph{not} desire. If $\theta^{-1}\dot M$ is self-conjugate for matter in bulk, then crystalline transmission of light cannot be explained on Maxwell's theory.

\indent To prove this, choose co-ordinates such that $\iota=\iota_4$; this requires co-ordinates which at the point under consideration reduce the medium to rest. Examine the case when the co-ordinates are Galilean, and Maxwell's equations become (1a), (2a), (3a) of §20. To translate $\dot M\alpha^\times$ into quaternion form, identify the multenion $\iota_2\iota_3$, $\iota_3\iota_1$, $\iota_1\iota_2$ with the quaternion $i$, $j$, $k$ and the multenion
\[
  \iota_1\iota_2\iota_3\iota_4=\upsilon
\]
with Hamilton's $\sqrt{-1}$ in his bi-quaternion $p+q\sqrt{-1}$. Thus $\dot\omega=\mathbf{H}+\upsilon\mathbf{D}$ and $\omega^\times=\mathbf{B}+\upsilon\mathbf{E}$. Indicate a ${}_nV_1$ such as $\alpha^\times$ by $(\sigma,s_4)$ where
\[
  \alpha^\times=\sum s_c\iota_c\quad(c=1,2,3,4),
  \qquad \sigma=s_1i+s_2j+s_3k,
\]
and indicate a ${}_nV_3$ such as $\upsilon\alpha^\times$ by $\upsilon(\sigma,s_4)=(\upsilon\sigma,\upsilon s_4)$. Thus [compare (3) and (3a) of §20]
\[
  \begin{aligned}
    V_1\alpha^\times\dot\omega
      &=V_1(\sigma,s_4)(\mathbf{H}+\upsilon\mathbf{D})\\
      &=(V\sigma\mathbf{H}+s_4\mathbf{D},-S\sigma\mathbf{D})=(\text{say})(\sigma',s_4'),\\
    V_1(V_1\alpha^\times\dot\omega)\omega^\times
      &=(V\sigma'\mathbf{B}+s_4'\mathbf{E},-S\sigma'\mathbf{E})\\
      &=(V[V\sigma\mathbf{H}+s_4\mathbf{D}]\mathbf{B}-\mathbf{E}S\sigma\mathbf{D},-S\mathbf{E}[\sigma\mathbf{H}+s_4\mathbf{D}]).
  \end{aligned}
\]
adding $-\tfrac12\alpha^\times V_0\dot\omega\omega^\times$ and making a slight quaternion reduction we obtain
\[
  \dot M\alpha^\times=\bigl(-\tfrac12V[\mathbf{H}\sigma\mathbf{B}+\mathbf{E}\sigma\mathbf{D}]+\tfrac12V\sigma V[\mathbf{H}\mathbf{B}+\mathbf{E}\mathbf{D}]+s_4V\mathbf{D}\mathbf{B},\quad S\sigma\mathbf{E}\mathbf{H}-\tfrac12s_4S[\mathbf{B}\mathbf{H}+\mathbf{E}\mathbf{D}]\bigr).
\]
\indent In the square array of sixteen scalars representing the linity $\dot M$ the fourth column contains the components of $V\mathbf{D}\mathbf{B}$, the fourth row contains those of $-V\mathbf{E}\mathbf{H}$, and the constituent in the fourth row and fourth column is $-\tfrac12S(\mathbf{B}\mathbf{H}+\mathbf{E}\mathbf{D})$. The other nine scalars are those of the vector linity of $\sigma$ in the expression for $\dot M\alpha^\times$ just written. Thus the conditions that $\theta^{-1}\dot M$ shall be self-conjugate are that
\[
  V\mathbf{D}\mathbf{B}=V\mathbf{E}\mathbf{H},\qquad V\mathbf{H}\mathbf{B}=V\mathbf{D}\mathbf{E}.
\]

\indent [These, of course, amount to saying that $\mathbf D$ is parallel to $\mathbf E$ and $\mathbf B$ to $\mathbf H$.]

\indent Now in a plane light wave $V\mathbf E\mathbf H$ is parallel to the ray and $V\mathbf D\mathbf B$ is normal to the wave front, so that to say that $V\mathbf D\mathbf B=V\mathbf E\mathbf H$ is to deny the possibility of double and conical refraction.

\indent Our assumption that $f$ is a function of $\dot\omega$ and $\theta$ only is contrary to Maxwell's, which asserts that $N$ and $I$ enter into the relations between $\dot\omega$ and $\omega^\times$, that is $\iota$, as well as $\dot\omega$ and $\theta$ is a constituent variable of $f$. We may take $f$ to be an explicit function of $\dot\omega$, $\dot\phi=I\dot\omega$ and $\theta$. Its general increment $df$, may be written
\[
  df=-V_0d\dot\omega\mathbin{\cdot}(\dot\omega\mathbin{\cdot}\nabla)f
     -V_0d\dot\phi\mathbin{\cdot}(\dot\phi\mathbin{\cdot}\nabla)f
     -V_0d\theta\zeta\,(\thetaPi)f\zeta.
\]
where $(\dot\omega\mathbin{\cdot}\nabla)$, $(\dot\phi\mathbin{\cdot}\nabla)$ and $(\thetaPi)$ are all \emph{partial} differential coefficient operators, whereas $\dot\omega\mathbin{\cdot}\nabla$ and $\thetaPi$ by contrast are \emph{total} differential operators. Thus $\dot\omega\mathbin{\cdot}\nabla f$ comes partly from $(\dot\omega\mathbin{\cdot}\nabla)f$ and partly from $(\dot\phi\mathbin{\cdot}\nabla)f$, and again $\thetaPi f$ comes partly from $(\thetaPi)f$ and partly from $(\dot\phi\mathbin{\cdot}\nabla)f$ (because $\theta$ is involved in $\dot\phi$). However, there is a part of $-V_0d\dot\omega\mathbin{\cdot}(\dot\omega\mathbin{\cdot}\nabla)f$, involving $d\iota$, still standing over.

\indent [In Maxwell's own case $(\dot\phi\mathbin{\cdot}\nabla)f=(-I+N)\omega^\times$ and we might assume this equation to hold in general, though this is not done below.] Writing this outstanding part as $-V_0d\iota\nabla f$ we shall have
\[
  df=-V_0d\iota\nabla f-V_0d\dot\omega\,\dot\omega\mathbin{\cdot}\nabla f-V_0d\theta\zeta\,\thetaPi f\zeta.
\]
Thus using $d\dot\omega'$, for the moment, to stand for the increment of $\dot\omega$ depending on $d\iota$ we have
\[
  \dot\omega'=V_2\iota V_1\theta\iota\dot\omega',
\]
\[
  d\dot\omega'=V_2d\iota V_1\theta\iota\dot\omega'+V_2\iota V_1\theta d\iota\dot\omega',
\]
\[
  -V_0d\iota\nabla f=df=-V_0d\iota\zeta\bigl\{V_1[V_1\theta\iota\dot\omega'](\dot\phi\mathbin{\cdot}\nabla)f+\theta V_1\dot\omega'V_1(\dot\phi\mathbin{\cdot}\nabla)f\zeta\bigr\}.
\]
Thus we may understand $\nabla f$ to mean as follows
\begin{equation}
  \left.\begin{aligned}
    \nabla f&=\xi^\bullet\iota
      =N\bigl\{V_1[V_1\theta\iota\dot\omega](\psi^\bullet\mathbin{\cdot}\nabla)f
       +\theta V_1\dot\omega V_1(\psi^\bullet\mathbin{\cdot}\nabla)f\iota\bigr\},\\
    \text{where}\quad \xi^\bullet\alpha&=N(\psi+\psi')\alpha,\qquad
      \psi\alpha=\theta V_1\dot\omega V_1(\psi^\bullet\mathbin{\cdot}\nabla)f\alpha,\\
    \text{and}\quad \xi^\bullet\iota&=((\xi^\bullet\iota)).
  \end{aligned}\right\}
  \tag{23a}
\end{equation}

\indent [\emph{Correction added September, 1921.}---$N$ was absent from equation (23a) as given in the original paper. If $N$ is omitted it is not necessarily true, but if $N$ is included it is necessarily true, that $\xi^\bullet\iota=((\xi^\bullet\iota))$. In the original presentation it was tacitly, and erroneously, assumed that $(\dot\omega\mathbin{\cdot}\nabla)f$, $(\psi^\bullet\mathbin{\cdot}\nabla)f$ and $(\thetaPi)f$ are always invariantive. They will doubtless be so when a proper invariantive form of $f$ has been selected, and it is desirable to make such a selection when possible and convenient. We will now show that, whether such a selection has been made or not, $\nabla f$, $\dot\omega\mathbin{\cdot}\nabla f$ and $\thetaPi f$ are always invariantive, that is to say $\nabla f$ is a covariant vector density, $\dot\omega\mathbin{\cdot}\nabla f$ is a covariant ${}_nV_2$, and $\thetaPi f$ is a

contraexco vector linity density. Let us put $\tau=x\iota$ where $x$ is an arbitrary invariant scalar function of position which we shall \emph{eventually} take to be $x=1$. Thus
\begin{equation}
  f=f(\dot\omega,\dot\psi,\theta)=f(\dot\omega,\Gamma'\dot\omega,\theta)
  \tag{A}
\end{equation}
$\Gamma'$ being that definite explicit function of $\tau$ and $\theta$ which is given by ($\Omega=$ a contravariant ${}_nV_2$ dummy)
\begin{equation}
  \Gamma'\Omega=V_2\tau V_1\theta\tau\Omega/V_0\tau\theta\tau.
  \tag{B}
\end{equation}
From this last $d\Gamma'$ can be expressed in terms of $d\theta$ and $d\tau$ (and it may be noticed in passing that since $\Gamma'$ is homogeneous and of degree zero both in $\theta$ and in $\tau$ we have by Euler's theorem for homogeneous functions
\begin{equation}
  V_0\tau_\tau\nabla_9\mathbin{\cdot}\Gamma_9'=0=V_0\theta\zeta_\theta\thetaPi_9\mathbin{\cdot}\Gamma_9'.
  \tag{C}
\end{equation}
the first of these explaining the presence of $N$ in (23a), for it leads to $V_0\tau_\tau\nabla f=0$). We now have
\begin{equation}
  \left.\begin{aligned}
    df&=-V_0d\dot\omega\mathbin{\cdot}(\dot\omega\mathbin{\cdot}\nabla)f-V_0d\dot\psi\mathbin{\cdot}(\psi^\bullet\mathbin{\cdot}\nabla)f-V_0d\theta\zeta\,(\thetaPi f)\zeta,\\
      &=-V_0d\dot\omega\,\dot\omega\mathbin{\cdot}\nabla f-V_0d\tau_\tau\nabla f-V_0d\theta\zeta\,\thetaPi f\zeta
  \end{aligned}\right\}
  \tag{D}
\end{equation}
the second being derived from the first by finding $d\dot\psi$ from
\[
  d\dot\psi=\Gamma'd\dot\omega+d\Gamma'\mathbin{\cdot}\dot\omega,
\]
$d\Gamma'$ coming from (B). In the second of (D) $d\dot\omega$, $d\tau$ and $d\theta$ are all \emph{independently, arbitrary and invariantive} and therefore their coefficients,
\[
  \dot\omega\mathbin{\cdot}\nabla f,\qquad \tau\nabla f,\qquad \thetaPi f,
\]
are \emph{without ambiguity and are invariantive}. These last important deductions have been made possible by using $\tau$ instead of $\iota$. $\tau$ (and therefore $d\tau$) is a general contravariant vector whereas $\iota$ is restricted by the equation $V_0\iota\theta\iota=1$.

\indent From the above we obtain by straightforward filling in of detail that
\[
  x_\tau\nabla f=(1-I)(\psi-\psi')\iota=\xi^\bullet\iota
\]
as given in equation (23a) above. Since $x_\tau\nabla f$ is independent in meaning of $x$, it is natural as in (23a) to write it in the form $\nabla f$.

\indent I take this opportunity of suggesting what is the proper form of $f$ \emph{in vacuo} and inside matter. \emph{In vacuo} we have
\[
  f=V_0\dot\omega\theta^{-1}\dot\omega=l^{-1}V_0\omega^\times\theta\omega^\times.
\]
Now $\theta=N\theta N'+I\theta I'$ so that

\[
  lf=V_0(N'\omega^\times)\theta(N'\omega^\times)+V_0(I'\omega^\times)\theta(I'\omega^\times).
\]
I would suggest that inside matter
\begin{equation}
  lf=V_0(N'\omega^\times)\xi(N'\omega^\times)+V_0(I'\omega^\times)\xi(I'\omega^\times)
  \tag{E}
\end{equation}
where $\xi$ is a function of $\theta$ and of position $(p)$ only. Thus
\begin{equation}
  lf=V_0\omega^\times(N\xi N'+I\xi I')\omega^\times.
  \tag{F}
\end{equation}

Without loss of generality $\xi$ may be replaced first by its self-conjugate part $\tfrac12(\xi+\xi')$ since (E) is thereby unchanged, and next by the part
\[
  \tfrac12[N(\xi+\xi')N'+I(\xi+\xi')I']
\]
since (F) is thereby unchanged. With $n=4$ the general $\xi$ would involve 36 scalars, but the part we have finally reached involves only 12 scalars, namely 6 for $N(\xi+\xi')N'$ and 6 for $I(\xi+\xi')I'$. We may, as a special case, take $\xi=A\theta A'$ where $A$ is a coexco linity function of position.]

\indent We have to add to the right of (24), (25), and (26) the following marked (24a), (25a), (26a) respectively:
\[
  V_0\epsilon\xi^\bullet\iota\mathbin{\cdot}V_0\iota\nabla_9-V_0\epsilon\nabla_8\mathbin{\cdot}V_0\iota_8\xi^\bullet\iota,\tag{24a}
\]
\[
  -\nabla_8\mathbin{\cdot}V_0\iota_8\xi^\bullet\iota,\tag{25a}
\]
\[
  \xi^\bullet\iota\mathbin{\cdot}V_0\iota\alpha^\times.\tag{26a}
\]
(15) and (19) also are modified by $V_1\kappa\omega^\times$ changing to
\[
  -lT_{(\zeta)}\zeta=V_1\kappa\omega^\times-\nabla_9V_0\iota_9\xi^\bullet\iota-(\xi^\bullet\iota)_9V_0\iota_9\nabla_9.\tag{15a and 19a}
\]
I think there can be no doubt\footnote{Especially in view of the paper twice referred to above: ``The Mechanical Forces \ldots\ in an Electromagnetic Field.''} that this amendment is the true expression for the force, cum power, per unit volume. It will be observed that no change occurs in the total energy linity $T^{\bullet(t)}=T^\bullet+2\thetaPi f$.

\indent Pass now to the variation $\delta\theta=\theta_0$. There is no trouble in supplementing (15) and (17) with the $\theta_0$ parts since no integration by parts is here involved. Thus since $\dot m=\sqrt{(V_0\dot\mu\theta\dot\mu)}$, $\delta\dot m=\tfrac12V_0\iota\theta_0\dot\mu$, and
\begin{equation}
  (\theta_0\text{ part})\quad\delta(\dot m+f)=-\tfrac12V_0\theta_0\zeta\theta^{-1}T^{\bullet(t)}\zeta.
  \tag{27}
\end{equation}
$G$ and $G^\bullet(=lG)$ are functions of $\theta$ and its first and second derivatives $(\nabla)$, and therefore of $\theta^\bullet$ and its first and second derivatives, where $\theta^\bullet=l\theta^{-1}$. Let $\theta^\bullet$ receive any infinitesimal increment $\delta\theta^\bullet$. At one part of our work we shall assume $\delta\theta^\bullet$ to have the meaning corresponding to $\delta\theta=\theta_0$ of equation (8); at another part $\delta\theta^\bullet$ will mean the actual increment $d\theta^\bullet=-V_0\alpha\nabla\mathbin{\cdot}\theta^\bullet$ of $\theta^\bullet$ on passing from the position $p$ to the position $p+d p=p+\alpha$ in the manifold; on some other occasion it may mean some combination of these; in a word, $\delta\theta^\bullet$ is any conceivable defined infinitesimal increment of $\theta^\bullet$. The following two equations are true:---
\[
  G^\bullet=-V_0\Omega\zeta\theta^\bullet\zeta=-V_0\nabla(F_\zeta\theta^\bullet\zeta-\theta^\bullet F_\zeta'\zeta)-V_0\zeta\theta^\bullet(F_{\zeta1}'F_{\zeta}'-F_{\zeta}'F_{\zeta1}')\zeta_1,\tag{28}
\]
\[
  -V_0\Omega\zeta\delta\theta^\bullet\zeta=-V_0\nabla(F_\zeta\delta\theta^\bullet\zeta-\delta\theta^\bullet F_\zeta'\zeta)-\delta\mathbin{\cdot}V_0\zeta\theta^\bullet(F_{\zeta1}'F_{\zeta}'-F_{\zeta}'F_{\zeta1}')\zeta_1.\tag{29}
\]
The great similarity between these equations is connected with the fact that the second term on the right of (28) is a function homogeneous and

quadratic in the derivatives $(\nabla_9,\theta^\bullet_9)$, and also homogeneous and of degree $-1$ in $\theta^\bullet$; (28) may, as a matter of fact, be derived from (29) by means of this property if attention is paid to symmetry in $\nabla_9$, $\theta^\bullet_9$, and $\nabla_8$, $\theta^\bullet_8$, in the quadratic terms of the form $(\nabla_9,\theta^\bullet_9,\nabla_8,\theta^\bullet_8)$. The proof of (29) is here outlined, and the proof of (28) is similar and simpler. From (24) of \S\,16 (remember $F_\zeta'\zeta=l^{-1}\nabla l$)
\[
  \begin{aligned}
  -V_0\Omega\zeta\delta\theta^\bullet\zeta={}&-V_0\nabla(F_\zeta\delta\theta^\bullet\zeta-\delta\theta^\bullet F_\zeta'\zeta)
    -V_0\delta\theta^\bullet_9\nabla_9F_\zeta'\zeta+V_0\delta\theta^\bullet_9\zeta F_\zeta'\nabla_9\\
   &{}-V_0\delta\theta^\bullet\zeta(-F_{\zeta1}'F_\zeta'+F_\zeta'F_{\zeta1}')\zeta_1.
  \end{aligned}
\]
Comparing this with (29), putting $F_\zeta'\zeta=\sigma$, and noting that
\[
  F_\zeta\theta^\bullet\zeta=-\theta^\bullet_9\nabla_9,\qquad
  \delta(F_\zeta\theta^\bullet\zeta)=-(\delta\theta^\bullet)_9\nabla_9
  \tag{30}
\]
it will be seen that we still have to prove that
\[
  \begin{aligned}
  -\delta\mathbin{\cdot}V_0\zeta\theta^\bullet(F_{\zeta1}'F_\zeta'\zeta_1-F_\zeta'\sigma)
  ={}&V_0\sigma\delta(F_\zeta\theta^\bullet\zeta)\\
   &{}+V_0\delta\theta^\bullet_9\zeta F_\zeta'\nabla_9
    +V_0\delta\theta^\bullet\zeta(F_{\zeta1}'F_\zeta'\zeta_1-F_\zeta'\sigma).
  \end{aligned}
\]
The only term that gives any trouble is the first on the left. In it $F_\zeta'\zeta_1$ and $\zeta$ may be interchanged because $F_\zeta'\zeta_1$ is self-conjugate in $\zeta$. Thus the term is
\[
  -\delta\mathbin{\cdot}V_0F_\zeta'\zeta_1\delta(\theta^\bullet F_{\zeta1}')\zeta
  =-2V_0F_\zeta'\zeta_1\delta(\theta^\bullet F_{\zeta1}')\zeta+V_0F_\zeta'\zeta_1(\delta\theta^\bullet)F_{\zeta1}'\zeta.
\]
Here, in the first term on the right, we need only consider the self-conjugate part of $\delta(\theta^\bullet F_{\zeta1}')$ because $F_\zeta'\zeta_1$ is self-conjugate in $\zeta$. This self-conjugate part is simple for
\[
  \theta^\bullet F_\beta'+F_\beta\theta^\bullet=lV_0\beta\nabla\mathbin{\cdot}\theta^{-1}=V_0\beta\nabla\mathbin{\cdot}\theta^\bullet-\theta^\bullet V_0\beta\sigma.
\]
When this is used the rest of the reduction is straightforward.

\indent The first of (30) may be used to modify three terms in (28) and (29). Similarly, $F_\zeta'\zeta=l^{-1}\nabla l$ may be used to modify four terms.

\indent Assuming now the truth of (28) and (29) the gravitational part of the action in (6) may be put
\[
  \int^n\!\int(cl-\tfrac12G^\bullet)\,db=-\tfrac12\int^{n-1}\!\int V_0(\theta^\bullet_9\nabla_9+l^{-1}\nabla l)\,d\alpha+\int^n\!\int g^\bullet\,db
  \tag{31}
\]
where
\[
  g^\bullet=cl+\tfrac12V_0\zeta\theta^\bullet(F_{\zeta1}'F_\zeta'-F_\zeta'F_{\zeta1}')\zeta_1.
  \tag{32}
\]
Using (27) and (29), the last for preparing $\delta g^\bullet$ for integration by parts, we have since
\[
  l^{-1}\delta l=-\tfrac12V_0\zeta\theta_0\theta^{-1}\zeta,\qquad
  \delta\theta^\bullet=-\theta^\bullet(\theta_0\theta^{-1}+\tfrac12V_0\zeta\theta_0\theta^{-1}\zeta)
  \tag{33}
\]
that
\[
  0=\delta\int^n\!\int Wdb=-\tfrac12\int^n\!\int db\,V_0\theta^\bullet\theta_0\zeta\{T^{(t)}+(C+\tfrac12V_0\zeta_1C\zeta_1)+c\}\zeta
\]
where $C$ stands for $\Omega\theta^{-1}$ and not for the conjugate $\theta^{-1}\Omega$ (we have in either case $V_0\zeta C\zeta=-G$). Hence
\[
  T^{(t)}+(C+\tfrac12V_0\zeta C\zeta)+c=0.
  \tag{34}
\]
Since $T_{(\zeta)}^{(t)}\zeta=0$ by (22) it follows that
\[
  (C-\tfrac12G)_{(\zeta)}\zeta=C_{(\zeta)}\zeta+\tfrac12V\nabla V_0\zeta C\zeta=0.
  \tag{35}
\]


\indent (35) has appeared incidentally in our application of stationary action, but it is really an identity satisfied by $C$ that has nothing to do with stationary action. That this is the case is proved by help of (28) and (29), from the fact that $G$ is an invariant in just the same manner as the identity (26), satisfied by $f$, was proved. Conversely, if we had chosen to establish this identity in the manner indicated (34) would furnish us with a second proof that $T_{(\zeta)}^{(t)}\zeta=0$.

\indent From (29) we get another identity of importance in connection with the conservation of energy. Let the $\delta$ of (29) stand for $-V_0\alpha\nabla\mathbin{\cdot}(\ )$, and let (29) thus come to read $V_0\alpha\sigma^\times=0$. We have, of course, $\sigma^\times=0$. If the reader will write out in full the meaning of this identity he will find it can be expressed thus. [Contrary to the conventions of these papers, $t^\bullet$ is here used to denote a vector linity, because $t$ is similarly used by writers on Relativity for the corresponding ``pseudo-tensor.''] If
\[
  2t^\bullet\alpha^\times=-\nabla_9V_0\alpha^\times(F_\zeta\theta^\bullet_9\zeta-\theta^\bullet_9F_\zeta'\zeta)-\alpha^\times V_0\zeta\theta^\bullet(F_{\zeta1}'F_\zeta'-F_\zeta'F_{\zeta1}')\zeta_1
  \tag{36}
\]
we have the identity
\[
  2t^\bullet_9\nabla_9=\nabla_9V_0\zeta\theta^\bullet\theta^{-1}(C-\tfrac12G^\bullet)\zeta=-2(C-\tfrac12G^\bullet)_9\nabla_9
  \tag{37}
\]
by (35) just written, and (19) \S\,15.

\indent The conservation of energy is regarded as the interpretation of the equation
\[
  0=(T^{\bullet(t)}+t^\bullet)_9\nabla_9=(T^\bullet+2\theta\thetaPi f+t^\bullet)_9\nabla_9
  \tag{38}
\]
where $T^\bullet$ comes from $\dot m$ (matter), $2\theta\thetaPi f$ comes from $f$ (the electromagnetic field), and $t^\bullet$ comes from the gravitational field. $t^\bullet$ is the only one of our many symbols which is not invariantive. [The reader might suppose that the terms in (19a) above, which involve $\nabla_9$, are not invariantive, but it is not difficult to show directly that they are.]

\indent A little reflection will show that \emph{the two identities (35) and (37) hold in Weyl's manifold} as well as in Riemann's. For $C$, $G$ and $F_\beta$ we have, in the notation of \S\,18, to read ${}^aC$, ${}^aG$ and ${}^aE_\beta$.

\indent \textsc{\S 22. \emph{A Matter that has been Overlooked.}}---I forgot to describe in the second paper a process I frequently use in my MS. work. It may prove useful to others, and is especially convenient when $n=4$. I am reminded of it by the $w^\bullet(=V_3w^\bullet)$ of (9) \S\,21. The use of $w^\bullet$, a contravariant \emph{density} may be replaced by a complementary (${}_nV_a$ replaced by a ${}_nV_{n-a}$) covariant multenion (not density); and, similarly, a covariant anti-density may be replaced by a complementary contravariant multenion; this latter is the replacement we are already familiar with of $d p_a$ by $d s_{n-a}$ and conversely.

\indent $l\dot v$ (where $\dot v=\iota_1\iota_2\ldots\iota_n$) is a covariant multenion. More generally, $xl\dot v$ where $x$ is invariant is a covariant multenion; and its reciprocal (of form

$yl^{-1}\dot v$ where $y$ is any invariant) is a contravariant multenion. It easily follows that the two linities
\[
  \Upsilon=l\dot v(\ ),\qquad \Upsilon^{-1}=l^{-1}\dot v^{-1}(\ )
  \tag{1}
\]
are the one a $c/\invc$ and the other a $\invc/c$, that is, of the same types as $\theta$ and $\theta^{-1}$. As an example, $\Upsilon\omega^\bullet$ above ($n=4$) is a covariant vector density; multiplying by $l^{-1}$, $\dot v\omega^\bullet$ or $\omega^\bullet\dot v^{-1}$ is a covariant vector, $=\sigma^\times$ say. Thus, in (9),
\[
  \delta\omega^\bullet=V_2\nabla\omega^\bullet=V_2\nabla(\sigma^\times\dot v)=V_2\nabla\sigma^\times\mathbin{\cdot}\dot v.
\]
[When $n=4$ the complement of a ${}_4V_2$ is a second ${}_4V_2$.] This serves to indicate that our use of $V_2\nabla w^\bullet$ is not far different from the usual use of $\omega^\times=V_2\nabla\lambda$, $\delta\omega^\times=V_2\nabla\delta\lambda$. [As a matter of fact, the $f(\omega^\bullet,\theta)$ used above, and the corresponding usual expression involving $\lambda$, can be shown to be ``reciprocal'' functions in Routh's sense.]

\indent [\emph{Added August, 1922.}---Let us illustrate both the methods of the present section and those of the ``digression'' in \S\,18 at the end of the second paper by considering the contractions of a certain identity,\footnote{``A New Identity affecting Questions in Relativity; and some Cognate Matters.'' A paper despatched to the \emph{Phil. Mag.}, in January, 1922, which as yet I have not seen in print.} namely,
\[
  V_3\zeta\varpi_\zeta'\omega=0,
  \tag{2}
\]
or as we will here write it
\[
  V_3\zeta\varpi_\zeta'\eta^{-1}V_2\gamma^\times\epsilon^\times=0=V_3\zeta(\varpi_\zeta'\eta^{-1})_\zeta V_2\gamma^\times\epsilon^\times.
  \tag{3}
\]
where $\varpi$ has the meaning given in \S\,16 above for a Weyl manifold. (The suffix $\zeta$ denotes absolute differentiation \emph{proper}.) This may be contracted twice, the contractions being obtained from the modification of (3),
\[
  V_0(V_3\alpha\beta\delta)(V_3\zeta\varpi_\zeta'\eta^{-1}V_2\gamma^\times\epsilon^\times)=0,
  \tag{3a}
\]
by first setting $\delta,\epsilon^\times=\zeta_1,\zeta_1$, and then setting $\beta,\gamma^\times=\zeta_2,\zeta_2$; finally writing $V_2\zeta_2\zeta_1=\zeta_2''$. The contractions are
\[
  V_2\zeta_1V_3\zeta\varpi_\zeta'\eta^{-1}V_2\gamma^\times\zeta_1=0.
  \tag{4}
\]
\[
  V_1\zeta_2''V_3\zeta\varpi_\zeta'\eta^{-1}\zeta_2''=0.
  \tag{5}
\]
(5) is the well-known identity (known, that is, in a Riemann manifold) numbered (35) in \S\,21 above. (3) asserts that a certain ${}_nV_3/{}_nV_2$ multenion linity (that is, a linity which acts solely on an arbitrary ${}_nV_2$ and converts it to a ${}_nV_3$) vanishes, (4) that a ${}_nV_2/{}_nV_1$ vanishes, and (5) that a vector vanishes.

\indent Consider the number of scalar equations involved in the case of the Riemann manifold. The three identities are relations among the $\tfrac12n^3(n^2-1)$ first differential coefficients of $\varpi$. (3) is equivalent to ${}_nC_3\mathbin{\cdot}{}_nC_2$ scalar equations

(4) to ${}_nC_2\mathbin{\cdot}n$ such equations; and (5) to $n$ such. The cases of $n=3$ and $n=4$ are particularly instructive. With $n=3$ there are 18 differential coefficients and the other three numbers in order are 3, 9, 3; and when $n=4$ there are 80 differential coefficients, and the three numbers in order are 24, 24, 4.

\indent In the case of $n=3$ there cannot, of course, really be nine independent equations resulting from the first contraction, since only three independent equations are given by (3). The details of this case are not difficult and may be left to the reader.

\indent In the case of $n=4$, the relativity case, the two numbers 24, 24 above suggest that perhaps the original equation (3) and its first contraction (4) are mathematically equivalent to one another, and this proves to be true. Put $F(\ )$ for the $V_3\zeta\varpi_\zeta'\eta^{-1}(\ )$ occurring in (3), (4), (5), so that $F$ is a $(c.{}_nV_3)/(c.{}_nV_2)$, from which we deduce that $F'$ is a $(\invc.{}_nV_2)/(\invc.{}_nV_3)$, that $\Upsilon F'\Upsilon^{-1}$ is a $(c.{}_nV_2)/(c.{}_nV_3)$, and so on.

\indent The contractions (4) and (5) of $F$ are
\[
  H\gamma^\times=V_2\zeta FV_2\gamma^\times\zeta,
  \tag{4a}
\]
where $H$ is a $(c.{}_nV_2)/(c.{}_nV_1)$, and
\[
  \sigma^\times=V_1\zeta H\zeta=V_1\zeta_2''F\zeta_2'',
  \tag{5a}
\]
$H$ and $\sigma^\times$, which are introduced for brevity, are \emph{defined} by (4a) and (5a).

\indent Now, presumptively, both $F$ and $H$ involve twenty-four independent scalars ($F$ may be supposed by definition to involve the twenty-four, but $H$ might possibly involve fewer, as when above we put $n=3$), so that it appears probable that in the special case $n=4$ (and also $n=3$, but \emph{certainly} not in general for $n>4$) that not only can $H$ be expressed in terms of $F$, but also $F$ can be expressed in terms of $H$. Actually I find that
\[
  \Upsilon H'\Upsilon^{-1}\omega^\times=\tfrac12V_3\omega^\times\sigma^\times+F\omega^\times,
  \tag{6}
\]
where $\omega^\times$ is a dummy, and $\sigma^\times$ may be expressed either in terms of $F$ or in terms of $H$. From (6), then, we easily express $F$ in terms of $H$. I do not think I have found the simplest way to (6), namely by aid of a rectangular array of twenty-four scalars specifying $F$. [In the verification of (6) we need not cumber ourselves with the inconvenience of the signs in
\[
  \iota_1^2=\iota_2^2=\iota_3^2=-\iota_4^2.
\]
It is legitimate to work with $\iota_1$, $\iota_2$, $\iota_3$, $\iota_4$, for which each $\iota_a^2=-1$.]

\indent The following can also be proved. If with $n=4$, $F$ is a $(c.{}_4V_2)/(c.{}_4V_2)$, that is a ${}_4V_3$ coexco linity, and if $H$ is its first contraction and $f$ its second according to the equations
\[
  \left.\begin{aligned}
    H\beta^\times&=V_1\zeta FV_2\beta^\times\zeta,\\
    f&=-V_0\zeta H\zeta=-V_0\zeta_2''F\zeta_2''
  \end{aligned}\right\}
  \tag{7}
\]

then, putting $C$ for the contractile part of $F$ (so that $F-C$ is the non-contractile part) we have
\[
  \left.\begin{aligned}
    C\omega^\times&=\tfrac12(F-\Upsilon F'\Upsilon^{-1}+\tfrac16f)\omega^\times\\
      &=\tfrac12V_2\zeta HV_1\omega^\times\zeta-\tfrac16f\mathbin{\cdot}\omega^\times
  \end{aligned}\right\}
  \tag{8}
\]
where $\omega^\times$ is a dummy. If in (7) and (8) $F$ is put equal to $\varpi\eta^{-1}$, or $\varpi'\eta^{-1}$, or $\varpi\theta^{-1}$, the contractions are the familiar contractions of the curvature.

\indent It should here be stated that throughout the above and also in general work with contractions the contractile parts can always be written down from the following easily verified statements. If $(\beta,\gamma)$, $(\beta,\gamma^\times)$ and $(\beta^\times,\gamma^\times)$ stand for any bilinear invariant forms in two vectors then the contractile parts of the forms are
\[
  -n^{-1}(\zeta,\eta^{-1}\zeta)V_0\beta\eta\gamma,\qquad
  -n^{-1}(\zeta,\zeta)V_0\beta\gamma^\times,\qquad
  -n^{-1}(\zeta,\eta\zeta)V_0\beta^\times\eta^{-1}\gamma^\times
\]
respectively.

\indent With the design of trying to avoid future chaos in notation I venture to recommend a development of notation to users of the present methods.

\indent The single letter symbol $\omega$ is not sufficient provision for the ${}_nV_2$ meaning. Also we have an insufficient supply of letter symbols for scalars. Let us release $p$, $q$, $r$, $s$, $u$, $v$, $w$ for scalar use and provide several letters for the ${}_nV_2$ meaning. I recommend that most of the small letters in Clarendon type be ear-marked as follows. Use
\[
  \mathbf{p},\ \mathbf{q},\ \mathbf{r},\ \mathbf{s}
\]
for general multenions;
\[
  \mathbf{a},\ \mathbf{b},\ \mathbf{c}
\]
for the ${}_nV_a$, ${}_nV_b$, ${}_nV_c$ meanings;
\[
  \mathbf{f},\ \mathbf{g},\ \mathbf{h},\ \mathbf{i},\ \mathbf{j},\ \mathbf{k}
\]
for the ${}_nV_2$ meaning. [With $n=4$ we may take
\[
  (\mathbf{i},\mathbf{j},\mathbf{k})=(\iota_2\iota_3,\iota_3\iota_1,\iota_1\iota_2).]
\]
Also use $\mathbf{l}$, $\mathbf{m}$, $\mathbf{n}$ for the ${}_nV_3$ meaning. $\mathbf{d}$, $\mathbf{e}$, $\mathbf{t}$, $\mathbf{u}$, $\mathbf{v}$, $\mathbf{w}$, $\mathbf{x}$, $\mathbf{y}$, $\mathbf{z}$ may be left for analogous uses as experience guides. Similarly occasional use may be made of capitals in Clarendon type, or perhaps such capitals had better be reserved for vectors in a quaternion system, as with $\mathbf{A}$, $\mathbf{B}$, $\mathbf{C}$, $\mathbf{D}$, $\mathbf{E}$, $\mathbf{F}$, $\mathbf{H}$, $\mathbf{K}$ above (and, in other papers recently forwarded, $\mathbf{V}=$ velocity and $\mathbf{U}=$ unit radial vector).

{\renewcommand\thefootnote{\fnsymbol{footnote}}\setcounter{footnote}{0}
\indent Is it a vain hope that the publication of the present papers may encourage some young mathematician to ``cultivate Hamilton's garden,'' his \emph{linear associative vector algebra}, in spite of the genial sneers of Weyl and Eddington who virtually thank the Lord that \emph{their} hands need not be soiled by this ``somewhat forbidding''\footnote{Eddington's \emph{Report on Relativity}, preface to first edition, reproduced in second edition.}---hah---``orgy of''\footnote{English translation by Brose of Weyl's \emph{Space, Time, Matter}, p. 54.} agricultural labour? During 30 years and more I have often glanced over the hedge and wondered sadly why fellow workers should be content to turn over the, perhaps lighter, but incomparably less fertile, patches in the immediate neighbourhood of the garden.}


\end{document}
